\documentclass[a4paper, twoside,openright]{report}

% Dimensione dei margini
\usepackage[a4paper,top=3cm,bottom=3cm,left=3cm,right=3cm]{geometry}

\usepackage[T1]{fontenc}
\usepackage[utf8]{inputenc}
\usepackage[fontsize=12pt]{scrextend}

% Lingua del testo
\usepackage[english]{babel}

\usepackage{microtype}
\usepackage{fancyhdr}
\usepackage[backend=biber,style=numeric,citestyle=science]{biblatex}
\usepackage{csquotes}
\usepackage{float}

% Librerie matematiche - SIMPLIFIED (no amssymb to avoid conflicts)
\usepackage{newtxtext}
\usepackage{amsmath}
\usepackage{newtxmath}
% Note: amssymb removed to avoid \Bbbk conflict - newtxmath provides most symbols
%\usepackage{amssymb}  % commented out due to conflict with newtxmath
%\usepackage{amsthm}   % uncomment if needed

%comments
\usepackage{comment}

% Uso delle immagini
\usepackage{graphicx}
% Uso dei colori
\usepackage[dvipsnames]{xcolor}

% Uso dei listing per il codice
\usepackage{listings}

% Per inserire gli hyperlinks tra i vari elementi del testo
\usepackage{hyperref}

% Diversi tipi di sottolineature
\usepackage[normalem]{ulem}

% Per numeri con virgola
\DeclareUnicodeCharacter{2212}{-}

% Additional packages
\usepackage{multirow}
\usepackage{tabularx}
\usepackage{rotating}
\usepackage{algorithm}
\usepackage{algorithmic}
%\usepackage{algpseudocode}
\usepackage{subfigure}
\usepackage{etoolbox}
\usepackage{comment}
\usepackage{nicematrix}
\usepackage{booktabs}

%APPENDIX
\usepackage[toc,page]{appendix}

% -----------------------------------------------------------------
% Modifica lo stile dell'header
%\pagestyle{fancy}
%\fancyhf{}
%\lhead{\rightmark}
%\rhead{\textbf{\thepage}}
%\fancyfoot{}
%\setlength{\headheight}{12.5pt}

% Rimuove il numero di pagina all'inizio dei capitoli
\fancypagestyle{plain}{
  \fancyfoot{}
  \fancyhead{}
  \renewcommand{\headrulewidth}{0pt}
}

% Modifica dello stile dei riferimenti, con il testo in cyano
\hypersetup{
    colorlinks,
    linkcolor=CornflowerBlue,
    citecolor=CornflowerBlue
}

% Aggiunti definizioni, teoremi, linea e listing
\newtheorem{definition}{Definition}[section]
\newtheorem{theorem}{Theorem}[section]
\newtheorem{assumption}{Assumption}

\providecommand*\definitionautorefname{Definition}
\providecommand*\theoremautorefname{Theorem}
\providecommand*{\listingautorefname}{Listing}
\providecommand*\lstnumberautorefname{Line}

\raggedbottom

%---------------------------------------------------------------
\makeatletter
\patchcmd{\chapter}{\if@openright\cleardoublepage\else\clearpage\fi}{}{}{}
\makeatother

\setlength\parindent{0pt}

\bibliography{bibliografia}

% -----------------------------------------------------------------
% -----------------------------------------------------------
% Convenience macros
% -----------------------------------------------------------
\newcommand{\R}{\mathbb{R}}
\newcommand{\C}{\mathbb{C}}
\newcommand{\E}{\mathbb{E}}
\newcommand{\Prob}{\mathbb{P}}
\newcommand{\F}{\mathcal{F}}
\newcommand{\dt}{\,\mathrm{d}t}
\newcommand{\dW}{\,\mathrm{d}W}
\newcommand{\dd}{\,\mathrm{d}}
\newcommand{\abs}[1]{\left|#1\right|}
\newcommand{\norm}[1]{\left\|#1\right\|}
\newcommand{\iunit}{\mathrm{i}}   % imaginary unit
\usepackage{xcolor} % Required for text colors

% Option B: Gray, italicized, right-aligned with a triangle (Classic Academic Style)
 \renewcommand{\algorithmiccomment}[1]{\textcolor{gray}{$\triangleright$ \textit{#1}}}



\begin{document}

\begin{titlepage}
\begin{figure}[!htb]
    \centering
    
\includegraphics[keepaspectratio=true,scale=1.0]{frontespizio.jpg}
\end{figure}


\begin{center}
    \LARGE{UNIVERSITY OF BOLOGNA}
    \vspace{5mm}
    \\ \large{DEPARTMENT OF STATISTICAL SCIENCES}
    \vspace{5mm}
    \\ \LARGE{Master degree in Quantitative Finance}
\end{center}

\vspace{15mm}
\begin{center}
    {\LARGE{\bf A Scoring Rule GAN for the generation of future probability distribution of financial processes}}
    
    
\end{center}
\vspace{30mm}

\begin{minipage}[t]{0.47\textwidth}
	{\large{Supervisors:}{\normalsize\vspace{3mm}
	\bf\\ \large{Pietro Rossi} \normalsize\vspace{2mm}\bf \\
     \large{Giovanni Della Lunga} \normalsize\vspace{3mm}\bf}}
\end{minipage}
\hfill
\begin{minipage}[t]{0.47\textwidth}\raggedleft
	{\large{Candidate:}{\normalsize\vspace{3mm} \bf\\ \large{Matteo Rofrano}}}
\end{minipage}

\vspace{30mm}
\hrulefill
\\\centering{\large{ACADEMIC YEAR 2025/2026}}

\end{titlepage}
\include{abstract}
\tableofcontents
\listoffigures
\listoftables

\include{intro_concepts/intro_concepts}
\chapter{Literature Review}
\label{chapter:literature_review}

Generative Adversarial Networks (GANs) represent one of the most influential
families of deep generative models introduced in recent years for image generation tasks.
Originally proposed by Goodfellow et al.~\cite{background:gan},
GANs provide a framework for learning complex probability distributions
through an adversarial optimization process between two neural networks:
a generator and a discriminator.\\
Unlike likelihood-based generative approaches, where the objective is to
explicitly estimate the probability density function of the observed data,
GANs belong to the class of implicit generative models.
The generator learns a transformation from a simple latent distribution into
samples that approximate the unknown data distribution, without requiring an
explicit formulation of the underlying probability density.

\section{Standard GAN model}
\label{sec:gans}
The objective of the generator (G) is to produce synthetic observations that are
indistinguishable from real data, while the discriminator (D) attempts to identify
whether a given observation originates from the true data distribution or from
the generator.
GANs are often formalized as a minimax game, where the generator seeks to minimize the value function V while the discriminator
aims to maximize it:

\begin{equation}
    \label{eq:GAN minimax}
\min _G \max _D V(D, G)=\mathbb{E}_{\boldsymbol{x} \sim p_{\text {data }}(\boldsymbol{x})}[\log D(\boldsymbol{x})]+\mathbb{E}_{\boldsymbol{z} \sim p_{\boldsymbol{z}}(\boldsymbol{z})}[\log (1-D(G(\boldsymbol{z})))]
\end{equation}

Where $p_{\text {data }}$ is the unknown data distribution and $p_{\boldsymbol{z}}$ is the known latent distribution, usually a multivariate Gaussian or uniform distribution.
The discriminator outputs a value $D(x)$, 
which represents the probability $p$ that $x$ is a real sample.
Its goal is to correctly classify the real samples as real and generated samples as fake,
effectively maximizing the likelihood of the observed data.
To measure this, the cross-entropy loss is adopted, as it is common in deep learning,
expressed as $plog(1-p)$. 
For real samples, the true label $y$ equals 1, 
while for generated samples the label is 0.
With this setup, the discriminator's objective becomes:
\begin{equation}
    \label{eq:discriminator objective}
    \max _D V(D)=\mathbb{E}_{\boldsymbol{x} \sim p_{\text {data }}(\boldsymbol{x})}[\log D(\boldsymbol{x})]+\mathbb{E}_{\boldsymbol{z} \sim p_{\boldsymbol{z}}(\boldsymbol{z})}[\log (1-D(G(\boldsymbol{z})))]
\end{equation}

where $\mathbb{E}_{\boldsymbol{x} \sim p_{\text {data }}(\boldsymbol{x})}[\log D(\boldsymbol{x})]$ define how good the discriminator recognize real samples.
and $\mathbb{E}_{\boldsymbol{z} \sim p_{\boldsymbol{z}}(\boldsymbol{z})}[\log (1-D(G(\boldsymbol{z})))]$ define how good the discriminator recognizes generated samples.\\
The generator's objective is to produce distribution parameters that maximize $D(x)$, pushing the discriminator to classify them as real:
\begin{equation}
    \label{eq:generator objective}
\min _G V(G)=\mathbb{E}_{\boldsymbol{z} \sim p_{\boldsymbol{z}}(\boldsymbol{z})}[\log (1-D(G(\boldsymbol{z})))]
\end{equation}

Where $\mathbb{E}_{\boldsymbol{z} \sim p_{\boldsymbol{z}}(\boldsymbol{z})}[\log (1-D(G(\boldsymbol{z})))]$ define how good the generator is in fooling the discriminator.


\begin{figure}[H]
    \centering
    \includegraphics[width=1.0\textwidth]{intro_concepts/images/gan_architecture.png}
    \caption{Standard GAN architecture}
    \label{GAN_architecture}
\end{figure}

Moreover, as proved in \cite{background:gan}, for a fixed generator, the optimal discriminator can be derived analytically, obtaining the following expression:

\begin{equation}
D_G^*(x)
=
\frac{p_{data}(x)}
{p_{data}(x)+p_g(x)}.
\label{eq:optimal_discriminator}
\end{equation}

Where $p_g(x)$ is the distribution induced by the generator, equivalent to the distribution of the generated samples $G(z)$ when $z \sim p_z(z)$. 
Substituting this expression into the original minimax objective gives:

\begin{equation}
V(G,D_G^*)
=
- \log(4)
+
2 \cdot
JSD(p_{data}||p_g),
\label{eq:GAN_JS}
\end{equation}

where $JSD$ denotes the Jensen--Shannon divergence between the real and
generated distributions. Therefore, minimizing the GAN objective is equivalent to minimizing the
Jensen--Shannon divergence between the generated distribution and the true
data distribution, so when the generator is also optimal $p_g = p_{data}$ then the objective function reaches its minimum at $-\log(4)$.
Unfortunately, this remains a theoretical result 
since it is not possible, usually, to compute the optimal discriminator in practice,
and the training of GANs is often unstable and prone to issues such as mode collapse and vanishing gradients which will be discussed more in detail later on.\\


Despite the simplicity of the original formulation, GAN training suffers from
several fundamental difficulties.
The main issues include unstable optimization, vanishing gradients, mode
collapse, and sensitivity to hyperparameter choices.\\
These problems arise mainly because the discriminator and generator are
optimized simultaneously in a non-convex game.
Unlike classical optimization problems, where a unique minimum can often be
identified, GAN training seeks a Nash equilibrium between two competing
players.\\

Arjovsky and Bottou~\cite{background:gan_train_problems} demonstrated that the original GAN
objective can lead to severe optimization difficulties when the generated and
real distributions lie on low-dimensional manifolds.
In high-dimensional spaces, the supports of $p_{data}$ and $p_g$ may have
little or no overlap, causing the Jensen--Shannon divergence to become
constant.
When the discriminator becomes optimal too quickly, it can perfectly
distinguish real and generated samples. In this situation 
$D(G(z))\rightarrow 0$
and the gradient received by the generator becomes extremely small, providing little information for improving its distributional approximation capabilities.
This phenomenon is known as the \textit{vanishing gradient problem}.\\

Furthermore, GAN optimization may exhibit oscillatory behaviour rather than
convergence.
Mescheder et al.~\cite{background:gan_train_problems2} showed that simultaneous gradient-based
optimization does not necessarily converge to the desired Nash equilibrium,
even when the theoretical optimum exists.\\

These limitations become particularly relevant when GANs are applied to
financial time series. Unlike image datasets, financial observations are characterized by non-stationarity,
 regime changes, noisy stochastic dynamics and heavy tails distribution.
Therefore, instability in adversarial training may directly affect the quality
of generated financial scenarios and the reliability of downstream applications
such as risk estimation and derivative pricing.

Academic findings highlighted in \cite{background:gan_time_series} suggest that one of the principal failure modes of GANs 
arises when the generator converges to a parameter configuration that yields only a limited variety of outputs.
Existing research indicates that the occurrence of mode collapse is attributed to the establishment of a local 
equilibrium within the minimax optimization framework.
A further critical challenge involves the vanishing gradient problem,
which emerges when the discriminator becomes overly proficient,
causing the generator's gradients to approach zero and consequently impeding its learning process.\\


The original GAN framework learns the unconditional distribution
$p_{data}(x)$ and generates samples independently of any additional
information. However, many practical applications require the generation of
samples conditioned on external variables. In particular, in financial
forecasting the objective is not simply to generate realistic asset prices,
but rather to estimate the conditional distribution of future prices given
historical market information.\\
Conditional Generative Adversarial Networks (CGANs), introduced by Mirza and
Osindero~\cite{background:cgan}, extend the original GAN architecture by incorporating
additional information into both the generator and discriminator networks.
The conditioning variable can represent any relevant contextual information,
such as class labels, historical observations, market indicators, or previous
time-series values.\\
Let $c$ denote the conditioning variable. In the context of financial
forecasting, $c$ may represent the observed trajectory of an asset over a
historical window $c_t=(x_{t-k+1},x_{t-k+2},...,x_t)$ where $k$ denotes the length of the historical observation window.\\
The objective of a conditional generator is to learn the conditional
distribution $p(x|c)$ and generate samples that are consistent with the provided context $c$.\\

The generator receives both the latent noise vector $z$ and the conditioning
information $c$ while the discriminator evaluates whether a sample is consistent with the
conditioning information.
The conditional discriminator therefore does not only determine whether a
sample is realistic, but also whether it is compatible with the given
condition.
The CGAN objective function is given by:


\begin{equation}
    \label{eq:CGAN minimax}
    \min _G \max _D V(D, G)=\mathbb{E}_{\boldsymbol{x} \sim p_{\text {data }}(\boldsymbol{x})}[\log D(\boldsymbol{x} \mid \boldsymbol{c})]+\mathbb{E}_{\boldsymbol{z} \sim p_{\boldsymbol{z}}(\boldsymbol{z})}[\log (1-D(G(\boldsymbol{z} \mid \boldsymbol{c})))]
\end{equation}

A standard way for training GANs can be summarized by the following algorithm:

\begin{algorithm}[H]
\caption{Standard Conditional CGAN Training}
\label{alg: cgan}
\begin{algorithmic}[1]
\REQUIRE Initial discriminator parameters $\theta_D$, initial generator parameters $\theta_G$, condition matrix $C$, output data $X$.
\WHILE{$\theta_G$ has not converged}
    \STATE Sample real data $\{x^{(i)}\}_{i=1}^m \sim \mathbb{P}_{data}$ and context $\{c^{(i)}\}_{i=1}^m$.
    \STATE Sample latent noise $\{z^{(i)}\}_{i=1}^m \sim p_z$.
    \STATE $\tilde{x}^{(i)} \leftarrow G(z^{(i)}, c^{(i)})$ \COMMENT{Generate synthetic samples}
    
    \STATE \COMMENT{Update Discriminator: Maximize log(D(x)) + log(1-D(G(z)))}
    \STATE $L_{D} \leftarrow -\frac{1}{m} \sum_{i=1}^m \left[ \log D(x^{(i)}, c^{(i)}) + \log(1 - D(\tilde{x}^{(i)}, c  ^{(i)})) \right]$
    \STATE $\theta_D \leftarrow \theta_D - \alpha \nabla_{\theta_D} L_{D}$ \COMMENT{Gradient descent and backpropagation}
    
    \STATE Sample latent noise $\{z^{(i)}\}_{i=1}^m \sim p_z$.
    \STATE \COMMENT{Update Generator: Maximize log(D(G(z))) } 
    \STATE $L_{G} \leftarrow -\frac{1}{m} \sum_{i=1}^m \log D(G(z^{(i)}, c^{(i)}))$
    \STATE $\theta_G \leftarrow \theta_G - \alpha \nabla_{\theta_G} L_{G}$ \COMMENT{Gradient descent and backpropagation}
\ENDWHILE
\end{algorithmic}
\end{algorithm}


\section{Wasserstein GANs and other improvements}
The instability of the original GAN and (inherently the CGAN as well) formulation motivated the development of
alternative training objectives.
A major contribution was introduced by Arjovsky, Chintala and Bottou
through the Wasserstein GAN (WGAN) framework~\cite{background:wgan}.
The key idea behind WGAN is to replace the Jensen--Shannon divergence used in
the original GAN objective with the Wasserstein distance, also known as the
Earth Mover distance which measures the minimum cost required to transform one
probability distribution into another.

The transition from traditional divergence-based
Generative Adversarial Networks to those utilizing the Wasserstein metric 
represents a significant shift in the geometric treatment of probability 
measures within deep learning. 
The primal definition of the Wasserstein-1 distance 
for two probability distributions, $\mathbb{P}_r$ (real) and $\mathbb{P}_g$ (generated) is given 
by:
\begin{equation*}
    W(\mathbb{P}_r, \mathbb{P}_g) = \inf_{\pi \in \Pi(\mathbb{P}_r, \mathbb{P}_g)} \mathbb{E}_{(x, y) \sim \pi} [\|x - y\|]
\end{equation*}
where $\Pi(\mathbb{P}_r, \mathbb{P}_g)$ denotes the set of all joint distributions $\pi(x, y)$ with marginals $\mathbb{P}_r$ and $\mathbb{P}_g$.
While this formulation provides a physically intuitive measure of "work" required to transform one distribution into another, in fact this calculates the minimum "cost" of transporting mass to transform one distribution into the other.
It is computationally intractable in the context of neural networks because the infimum over the space of 
all possible transport plans $\pi$ cannot be efficiently optimized via backpropagation as suggested in \cite{background:wgan}.
Moreover, even if the two distributions are discretized through a binning strategy, an high computational cost is associated with the simulation of terminal values which define the empirical distribution and the following
discretization and evaluation of the Wasserstein distance.\\

To bridge this gap, CWGAN models rely on the \emph{Kantorovich-Rubinstein duality} \cite{intro:duality}, which transforms the constrained minimization over transport plans into a maximization over a set of 1-Lipschitz continuous functions.
The application of this duality is critical because it replaces the search for a transport plan with a search for a function $f$ (often called the Critic) that maximizes the difference between the expected values of the two distributions.\\


The main changes that characterize the CWGAN with respect to the CGAN are:
\begin{itemize}
    \item The loss function of a standard GAN is based on Jensen-Shannon divergence minimization while the CWGAN is based on the Wasserstein distance, which provides a more stable training process and mitigates issues such as mode collapse.
    \item The discriminator, which in Wasserstein GAN architecture is called critic, is designed to output real-valued scores instead of probabilities, allowing it to better capture the differences between real and generated data distributions.
    For this reason, the discriminator network does not use a sigmoid activation function in the output layer.
    \item The use of gradient penalty term in the loss function has the objective to enforce the Lipschitz constraint required for the Wasserstein distance computation.    
\end{itemize}

\begin{theorem}[Kantorovich duality]
\label{eq:kantarovich}
Let $(\mathcal{X}, \mu)$ and $(\mathcal{Y}, \nu)$ be two Polish probability spaces and let $c : \mathcal{X} \times \mathcal{Y} \to \mathbb{R} \cup \{+\infty\}$ be a lower semicontinuous cost function, such that
\begin{equation*}
    \forall(x,y) \in \mathcal{X} \times \mathcal{Y}, \quad c(x,y) \geq a(x) + b(y)
\end{equation*}
for some real-valued upper semicontinuous functions $a \in L^1(\mu)$ and $b \in L^1(\nu)$. Then
\begin{itemize}
    \item[(i)] There is duality:
\end{itemize}
\begin{align*}
    \min_{\pi \in \Pi(\mu, \nu)} & \int_{\mathcal{X} \times \mathcal{Y}} c(x, y) \, d\pi(x, y) \\
    &= \sup_{(\psi, \phi) \in C_b(\mathcal{X}) \times C_b(\mathcal{Y}); \,\, \phi-\psi \leq c} \left( \int_{\mathcal{Y}} \phi(y) \, d\nu(y) - \int_{\mathcal{X}} \psi(x) \, d\mu(x) \right) \\
    &= \sup_{(\psi, \phi) \in L^1(\mu) \times L^1(\nu); \,\, \phi-\psi \leq c} \left( \int_{\mathcal{Y}} \phi(y) \, d\nu(y) - \int_{\mathcal{X}} \psi(x) \, d\mu(x) \right) \\
    &= \sup_{\psi \in L^1(\mu)} \left( \int_{\mathcal{Y}} \psi^c(y) \, d\nu(y) - \int_{\mathcal{X}} \psi(x) \, d\mu(x) \right) \\
    &= \sup_{\phi \in L^1(\nu)} \left( \int_{\mathcal{Y}} \phi(y) \, d\nu(y) - \int_{\mathcal{X}} \phi^c(x) \, d\mu(x) \right),
\end{align*}
and in the above suprema one might as well impose that $\psi$ be $c$-convex and $\phi$ $c$-concave.
\end{theorem}

To formalize the components of the Kantorovich duality
within the context of Wasserstein GANs (WGANs), 
we define the underlying mathematical structures and their functional roles.
In this framework, the pairs $(\mathcal{X}, \mu)$ and $(\mathcal{Y}, \nu)$ 
represent Polish probability spaces, where in the specific case of generative modeling,
$\mu$ corresponds to the real data distribution $P_r$ and $\nu$ represents the model's 
generated distribution $P_g$. The cost function $c(x, y) \colon \mathcal{X} \times \mathcal{Y} \to \mathbb{R} \cup \{+\infty\}$ serves as the metric 
for the "work" required to transport a unit of mass from a point $x$ in the real domain to a point $y$ 
in the generated domain. The central object of the primal problem is the transport plan $\pi \in \Pi(\mu, \nu)$, 
which is a joint distribution whose marginals are constrained to $\mu$ and $\nu$. 
This plan describes the optimal allocation of mass across the product space.
Transitioning to the dual formulation introduces the dual potentials $\phi$ and $\psi$.
In the architecture of a WGAN, the potential $\phi$ (often denoted as $f$) is parameterized as the Critic network.
These functions assign scalar values to points in the manifold, effectively measuring the "potential energy" difference used to calculate the Wasserstein distance. Mathematically,
these potentials reside in $L^1(\mu)$, the space of functions that are integrable with respect 
to the measure $\mu$. A critical component for achieving the final loss function is the $c$-transform (or $c$-conjugate) of $\phi$, defined as:

\begin{equation}
    \phi^c(x) = \inf_{y \in \mathcal{Y}} [c(x, y) - \phi(y)]
\end{equation}

This transform identifies the optimal "mate"
function that satisfies the feasibility constraint
$\phi(y) - \psi(x) \leq c(x, y)$. When the cost function $c(x, y)$
is chosen as the $L_1$ distance $\|x - y\|$, the $c$-transform forces 
$\phi$ to satisfy a Lipschitz continuity constraint. 
This transformation is what allows the WGAN to replace the intractable search over all transport plans $\pi$ with a simpler maximization over $1$-Lipschitz functions, 
providing the stable gradients necessary for training deep generative models.

So finally, the Kantorovich-Rubinstein duality can be expressed in the specific case of the Wasserstein-1 distance, given
$P_r$ and $P_g$ two probability distributions defined on a compact metric space $\mathcal{X}$ and where
the supremum is taken over all $1$-Lipschitz continuous functions $f: \mathcal{X} \to \mathbb{R}$, such that for all $x, y \in \mathcal{X}:   |f(x) - f(y)| \leq \|x - y\|$, then:
\begin{equation}
    \label{eq: emd_dual}
    W(P_r, P_g) = \sup_{\|f\|_L \leq 1} \left( \mathbb{E}_{x \sim P_r}[f(x)] - \mathbb{E}_{x \sim P_g}[f(x)] \right)
\end{equation}


In this duality, the "discriminator" is redefined as a "critic" ($D$), whose task is not to classify
inputs as real or fake, but to act as the optimal function $f$ that maximizes the distance between the 
expected scores of the real and synthetic distributions.
In the following thesis, $f$ is parametrized by using a neural network $f_w$ with weights $w$ restricted to a compact space $\mathcal{W}$ to enforce the Lipschitz constraint.
Substituting the generator output $G(z)$ for samples from $P_g$, the optimization objective function for the Wasserstein GAN is obtained.
The minimax objective function is the following:
\begin{equation}
    \label{eq:wgan_objective_function}
    L(G, f_w) = \min_G \max_{w \in \mathcal{W}} \mathbb{E}_{x \sim P_r}[f_w(x)] - \mathbb{E}_{z \sim p(z)}[f_w(G(z))]
\end{equation}
By enforcing $\|f_w\|_L \leq K$ (where $K$ is a constant, typically through weight clipping or gradient penalty), the Critic provides a linear, non-saturating gradient to the Generator, ensuring stable training even when the Generator is far from the target distribution.

In a Conditional GAN (CGAN), this distance is 
extended to account for auxiliary information $c$ (such as past history, drift, or volatility in 
time-series). So in the conditional setting, 
mathematically, the EMD is calculated over the pairs $(x, c)$, where $c$ is the history and $x$ is the generated output.
The real pair $(x_{i, real}, c_i)$ and fake pair $(G(z_i, c_i), c_i)$ are produced. The "Earth mass" is not being moved on a 1D line but instead it is being moved in a higher-dimensional space where one axis is the context and the other is the generated output. 
The distance is only minimized when the generator produces the correct $x$ for that specific $c$.

The objective becomes the minimization of the Wasserstein distance between the 
conditional distributions $\mathbb{P}_r(x|c)$ and $\mathbb{P}_g(x|c)$. 
By conditioning both the generator $G(z, c)$ and the critic $D(x, c)$ on the same context $c$,
the critic learns a "score surface" that is specific to the underlying stochastic process defined by $c$.

The use of the critic's output to compute the loss is mathematically 
justified because, under the 1-Lipschitz constraint, 
the difference between the average scores of the real batch and 
the fake batch directly approximates the Wasserstein distance as demonstrated in \cite{background:wgan}.
This is a superior training signal compared to Jensen-Shannon divergence because
$W(\mathbb{P}_r, \mathbb{P}_g)$ is
continuous and differentiable even when the supports of the two
distributions are disjoint, a common occurrence in the early stages of training.
Consequently, the generator receives stable, non-vanishing gradients that indicate 
not just if the generated sample is wrong,
but how far it must move in the feature space to align with the real
data manifold. 
To enforce the Lipschitz constraint required by the duality, modern implementations utilize a Gradient Penalty (GP),
which regularizes the norm of the critic's gradient with respect to its input toward unity, ensuring the function remains 
within the space of 1-Lipschitz functions $\mathcal{D}$.

\begin{algorithm}[H]
\caption{Conditional Wasserstein GAN with Gradient Penalty (CWGAN-GP)}
\label{alg:cwgan_gp}
\begin{algorithmic}[1]
\REQUIRE The gradient penalty coefficient $\lambda$, the number of critic iterations per generator iteration $n_{critic}$, the batch size $m$, Adam hyperparameters $\alpha, \beta_1, \beta_2$.
\REQUIRE Initial critic parameters $\theta_D$, initial generator parameters $\theta_G$.
\WHILE{$\theta_G$ has not converged}
    \FOR{$t = 1$ \TO $n_{critic}$}
        \STATE Sample real data $\{x^{(i)}\}_{i=1}^m \sim \mathbb{P}_r$ and context $\{c^{(i)}\}_{i=1}^m$.
        \STATE Sample latent noise $\{z^{(i)}\}_{i=1}^m \sim \mathbb{P}_z$.
        \STATE $\tilde{x}^{(i)} \leftarrow G(z^{(i)}, c ^{(i)})$ \COMMENT{Generate synthetic samples}
        \STATE $\epsilon \sim U[0, 1]$ \COMMENT{Random number for interpolation}
        \STATE $\hat{x}^{(i)} \leftarrow \epsilon x^{(i)} + (1 - \epsilon) \tilde{x}^{(i)}$ \COMMENT{Interpolated samples} 
        \STATE $L_{D} \leftarrow \frac{1}{m} \sum_{i=1}^m D(\tilde{x}^{(i)}, c^{(i)}) - \frac{1}{m} \sum_{i=1}^m D(x^{(i)}, c^{(i)}) + \lambda \mathbb{E}_{\hat{x}, c} [(\|\nabla_{\hat{x}} D(\hat{x}, c)\|_2 - 1)^2]$  
        \STATE $\theta_D \leftarrow \text{Adam}(\nabla_{\theta_D} L_{D}, \theta_D, \alpha, \beta_1, \beta_2)$
    \ENDFOR
    
    \STATE Sample latent noise $\{z^{(i)}\}_{i=1}^m \sim \mathbb{P}_z$.
    \STATE $L_{G} \leftarrow -\frac{1}{m} \sum_{i=1}^m D(G(z^{(i)}, c^{(i)}))$
    \STATE $\theta_G \leftarrow \text{Adam}(\nabla_{\theta_G} L_{G}, \theta_G, \alpha, \beta_1, \beta_2)$
\ENDWHILE
\end{algorithmic}
\end{algorithm}

An import part of the training procedure is the gradient penalty calculation, which is described in Algorithm \ref{alg:gradient_penalty}.
By sampling $\epsilon$ and creating $\hat{x}$, the Lipschitz constraint is enforced specifically on the straight lines between the real and fake data manifolds.

\begin{algorithm}[H]
\caption{Gradient Penalty Calculation for CWGAN-GP}
\label{alg:gradient_penalty}
\begin{algorithmic}[1]
\REQUIRE Real samples $x$, generated samples $\tilde{x}$, conditioning context $c$, critic network $D$.
\STATE $m \leftarrow \text{batch size}$
\STATE Sample $\epsilon \sim U[0, 1]^m$ \COMMENT{Uniform random interpolation weights}
\STATE $\hat{x} \leftarrow \epsilon x + (1 - \epsilon) \tilde{x}$ \COMMENT{Compute interpolated samples}
\STATE Enable gradient tracking for $\hat{x}$
\STATE $\hat{S} \leftarrow D(\hat{x}, c)$ \COMMENT{Evaluate critic on interpolates}
\STATE $\nabla_{\hat{x}} D \leftarrow \nabla_{\hat{x}} \hat{S}$ \COMMENT{Compute gradients of scores w.r.t. interpolates}
\STATE $L_2 \text{-norm} \leftarrow \|\nabla_{\hat{x}} D(\hat{x}, c)\|_2$ \COMMENT{Calculate norm across feature dimensions}
\STATE $GP \leftarrow \frac{1}{m} \sum (L_2 \text{-norm} - 1)^2$ \COMMENT{Compute mean squared error from unity}
\RETURN $GP$
\end{algorithmic}
\end{algorithm}

So the final objective function for the Conditional Wasserstein GAN with gradient penalty (CWGAN-GP) becomes:

\begin{equation}
    \begin{aligned}
        \min _G \max _D V(D, G) = & \mathbb{E}_{\boldsymbol{x} \sim p_{\text {d}}}[D(\boldsymbol{x} \mid \boldsymbol{c})] - \mathbb{E}_{\boldsymbol{z} \sim p_{\boldsymbol{z}}}[D(G(\boldsymbol{z} \mid \boldsymbol{c}))] + \lambda GP
    \end{aligned}
\end{equation}
where $\lambda$ is a hyperparameter that controls the strength of the penalty and it is usually set to $10$ as suggested in \cite{background2:gp_wgan}.\\
An additional advantage of using the Wasserstein distance is that it provides a meaningful measure of the distance between the real and generated data distributions, which allows for the implementation of an early stopping technique which
improve the training efficiency.\\

When the generated distribution and real distribution are far apart, the
Jensen--Shannon divergence often becomes constant.
Consequently, the discriminator saturates and the generator receives
vanishing gradients. Conversely, the Wasserstein distance changes continuously as the generated
distribution moves toward the target distribution.
Therefore, even when the two distributions have limited overlap, the critic
provides meaningful gradient information.\\

Other tentatives to improve the stability of GAN training. One of the most successful techniques are the spectral normalization and feature matching.
Miyato et al.~\cite{background:spectral_norm} proposed Spectral Normalization as a method to
control the Lipschitz constant of neural networks. The weight matrix of each discriminator layer is normalized according to its
spectral norm: $\bar{W}=\frac{W}{\sigma(W)}$
where $\sigma(W)$ represents the largest singular value of $W$.
Spectral normalization constrains discriminator updates and improves
stability without requiring additional gradient penalties.
It has become one of the most effective stabilization techniques for modern
GAN architectures.\\

Salimans et al.~\cite{background:feature_matching} introduced feature matching as a method to
reduce mode collapse. Instead of directly maximizing discriminator classification error, the
generator minimizes the difference between statistics of intermediate
discriminator features:

\begin{equation*}
L_G
=
\left|
\mathbb{E}_{x\sim p_{data}}D(x)
-
\mathbb{E}_{z\sim p_z}D(G(z))
\right|^2 .
\end{equation*}

This encourages the generator to reproduce multiple modes of the target
distribution rather than exploiting weaknesses of the discriminator.


\section{GANs for Time Series}
The architectures reviewed so far in this chapter, 
the original GAN, its conditional extension, and the variance-reduction/stabilization techniques, were developed 
primarily with static, i.i.d. data, most prominently images, in mind. 
Extending the adversarial framework to \emph{sequential} data introduces a distinct
set of architectural requirements: 
the generator and discriminator must be able to represent
temporal dependence, like autocorrelation, rather than merely a single, 
unconditional multivariate distribution, and, for variable-length or streaming applications, 
must do so with a fixed-size internal state rather than a fixed-size input/output tensor. 
This section reviews the principal families of sequential GAN architectures, 
culminating in the financial-time-series-specific 
proposals that directly motivate the modelling choices made later in this thesis.\\

The most direct way to extend a GAN to sequential data is to replace the feed-forward generator and discriminator networks of the original architecture with recurrent networks (LSTMs or GRUs) that process a sequence one step at a time,
propagating a hidden state $\mathbf{h}_\tau$ forward across the sequence index $\tau$. 
A generic recurrent generator produces a sequence $\hat{\mathbf{x}}_{1:H} = (\hat x_1, \dots, \hat x_H)$ autoregressively,
\begin{equation*}
\hat x_\tau = g_\theta\left(\mathbf{h}_{\tau-1}, z_\tau\right), \qquad \mathbf{h}_\tau = f_\theta\left(\mathbf{h}_{\tau-1}, \hat x_\tau, z_\tau\right), \qquad \tau = 1,\dots,H,
\end{equation*}

\cite{background:rcgan} propose the Recurrent Conditional GAN (RCGAN), developed originally for the generation of realistic multivariate medical time series. RCGAN uses an RNN (LSTM) for both generator and discriminator, with the generator receiving \emph{independent} noise draws at every time step rather than a single noise vector at initialization,
\begin{equation*}
\hat x_\tau = G_\theta\left(\mathbf{h}_{\tau-1}, z_\tau\right), \qquad z_\tau \stackrel{\text{iid}}{\sim} \mathcal{N}(0, I),
\end{equation*}
a design choice the authors show empirically increases sample diversity relative to single-noise-vector variants, directly counteracting the temporal mode-collapse tendency whereby a generator conditioned on a single static noise draw tends to produce sequences that are smooth, low-entropy interpolations rather than genuinely stochastic paths. The discriminator likewise processes the full sequence recurrently and is trained with the standard (or, in later variants, Wasserstein) adversarial loss applied either to the final hidden state or averaged across per-step outputs. RCGAN additionally supports conditioning on auxiliary static or time-varying labels.\\

Also \cite{background:timegan} propose a novel architecture for time series generation that addresses a specific shortcoming of purely adversarial recurrent GANs: 
the adversarial loss alone provides no explicit incentive for the 
generator to preserve the \emph{stepwise conditional} dynamics 
of the training data (i.e. the transition structure $p(x_\tau \mid x_{1:\tau-1})$), since the discriminator, 
in its simplest form, only ever evaluates entire generated sequences 
(or fixed-length windows) against real ones in aggregate, without an explicit term that penalizes mismatches in 
the one-step-ahead transition law. TimeGAN addresses this by combining the adversarial objective with a 
\emph{supervised} autoregressive loss, and by decoupling the 
space in which the generator operates from the original data space via a 
learned embedding.\\

The proposed architectures for time series proposed so far are designed for the generation of realistic sequences, but do not directly address the problem of \emph{forecasting}, 
where the goal is to estimate the conditional distribution of future values given a history. One of the most influential architectures in this space is ForGAN \cite{background:forgan}, 
which combines the conditional adversarial framework with a recurrent conditioning network.

\section{Probabilistic Forecasting using GANs}
The overwhelming majority of forecasting models developed in the quantitative finance and machine learning literature are trained to minimize a pointwise loss function, typically the mean squared error (MSE) or the mean absolute error (MAE), between a predicted scalar or vector and an observed realization. Formally, given a conditioning history $\mathbf{x}_{1:t} = (x_1, \dots, x_t)$ of a stochastic process, a point forecaster learns a deterministic mapping
\begin{equation*}
\hat{x}_{t+h} = f_\theta(\mathbf{x}_{1:t}), \qquad \theta^\star = \arg\min_\theta \; \mathbb{E}\left[ \ell\left(x_{t+h}, f_\theta(\mathbf{x}_{1:t})\right) \right],
\end{equation*}
where $\ell(\cdot,\cdot)$ is typically the squared error $\ell(x,\hat x) = (x-\hat x)^2$.
When the underlying data-generating process is stochastic, as is the case for essentially every financial time series, since asset prices, volatilities, and interest rates evolve under genuine aleatoric uncertainty, a point forecast conveys only the conditional mean (or, depending on the loss, the conditional median) of the future state, and discards all information about the shape, skewness, and tail behaviour of the conditional distribution.
This is a critical shortcoming for financial applications, where risk management, derivatives pricing, and portfolio construction depend fundamentally on the full conditional law of future returns rather than on a single expected value.\\

The correct object of interest is therefore the conditional predictive distribution
\begin{equation}
p\left(x_{t+h} \mid \mathbf{x}_{1:t}\right),
\end{equation}
or, in the multi-step-ahead case relevant to path-dependent risk metrics, the joint conditional distribution
\begin{equation}
p\left(x_{t+h}, x_{t+2}, \dots, x_{t+H} \mid \mathbf{x}_{1:t}\right)
\end{equation}
over a forecast horizon $H$. 
Estimating this object is the task of \emph{probabilistic forecasting} \cite{background:probabilistic_forecasting}. 
Classical approaches to probabilistic forecasting in econometrics rely on explicit parametric assumptions, for example, GARCH-family models coupled with a Gaussian or Student-$t$ innovation distribution,
or the Heston stochastic volatility model itself, 
which specifies a closed-form (up to Fourier inversion) 
characteristic function for the conditional log-price distribution. 
These approaches are interpretable and computationally tractable, 
but their accuracy is contingent on the correctness of the parametric assumption; 
misspecification of the innovation distribution or of the volatility dynamics propagates directly into biased risk estimates.\\

Generative Adversarial Networks 
offer an attractive alternative precisely
because they are implicit generative models
where the generator $G(\mathbf{c}, \mathbf{z})$, conditioned on a history $\mathbf{c} = \mathbf{x}_{1:t}$ 
and driven by an exogenous noise vector $\mathbf{z} \sim p_z$, is trained so that its induced output distribution approximates $p(x_{t+h} \mid \mathbf{x}_{1:t})$
without any parametric assumption on the shape of that distribution. 
Repeated sampling of $\mathbf{z}$ for a fixed conditioning history $\mathbf{c}$ produces an empirical ensemble of trajectories, from which arbitrary functionals of the predictive distribution, quantiles, moments, tail probabilities and so on,
can be estimated by Monte Carlo. This is the conceptual bridge between adversarial generative modelling and probabilistic time series forecasting,
and it is the foundation of the ForGAN architecture discussed below.\\

ForGAN, introduced by \cite{background:forgan},
is a conditional GAN architecture designed
explicitly for one-step-ahead probabilistic
forecasting of univariate time series. 
It combines the conditional adversarial training framework
with a recurrent conditioning network, so that the generator 
can encode an arbitrarily long history into a fixed-dimensional context vector 
before producing a sample of the next-step distribution.\\

The ForGAN's architecture is based on a generator and a discriminator, both utilizing Recurrent Neural Network (RNN) layers, specifically LSTM or GRU variants, in order to process sequential data.\\
The generator's objective is to forecast a future next value $x_{t+h}$ by transforming the noise vector $z$, 
conditioned on the historical window $c$.\\
The probability distribution of $x_{t+h}$ is modeled as $P(x_{t+h}|c)$ and obtained by running the generator multiple times with different noise inputs and same condition window $c$.\\
Structurally, the generator processes the condition window $c$ through an RNN layer to construct a "condition representation" whose size is given by hidden state size.
This representation is subsequently concatenated with the noise vector $z$.
The combined feature vector is then passed through two dense layers, culminating in a single scalar output representing the predicted $x_{t+h}$.\\
On the other hand, the discriminator has the role to inspect the validity of a given value $x_{t+h}$ in the context of the history $c$.
Unlike standard discriminators that might process inputs separately, the ForGAN discriminator takes the candidate $x_{t+h}$, which is sourced either from the ground truth dataset or the generator, and concatenates it to the end of the condition window $c$, forming an extended sequence $\{x_{0},...,x_{t+h}\}$.
This extended sequence is fed into an RNN layer, followed by a dense layer that outputs a validity score, determining whether the sequence is real or generated.
So the innovation is given by the presence of a conditioner which is implemented as a recurrent neural network that processes the conditioning window $\mathbf{c} = (x_{t-\ell+1}, \dots, x_t)$ of length $\ell$ and returns the terminal hidden state as a fixed-length embedding,
\begin{equation}
C_\phi(\mathbf{c}) = \mathrm{RNN}_\phi(x_{t-\ell+1}, \dots, x_t) \in \mathbb{R}^{d_c}.
\end{equation}\\

\begin{figure}[H]
    \centering
    \includegraphics[width=1.0\textwidth]{Introduction/images/forgan_rnn.png}
    \caption{ForGAN building blocks}
    \label{img:forgan_rnn}
\end{figure}

\begin{figure}[H]
    \centering
    \includegraphics[width=1.0\textwidth]{Introduction/images/forgan.png}
    \caption{ForGAN architecture}
    \label{img:forgan}
\end{figure}

The model is trained with the standard conditional minimax objective adapted to the time series conditioning structure. In fact, as proposed in \cite{background:forgan}, the objective is to minimize the Jensen-Shannon divergence between the real and generated data as seen before \ref{eq:CGAN minimax}.
Because the original ForGAN objective inherits the training pathologies
of the vanilla GAN, like vanishing gradients when the discriminator saturates,
and susceptibility to mode collapse (see Section~\ref{sec:limitations-gans}), 
a natural extension replaces the Jensen--Shannon-based minimax loss with the Wasserstein critic loss \ref{eq:wgan_objective_function}.
In this last case the model can be referenced as ForWGAN which represent a novel 
approach to increase the performance of the standard ForGAN.\\

While ForGAN 
was shown to perform strongly on the sensory
and synthetic benchmarks
for which it was originally developed, \cite{forgan:fin_gan} directly investigate 
its behaviour on financial return data, 
and show that a ForGAN trained with its original, 
generic binary-cross-entropy (BCE) adversarial loss is \emph{systematically outperformed},
on economically meaningful metrics, by both classical econometric baselines (ARIMA)
and by a variant of the same ForGAN architecture trained with a purpose-built, 
economics-driven loss function, the proposed Fin-GAN model.
This result is important for the present thesis because it demonstrates 
that the limitations catalogued in Section~\ref{sec:limitations-gans} 
are not merely a theoretical curiosity: 
they have a directly measurable, negative impact on downstream financial performance,
specifically in the ForGAN architecture that the SRForGAN model developed in this thesis builds upon.\\
The core diagnosis in \cite{forgan:fin_gan} 
is a form of \emph{objective misalignment} between what the ForGAN 
adversarial loss actually optimizes and what a 
financial forecasting application actually requires. 
In the Fin-GAN paper, the authors solve the inherent misalignment of the BCE-trained ForGAN by replacing the generic adversarial loss with a composite loss function that explicitly encodes 
the financial decision-making objectives of interest, 
including mean-squared error, sign accuracy, profit-and-loss, and risk-adjusted return. 
So, they worked on the loss function to make it more aligned with the financial forecasting task, 
rather than on the architecture itself. It is indeed, the same approach that is pursued in this thesis.\\


The Fin-GAN critique of BCE-trained ForGAN is,
in an important sense, structurally analogous to the motivation 
for the scoring-rule-based training paradigm underlying the SRForGAN architecture developed in this thesis.
Both approaches share the diagnosis that the original ForGAN training objective,
precisely because it is a generic adversarial divergence with no reference 
to the specific quantity the downstream application 
cares about, is an unreliable foundation for financial forecasting,
and both approaches respond by replacing that generic objective with one explicitly aligned 
to a well-defined target: 
in Fin-GAN's case, an economics-driven, 
decision-rule-specific composite loss 
(MSE, sign accuracy, PnL, and risk-adjusted-return terms); 
in the case of this thesis, a strictly proper scoring rule whose global minimizer is, 
by construction, the true conditional \emph{distribution} itself, 
rather than a decision-rule-specific functional of it.

\section{Scoring Rules and ForGAN advances}



%-------------------------------------------------------------
%-------------------------------------------------------------




\label{sec:energy_score_theory}
A scoring rule is any function $S:\mathcal P\times\Omega\to\mathbb R$ 
that assigns a numerical reward $S(P^\phi,y)$ to a 
probabilistic forecast $P\in\mathcal P$ once the 
outcome $y$ materialises \cite{background:scoring_rules}. 
Scoring rules originated as tools for \emph{forecast evaluation} and 
\emph{elicitation}. A forecaster who 
believes the true law is $P^\star$ is rewarded, 
in expectation, according to $S(P^\phi,P^\star)$ for
 whichever forecast $P^\phi$ is quoted, and 
 propriety is precisely the condition 
 guaranteeing that quoting $P^\phi =P^\star$ is the 
 forecaster's optimal, honest strategy.  
 Beyond evaluation, any strictly proper $S$ 
 can be repurposed as a \emph{training objective}: 
 minimising $-S(P^\phi, P^\star)$ over a parametric or 
 generative family $\{P^\phi\}$ pushes $P^\phi$ 
 toward the data-generating law, which is the 
 perspective adopted throughout this thesis.\\

This section formalises why minimising the empirical 
Energy Score of Eq.~\eqref{eq:energy_score_estimator} 
drives the generator$G_\phi$ toward the 
true conditional law $P^\star(\cdot|c)$, 
and states the regularity conditions under 
which this guarantee holds for the i.i.d.\ 
trajectory design of the subsequent experiments discussed in the following chapter.\\

Propriety is the most important charactersistic 
that a scoring rule can have and it is extensively leveraged during this work.
The authors of \cite{background:scoring_rules} show that propriety 
 can be fully characterized:
 \cite{background:scoring_rules} 
 show that every proper
 scoring rule derives
from a convex function $F$ on 
$\mathcal P$, where $F(P):=S(P,P)$ 
is the expected score under the forecast $P$ itself. 
Propriety is equivalent to $F$ being convex and 
to $S(P,\cdot)$ acting as a subtangent of 
$F$ at $P$; strict propriety corresponds 
to $F$ being strictly convex. $F$ 
thus plays the role of a generalized 
entropy, and the divergence $D(P^\star\|P^\phi)$ 
becomes, on suitable spaces, the 
Bregman divergence induced by $F$. 
Any admissible convex $F$ therefore 
yields a propriety-preserving training loss that can be optimized by gradient descent,
which is precisely what allows scoring rules 
to be tailored to a given problem.

\paragraph{Strict propriety.} 
For a scoring rule $S(P^\phi,y)$, 
define the expected score 
$S(P^\phi,P^\star):=\mathbb{E}_{Y\sim P^\star}S(P^\phi,Y)$. 
Following \cite{background:scoring_rules},
 $S$ is \emph{proper} relative to a class $\mathcal{P}$ 
 if $S(P^\star,P^\star)\le S(P^\phi,P^\star)\ \forall P^\phi\in\mathcal{P}$, 
 and \emph{strictly proper} if 
 equality holds only at $P^\phi=P^\star$. 
 Propriety is what makes $S$ a sound training 
 objective: it guarantees
$$D(P^\star\|P^\phi):=S(P^\phi,P^\star)-S(P^\star,P^\star)\ \ge\ 0,$$
is a valid statistical divergence, vanishing
iff $P^\phi=P^\star$.\\

The authors of \cite{background:scoring_rules} 
uses this property to demonstrate that given a fitting problem
of a parametric model $P^\phi$ where $Y_1, Y_2, \ldots, Y_n$ are the observed i.i.d. 
samples, and $\phi_0$ the true parameter value, then by 
maximizing the goodness-of-fit given by the expected score:

\begin{equation*}
    \mathcal{S}_n(\phi)=\frac{1}{n} \sum_{i=1}^n S\left(P_\phi, Y_i\right),
\end{equation*} 

then: $\arg \max _\phi \mathcal{S}_n(\phi) \rightarrow \phi_0 \text { as } n \rightarrow \infty$.\\


%Any $\phi$ minimising $\mathbb{E}_{c\sim\Pi}\,S(P^\phi(\cdot|c),P^\star(\cdot|c))$ therefore recovers the \emph{full}  conditional law for every $c\in\mathrm{supp}(\Pi)$,  not merely a point estimate or a moment of it.
 
One of the most widely used strictly proper scoring rules for multivariate continuous distributions is the Energy Score, which generalises the Continuous Ranked Probability Score (CRPS), which is the main scoring rule used in this work, to higher dimensions.\\
$S_E^{(\beta)}$ is defined as:
\begin{equation}
\mathrm{S}_E^{(\beta)}(P^{\phi}, \mathbf{y}) = \mathbb{E}\left\| \mathbf{X} - \mathbf{y} \right\|_d^\beta - \tfrac{1}{2}\, \mathbb{E}\left\| \mathbf{X} - \mathbf{X}' \right\|_d^\beta, \qquad \mathbf{X}, \mathbf{X}' \stackrel{\text{iid}}{\sim} P^{\phi},
\end{equation}

which is basically the Kernel Score $S_k(P^\phi,y)=\mathbb{E}[k(X,X')]-2\mathbb{E}[k(X,y)]$ 
with the negative-power kernel $k(x,y)=-\|x-y\|_2^\beta$. 
In \cite[Thm.~4]{background:scoring_rules}, authors show that
$S_k$ is strictly proper on the class of distributions with 
finite $\mathbb{E}\|X\|^\beta$ whenever $k$ is (conditionally) 
strictly positive-definite. For $\beta\in(0,2)$ the negative-power kernel satisfies this,
so $S_E^{(\beta)}$ is strictly proper on that class and by consequence also the univariate CRPS which is obtained by setting $\beta=1$.\\ 


%----------------------------------------------
Before using $S_E^{(\beta)}$ as a training loss, \cite{srforgan:srforgan} and \cite{background:scoring_rules} require:
 
\begin{itemize}
\label{scoring_rule_assumptions}
\item $\beta\in(0,2)$: the negative-power kernel $k(x,y)=-\|x-y\|^\beta$ is (conditionally) strictly positive-definite only in this range, which is what makes $S_E^{(\beta)}$ strictly proper rather than merely proper.
\item Finite $\beta$-th moments: $\mathbb{E}_{Y\sim P^\star}\|Y\|^\beta<\infty$ and $\mathbb{E}_{X\sim P^\phi}\|X\|^\beta<\infty$, without this the score, and hence the objective, may not even be finite.
\item A \emph{well-specified} model class: some $\phi^\star\in\Phi$ with $P^{\phi^\star}(\cdot|c)=P^\star(\cdot|c)$ for all $c$, needed for the population minimiser to be unique (\cite{srforgan:srforgan}, Lemma~12). This condition is satisfied by the theorem of the universal approximation of neural networks, 
which guarantees that a sufficiently large generator network can approximate any continuous function to arbitrary precision.
\item i.i.d.\ training pairs $(c_j,y_j)\sim\Pi\times P^\star(\cdot|c)$: if pairs are instead dependent (e.g.\ sliding windows of one sequence), consistency additionally requires stationarity and mixing conditions (\cite{srforgan:srforgan}, Thm.~3), which our design avoids by construction.
\item A generator admitting the reparameterisation trick: $X=G_\phi(z,c)$, $z\sim\mathcal N(0,I)$, differentiable in $\phi$, so $\nabla_\phi S_E$ can be estimated unbiasedly (App.~C, \cite{srforgan:srforgan}).
\end{itemize}



For our i.i.d.\ simulation 
design, this is not a modelling risk but an analytically verifiable property: 
with $\beta=1$, the log-price increment over $[X_{Ndt},X_{Ndt+\tau}]$ is Gaussian
 with variance $\sigma_j^2\Delta t$ and $\sigma_j\in(0,1.2]$ bounded above by 
 construction (Algorithm~\ref{alg:y condition}), so $\mathbb{E}|Y|<\infty$ trivially, 
 and $\mathbb{E}\|X\|<\infty$ follows for any generator with bounded output layer or 
 empirically bounded activations. This is a strictly easier position than the 
 prequential time-series case of \cite{srforgan:srforgan}, where the same 
 moment condition must hold \emph{uniformly over $t$} under stationarity and 
 mixing assumptions (their Theorem 3), an assumption we do not need, since our 
 $J$ trajectories are simulated i.i.d.\ rather than drawn as overlapping windows 
 of one sequence.\\

The proper scoring rules most relevant to this literature and in particular to this work is the Continuous Ranked Probability Score (CRPS) which is derived from the Energy Score by using a beta parameter equal to 1.
For a univariate predictive CDF $\hat F$ and a scalar outcome $y$:
\begin{equation}
    \label{eq:CRPS}
\mathrm{CRPS}(\hat F, y) = \int_{-\infty}^{\infty} \left( \hat F(x) - \mathbf{1}_{x \ge y} \right)^2 \, dx,
\end{equation}
which admits the equivalent, sample-based (kernel) representation in terms of two independent draws $X, X' \sim \hat F$,
\begin{equation}
\mathrm{CRPS}(\hat F, y) = \mathbb{E}\left| X - y \right| - \tfrac{1}{2}\, \mathbb{E}\left| X - X' \right|.
\end{equation}
This representation is directly usable as a training loss for a GAN generator: given $N$ samples $\{\hat x^{(i)}\}_{i=1}^N$ drawn from $G([\mathbf{c},\mathbf{z}^{(i)}])$, an unbiased estimator of the CRPS is
\begin{equation}
    \label{eq:CRPS_estimator}
\widehat{\mathrm{CRPS}} = \frac{1}{N} \sum_{i=1}^{N} \left| \hat x^{(i)} - x_{t+h} \right| - \frac{1}{2 N (N-1)} \sum_{i=1}^{N} \sum_{j \neq i} \left| \hat x^{(i)} - \hat x^{(j)} \right|,
\end{equation}


The Energy Score, and by consequence the CRPS, therefore plays a role directly analogous to the minimax value function $V(D,G)$ of the original GAN, but with the crucial advantage that it is a proper scoring rule with a known, well-behaved minimizer, removing the discriminator-induced gradient pathologies discussed in Section~\ref{sec:limitations-gans} at the cost of a bias--variance trade-off governed by the ensemble size $N$ (the estimator above is unbiased for the population Energy Score, but its variance decreases only at rate $O(1/N)$, so small ensembles yield noisy gradients.


\section{Applications in Quantitative Finance}
 
The appeal of adversarially or scoring-rule
trained generative models in quantitative 
finance stems from a structural feature 
shared by nearly every downstream task in the field: 
risk management, derivative pricing, and portfolio 
optimization all reduce, at some level, 
to computing an expectation (or a quantile, or a full distribution) 
of a functional of future asset paths. 
Classical Monte Carlo methods evaluate 
such expectations by simulating paths 
from an explicit, calibrated stochastic model 
(e.g. Heston, SABR, or a local-volatility model); 
GAN-based methods replace 
this explicit simulator with a 
learned, data-driven generator, 
potentially side-stepping model 
misspecification at the cost of 
requiring large amounts of 
representative historical data and 
forfeiting the interpretability of a closed-form model.\\

Distribution-valued
forecasting has a growing dedicated literature in finance.
\cite{background:quantile_regression} use a deep quantile-regression
network to forecast Value-at-Risk directly, 
reporting significant accuracy gains over linear quantile regression. 
A recent study \cite{background:applications} 
trains CNN and LSTM networks to output the parameters 
of skewed-$t$ predictive distributions, 
evaluating them with CRPS and PIT diagnostics 
and showing they are competitive with GARCH for 
VaR estimation. 
\cite{background:applications2} combine deep quantile regression with GAN-based scenario generation specifically to estimate VaR and Expected Shortfall jointly, explicitly motivating GAN-generated tail scenarios as a remedy for the difficulty classical models have capturing stressed regimes.\\

On the pricing side, generative and simulation-based deep models have been applied to option-market 
simulation and hedging, 
notably \cite{background:applications3}'s Deep Hedging framework and 
 Quant GAN for scenario generation. 
Also the work of \cite{background:applications4}, Tail-GAN, modify the adversarial objective itself to 
target VaR/ES elicitability directly, 
reflecting the same insight, 
that generic distributional losses under-serve the tail, 
that motivates using a proper scoring rule rather than a raw adversarial loss in SRForGAN. 
These works establish both the demand for, and the 
empirical viability of, GAN and neural-network-based 
probabilistic forecasters as inputs to standard risk and 
pricing pipelines, and SRForGAN can be positioned within 
this literature representing a suitable GAN-based architecture
capable of capturing the full conditional distribution of financial processes,
representing a valid alternative to parametric models and other generative architectures.


\section{Limitations of standard forecast techniques}
\label{sec:literature_gap}

Classical forecasting casts $x_{t+h}$ prediction as a regression problem,
with the model trained to minimize a pointwise loss such as RMSE. 
Under squared-error loss, the Bayes-optimal predictor is the conditional expectation,
\begin{equation}
\label{eq:mean_regression_target}
\hat{x}_{t+h} = \arg\min_{f} \; \mathbb{E}_{x_{t+h}}\big[(x_{t+h} - f(c))^2\big] = \mathbb{E}[x_{t+h} \mid c] = \mu(\rho(x_{t+h}\mid c)),
\end{equation}
i.e. the model converges
to the mean of the conditional distribution 
regardless of its shape. 
This is only informative when $\rho(x_{t+h}\mid c)$ 
is unimodal and concentrated, approximating a Dirac delta. 
When the true conditional distribution is multimodal, Eq.~\ref{eq:mean_regression_target} 
produces a value that need not lie near \emph{any} 
realizable outcome. 
\cite{background:forgan} illustrate this with a synthetic bimodal target where $x_{10}$ 
realizes one value in $80\%$ of cases and a second, 
disjoint value in the remaining $20\%$: 
the trained regressor converges to the 
weighted average $0.8$, 
a point with zero support 
under the true data-generating process. 
This failure mode is structural rather 
than a symptom of under-fitting or 
insufficient capacity. 
No amount of additional training data or 
model complexity corrects it, since the objective
in Eq.~\ref{eq:mean_regression_target} is itself misspecified 
relative to the estimand of interest 
whenever $\rho(x_{t+h}\mid c)$ deviates 
from near-degeneracy. 
The same argument applies to the Heston-driven volatility clusters targeted in this thesis, 
where the conditional density of log-price increments is heteroskedastic 
and can exhibit heavy tails and skewness 
inconsistent with a single-mode Gaussian-like summary.

Quantile regression \cite{background:quantile_regression} partially 
addresses this limitation by estimating 
conditional quantile functions $Q_{\alpha}(x_{t+h}\mid c)$ directly, 
via the asymmetric piecewise-linear (pinball) loss
\begin{equation}
\label{eq:pinball_loss}
\mathcal{L}_{\alpha}(x_{t+h}, \hat{x}_{t+h}) = \big(x_{t+h} - \hat{x}_{t+h}\big)\big(\alpha - \mathbf{1}[x_{t+h} < \hat{x}_{t+h}]\big), \qquad \alpha \in (0,1),
\end{equation}

consistent for the $\alpha$-quantile. 
Fitting a family $\{\hat{Q}_{\alpha_1}, \dots, \hat{Q}_{\alpha_K}\}$ over a grid of 
quantile levels in principle recovers an approximation 
of the full conditional distribution. 
In practice, however, each $\hat{Q}_{\alpha_k}$ 
is typically estimated by an independent regression, 
with no constraint enforcing monotonicity in $\alpha$. 
This gives rise to the \emph{quantile-crossing} problem: 
for $\alpha_i < \alpha_j$ the fitted quantiles 
can satisfy $\hat{Q}_{\alpha_i}(c) > \hat{Q}_{\alpha_j}(c)$, 
violating the definition of a valid distribution function. 
The pathology is most pronounced with dense quantile grids or 
small sample sizes \cite{background:mean_regression}, 
precisely the regime relevant to derivative-pricing 
and tail-risk applications where fine resolution in 
the extreme quantiles (VaR, ES) is required. Post-hoc 
rearrangement or monotonicity-constrained estimators 
\cite{background:quantile_regression2,background:quantile_regression3} 
mitigate but do not eliminate the problem, and add 
estimation overhead that scales with the number of 
quantile levels considered.

Both failure modes share a common root: 
they optimize a functional of $\rho(x_{t+h}\mid c)$ 
(its mean, or a finite set of quantiles) 
rather than the distribution itself, 
and the functional's fidelity to the 
underlying density degrades precisely in 
the multimodal, heteroskedastic regimes characteristic 
of stochastic volatility models. 
This motivates the shift toward implicit, 
sample-based conditional density estimation 
pursued by ForGAN and, subsequently, 
by the SRForGAN architecture proposed in this thesis.




\section{Limitations of GANs}
\label{sec:limitations-gans}
Despite their theoretical appeal as flexible, likelihood-free generative models, GANs are notoriously difficult to train in practice, and the literature documenting and diagnosing their failure modes is nearly as extensive as the literature proposing new architectures. This section reviews the principal sources of instability, with particular attention to the manifestations that are most consequential for financial time-series applications.\\

The original GAN formulation establishes tha, in the nonparametric limit, that is when both $G$ and $D$ have unlimited capacity and are optimized exactly, alternately, to their respective optima at each step,
the minimax game converges to a unique Nash equilibrium at which $p_g = p_{\mathrm{data}}$. This guarantee, however, does not extend to the practically relevant setting in which $G$ and $D$ 
are finite-capacity neural networks trained by simultaneous or alternating stochastic gradient descent. 
In this setting, the minimax objective \ref{eq:GAN minimax}
defines a non-convex, two-player game,
for which simultaneous gradient descent is not guaranteed to converge even locally,
let alone to the desired Nash equilibrium. 
The resulting optimization landscape is known to exhibit oscillatory, non-convergent trajectories even on simple, low-dimensional toy problems,
where the generator and discriminator can enter a persistent limit cycle rather than converging to the equilibrium.\\

One of the most widely documented failure modes is \emph{mode collapse}, 
refers to the failure mode in which the generator $G$ learns to produce outputs concentrated on a small subset of the modes
(or, in the extreme case, a single point or a very low-dimensional manifold) of the true data distribution
$p_{\mathrm{data}}$,
rather than covering its full support.
The phenomenon arises, as discussed in \cite{background:mode_collapse}, because the minimax objective, when optimized via alternating gradient steps
rather than the idealized best-response dynamic assumed in the original convergence proof,
provides no direct incentive for the generator to maintain output diversity:
if a subset of the noise space $\mathcal{Z}$ maps to outputs that reliably fool the current discriminator,
gradient descent on parameters of $G$ has no mechanism to \emph{penalize} 
the generator for ignoring the remainder of $\mathcal{Z}$, since the objective is evaluated only in expectation over $\mathbf{z}$ and contains no explicit diversity or entropy term.\\
Actually, gradient descent effectively optimizes the cases in which the generator 
need only to find, for each fixed discriminator, \emph{some} output that fools it, rather than a full distribution that matches $p_{\mathrm{data}}$
under every possible discriminator. 
\cite{background:wgan} further connect mode collapse to a distributional-support argument:
when $p_{\mathrm{data}}$ and $p_g$ have disjoint or low-dimensional-overlapping supports
(a generic situation when both are supported on low-dimensional manifolds embedded in a high-dimensional ambient space, as is typically the case for both images and, arguably, financial processes),
the Jensen--Shannon divergence between them is \emph{constant} (equal to $\log 2$)
and its gradient with respect to the generator's parameters is zero almost everywhere, 
providing no useful training signal and creating an optimization landscape in which the generator can become ``stuck'' 
having covered only a subset of the target modes, 
with no local gradient information indicating the direction of the missing modes.\\
In the context of probabilistic financial forecasting, 
mode collapse manifests specifically as \emph{underdispersion}: 
a generator conditioned on a given history $\mathbf{x}_{1:t}$ 
collapses across noise draws $\mathbf{z}$ 
to a narrow cluster of outputs, so that the induced 
predictive distribution $\hat p_G(x_{t+h} \mid \mathbf{x}_{1:t})$ 
has substantially smaller variance than the true conditional distribution.\\

A second, related instability is the \emph{gradient vanishing} that arises from gradient saturation at the discriminator.
Recall the original generator loss $\mathbb{E}_{\mathbf z}[\log(1 - D(G(\mathbf z)))]$:
when the discriminator is properly trained, 
it rapidly learns to assign $D(x) \approx 1$
for real samples and $D(G(\mathbf z)) \approx 0$ 
for generated samples that are not yet realistic. 
In this regime,
\begin{equation}
\frac{\partial}{\partial \theta} \, \mathbb{E}_{\mathbf z}\left[ \log\left(1 - D_\psi(G_\theta(\mathbf z))\right) \right] \; \longrightarrow \; 0,
\end{equation}
with $\theta$ the discriminator parameters and $\phi$ the generator parameters. This happens because
the function $\log(1-y)$ has near-zero slope as $y \to 0$, 
so that precisely when the generator is performing poorly and most needs a strong corrective gradient,
the gradient it receives vanishes. 
\cite{background:gan} proposed the standard workaround, training the generator to instead maximize $\mathbb{E}_{\mathbf z}[\log D_\psi(G_\theta(\mathbf z))]$, which restores a non-vanishing gradient in the early-training, 
poorly-performing regime, but at the cost of abandoning the interpretation of the generator's objective 
as the minimization of a well-defined statistical divergence between $p_g$ and $p_{\mathrm{data}}$,
and reintroducing exactly the mode-seeking, 
diversity-agnostic behaviour discussed above. 
\cite{background:wgan} show more generally that, 
whenever the discriminator is allowed to approach optimality 
(which is precisely what is needed for it to provide an accurate estimate of any $f$-divergence, including the Jensen--Shannon divergence underlying the original GAN), 
the generator's gradient with respect to the Jensen--Shannon-based objective vanishes in the limit, 
establishing vanishing gradients not as an incidental implementation 
artifact but as a structural property of divergence-minimizing GAN objectives whenever $p_g$ and $p_{\mathrm{data}}$ 
have (near-)disjoint support, precisely the motivation for the Wasserstein reformulation discussed, since the Earth-Mover distance, unlike the Jensen--Shannon divergence, 
remains a smooth, non-vanishing function of the generator's parameters even under disjoint supports.\\

All these pathologies are the exact reason why the introduction
of scoring-rule-based GANs is so important for financial forecasting applications and 
motivate the development of the SRForGAN architecture in this thesis.
































\chapter{Data, Models and Experimental Setup}
Experiments are conducted to evaluate the performance of the proposed SRForGAN model in comparison to the original ForGAN architecture in order to assess their ability to produce coherent conditional distributions given well known financial processes, 
in particular, the experiments adopt trajectories obtained from Geometric Brownian Motion and Heston stochastic volatility model. 
This chapter outlines how training data are computed, model configurations, training procedures, and evaluation metrics used in the experiments.

\section{Training Data}
In order to effectively train the Conditional GAN models, two kinds of data are required for the training phase:
\begin{itemize}
    \item Condition input data $C$: it is a $J \times N$ matrix, where $J$ is the number of trajectories and $N$ is the number of time steps of each trajectory that goes from time 0 up to $T$.
     This data represents the past information available to the model to infer a coherent output of the asset price.
     \item Target output data $X$: Since in this thesis we focus on a single horizon output, this data is a vector of size $J$, where $J$ represents the number of trajectories.
    So, each entry of the vector contains the future output at time $T + \tau$ associated with the respective past trajectory. 
    \end{itemize}

However, in the test phase, a third kind of data is required, that is the ground truth data distribution $D_\tau(c)$ which is the 
distribution of the future output $\tau$ steps ahead given the past trajectory $c$. In this experimental setting, this distribution can be 
computed analytically in the case of the Black-Scholes model and 
numerically in the case of the Heston model (or alternatively using the inverse fast Fourier transform to obtain the probability density function as shown in \ref{alg:heston_pdf}).\\

\section{The ForGAN Architecture}
\label{sec:forgan}

ForGAN \cite{background:forgan} reformulates the $h$-step-ahead time series forecasting as a conditional generative modeling problem,
replacing the point estimate $\hat{x}_{t+h}$ of classical regression with an implicit sampler for the full conditional density $\rho(x_{t+h} \mid c)$, where $c = \{x_0, \dots, x_t\}$ denotes the historical condition window. 
This design directly targets the two structural weaknesses of 
mean-regression forecasters discussed in previous chapter \ref{chapter:literature_review}: convergence to the conditional mean under multimodal targets, and the quantile-crossing pathology of explicit quantile regressors.

ForGAN builds on the conditional GAN (CGAN) framework of \cite{background:cgan}, 
in fact, the mimised loss function is the same of a standard conditional GAN model:
\begin{equation*}
\min_G \max_D V(D,G) = \mathbb{E}_{x_{t+h} \sim \rho_{\text{data}}(x_{t+h})}[\log D(x_{t+h} \mid c)] + \mathbb{E}_{z \sim \rho_z(z)}[\log(1 - D(G(z \mid c)))].
\end{equation*}
Where the target data to generate is no more a matrix representing an image but the next value of a sequential process.
At the Nash equilibrium of this game, the optimal generator $G^*$ implicitly represents the true conditional distribution 
$\rho(x_{t+h}\mid c)$, so that repeated sampling $z \sim \mathcal{N}(0, I)$ yields draws from the forecast distribution rather than a single summary statistic.

Architecturally, both $G$ and $D$ are built around a single recurrent layer (LSTM or GRU). 
The generator first encodes the condition window through the RNN layer to 
obtain a fixed-size condition representation, 
concatenates it with a noise vector 
$z \sim \mathcal{N}(0, I)$ of size $N$, 
and passes the result through two dense 
layers to produce the scalar forecast $\hat{x}_{t+h}$ as represented in schemas \ref{img:forgan} and \ref{img:forgan_rnn}. 
The discriminator appends the candidate value (either $\hat{x}_{t+h}$ from $G$ or the ground-truth $x_{t+h}$) 
to the end of the condition window, 
forming the extended sequence $\{x_0,\dots,x_{t+h}\}$, 
and feeds it through its own RNN layer followed by a dense 
layer that outputs a single real/fake validity score. 
Training proceeds asymmetrically: at each epoch the discriminator 
is updated for $D_{\text{iter}}$ steps while the generator
 is updated once, a ratio tuned alongside the RNN cell type, 
 hidden sizes $R_G, R_D$, noise dimension $N$, and condition length $C$ via hyperparameter optimization techniques like genetic algorithms. 
 Algorithm~\ref{alg:forgan_training} summarizes one training epoch.

\begin{algorithm}[H]
\caption{ForGAN Training (single epoch)}
\label{alg:forgan_training}
\begin{algorithmic}[1]
\REQUIRE Generator $G_\phi$, discriminator $D_\theta$, condition window $c$, target $x_{t+h}$, discriminator iterations $D_{\text{iter}}$
\FOR{$i = 1$ to $D_{\text{iter}}$}
\STATE Sample noise $z \sim \mathcal{N}(0, I)$
\STATE $\hat{x}_{t+h} \leftarrow G_\phi(z, c)$ \COMMENT{generated forecast}
\STATE $\ell_D \leftarrow -\big[\log D_\theta(x_{t+h} \mid c) + \log(1 - D_\theta(\hat{x}_{t+h} \mid c))\big]$
\STATE Update $\theta$ by descending $\nabla_\theta \ell_D$
\ENDFOR
\STATE Sample noise $z \sim \mathcal{N}(0, I)$
\STATE $\hat{x}_{t+h} \leftarrow G_\phi(z, c)$
\STATE $\ell_G \leftarrow -\log D_\theta(\hat{x}_{t+h} \mid c)$
\STATE Update $\phi$ by descending $\nabla_\phi \ell_G$
\end{algorithmic}
\end{algorithm}
As can be denoted from the following algorithm is that 
it reproduce the same logic of the standard conditional GAN training process like in
\ref{alg: cgan}. In fact, an easy extension of the ForGAN model using the Wasserstein Loss
can be easily deduced by adopting the algorithm of the Wasserstein GAN with Gradient Penalty \ref{alg:cwgan_gp}, which is denominated ForWGAN during this thesis.\\

To isolate the contribution of adversarial training from that of the architecture itself, 
the original authors also define a \emph{G-regression} baseline: a network 
structurally identical to $G$, 
but trained by minimizing RMSE directly instead of playing the minimax game. 
Across the Lorenz, Mackey-Glass, and Internet-traffic (A5M) benchmarks, 
ForGAN is shown to match or beat G-regression and prior state-of-the-art 
on point-wise metrics (RMSE, MAE, MAPE) 
while dramatically outperforming both on the 
Kullback-Leibler divergence between the sampled 
forecast distribution and the empirical target distribution.\\
Despite this success, 
vanilla adversarial training inherits well-known 
instabilities: vanishing gradients when $D$ saturates, 
and mode collapse, 
whereby $G$ learns to reproduce only a subset of the 
modes of $\rho(x_{t+h}\mid c)$ while still fooling $D$ as described in section \ref{sec:limitations-gans}. 
These pathologies motivate the design choices underlying SRForGAN, 
introduced in the following section, 
which departs from the adversarial signal of Eq.~\ref{eq:CGAN minimax} 
in favor of a scoring-rule-based objective, 
retaining ForGAN's conditional generative structure 
while removing its reliance on a saturating discriminator game.


\section{The SRForGAN Architecture}
Mode collapse, also known as the "Helvetica scenario," 
occurs when the generator learns to produce only a 
limited subset of the data distribution as already explained in section \ref{sec:limitations-gans}. 
It can be categorized into two forms, especially in the context of 
conditional time-series GANs like ForGAN:
\begin{itemize}
    \item Distribution Mode Collapse: A collapse in the variety of samples generated for a single specific condition. The model fails to capture the probabilistic uncertainty of the future process.
    \item Mean Mode Collapse: A collapse in the distribution of the means of generated values across all different conditions. This effectively reduces the model to producing similar point estimates regardless of the historical input
\end{itemize}

When both occur simultaneously, 
the model yields a single point estimate for every input, 
completely losing its probabilistic forecasting capabilities, a scenario similar to outputs seen in the classical ForGAN experiments \ref{forgan_simple_sample_sigma}. 
Research indicates that mode collapse is often tied to the generator
converging toward sharp local minima. 
Once $G$ finds a sample that receives a high score from $D$,
it may continue to reproduce that sample. 
On the other hand,
if $D$ cannot identify this repetitive behavior (due to converging to its own local minima), 
the gradients for $G$ vanish, and the model becomes stuck.



The Fin-GAN paper \cite{forgan:fin_gan} identifies specific triggers and solutions for these performance issues caused by mode collapse in the ForGAN model. The key findings include:
\begin{itemize}
    \item Weight Initialisation: Using He initialisation led to mode collapse in 67\% of simulations for ForGAN; however, switching to Xavier initialisation significantly alleviated the problem.
    \item Loss Function Regularization: Including standard supervised terms like Mean Squared Error (MSE) or Sharpe Ratio (SR) with high coefficients can actually induce mode collapse by encouraging the model to only output values close to the target, so new terms are required to avoid this behaviour.
    \item The paper suggests that adding a Profit and Loss (PnL) term to the generator's loss function acts as an effective regularizer. This term shifts the generated distributions and prevents the model from collapsing into single-point estimates, ensuring a more robust probabilistic forecast.
\end{itemize}

Also in \cite{srforgan:srforgan}, the authors propose replacing the adversarial loss with a strictly proper scoring rule, specifically the Energy Score, which is designed to directly optimize the generator for accurate probabilistic forecasting.
This approach addresses the fundamental issues of mode collapse and training instability by providing a theoretically sound objective that encourages the generator to produce well-calibrated distributions rather than just realistic samples.\\

In particular, the paper proposes replacing the adversarial objective with a scoring rule whose theory is covered in \ref{sec:energy_score_theory}.
To sum up, it is a function $S(P^\phi, y)$ 
that takes a forecast distribution $P^\phi$ and a single observed realization $y$ and returns, in our framework, a penalty score where the lower the score, the better the forecast.\\
The key definition is strict propriety. A scoring rule $S$ is strictly proper relative to a class $\mathcal{P}$ if:
$$\mathbb{E}_{Y \sim P^*}[S(P^*, Y)] < \mathbb{E}_{Y \sim P^*}[S(P^\phi, Y)] \quad \forall P^\phi \neq P^*$$
Where the expected score is uniquely minimized when the forecast distribution equals the true data distribution $P^*$. No other distribution achieves a lower expected score. This is the property that makes a scoring rule the right objective for distribution learning because
it is theoretically guaranteed to converge to the right answer, provided that some conditions are met, not just some approximation like the 
GAN objective that minimizes a divergence between distributions which has a similar guarantee in principle, but the GAN only minimizes a variational lower bound of that divergence using a finite-capacity critic, which breaks the guarantee in practice.\\



%-----------------------------------------------------


The proposed model follows the scoring-rule-based conditional generative framework introduced by Pacchiardi et al. \cite{srforgan:srforgan}. Let $c\in\mathcal{C}$ denote the conditioning variable and let $y\in\mathcal{Y}$ denote the corresponding future realization. The objective is to learn a conditional probability distribution
\[
P^\phi(\cdot\mid c),
\]
parameterized by a neural network with weights $\phi$, such that it approximates the true conditional distribution
\[
P^*(\cdot\mid c).
\]

Unlike adversarial approaches, the model is trained by directly minimizing a strictly proper scoring rule. Given a scoring rule $S(P,y)$, the population objective is
\begin{equation}
\label{eq:population_objective}
\mathcal{L}(\phi)
=
\mathbb{E}_{(C,Y)\sim P^*}
\left[
S\!\left(P^\phi(\cdot\mid C),Y\right)
\right].
\end{equation}

In the conditional setting the strictly propriety of scoring rule is satisfied if:
\[
\mathbb{E}_{Y\sim P^*(\cdot\mid c)}
S(P^*(\cdot\mid c),Y)
<
\mathbb{E}_{Y\sim P^*(\cdot\mid c)}
S(P^\phi(\cdot\mid c),Y)
\]
for every conditional distribution
$
P^\phi(\cdot\mid c)\neq P^*(\cdot\mid c).
$
Consequently, the unique minimizer of (\ref{eq:population_objective}) satisfies
\[
P^\phi(\cdot\mid c)=P^*(\cdot\mid c)
\qquad
\text{for almost every }c,
\]
which provides the statistical consistency of the learning objective.\\


Before using any scoring rule as a training loss, \cite{srforgan:srforgan} and \cite{background:scoring_rules} require:

\begin{enumerate}
    \label{ES_conditions}
\item $\beta\in(0,2)$: the negative-power kernel $k(x,y)=-\|x-y\|^\beta$ is (conditionally) strictly positive-definite only in this range, which is what makes $S_E^{(\beta)}$ \ref{eq:energy_score} strictly proper rather than merely proper.
\item Finite $\beta$-th moments: $\mathbb{E}_{Y\sim P^\star}\|Y\|^\beta<\infty$ and $\mathbb{E}_{X\sim P^\phi}\|X\|^\beta<\infty$, without this the score, and hence the objective, is not even be finite.
\item A \emph{well-specified} unbiased estimator of the population scoring rule.
\item i.i.d.\ training pairs $(c_j,y_j)\sim\Pi\times P^\star(\cdot|c)$ 
\item A generator admitting the reparameterisation trick, $X=G_\phi(z,c)$, $z\sim\mathcal N(0,I)$, differentiable in $\phi$, so $\nabla_\phi S_E$ can be estimated unbiasedly (App.~C, \cite{srforgan:srforgan}).
\end{enumerate}
The discussion and proofs of these conditions are met in out setting are provided in the appendix \ref{ch:condition_proofs}.\\


The original formulation of Pacchiardi et al. \cite{srforgan:srforgan} considers a prequential forecasting problem in which the conditioning variables are overlapping sliding windows extracted from a single observed time series. 
In contrast, the present work constructs the training dataset by simulating $J$ independent trajectories of the underlying stochastic differential equation. 
Each trajectory produces one condition--target pair
\[
(c_j,y_j),
\qquad
j=1,\ldots,J,
\]
where
\[
c_j
=
(X_0^j,\ldots,X_N^j)
\]
is the observed trajectory up to the forecasting horizon and
\[
y_j
=
X_{N+\tau}^j
\]
is the future realization.



Although different trajectories may be generated using different model parameters (for example different values of the volatility parameter in the Black--Scholes model or different Heston parameters), this does not violate the theoretical assumptions underlying scoring-rule minimization. Indeed, if the parameters $\vartheta$ of the underlying data generation process are independently 
sampled from a probability distribution $\Pi(\vartheta)$ before simulating each trajectory, every observation is generated according to a hierarchical model. In fact, the first step of the generative process is to sample the parameters $\vartheta_j$ for each trajectory $j$ from $\Pi$, and then to sample the condition--target pair $(c_j,y_j)$ from the conditional law $P(C,Y \mid\vartheta)$ defined by those parameters. Formally:
\[
\vartheta_j\sim\Pi,
\qquad
(c_j,y_j)\sim P(C,Y \mid\vartheta),
\]
which induces the marginal joint distribution:
\[
(c_1,y_1),\ldots,(c_J,y_J)
\stackrel{\mathrm{i.i.d.}}{\sim}
P(C,Y),
\]

So, even though each trajectory may have been simulated with different underlying stochastic parameters. The variation in the parameters is absorbed into the marginal data distribution and does not affect the strict propriety of the scoring rule, which is defined conditionally on the observed conditioning variable.\\




The empirical objective optimized during training is therefore
\begin{equation}
\label{eq:empirical_sr}
\hat{\mathcal{L}}(\phi)
=
\frac{1}{J}
\sum_{j=1}^{J}
S_E\!\left(
P^\phi(\cdot\mid c_j),
y_j
\right),
\end{equation}
which is the standard empirical estimator of (\ref{eq:population_objective}) for independent condition--target observations.


%----------------------------------------------------



The paper \cite{srforgan:srforgan} uses the Energy Score as its primary scoring rule. 
For $\beta \in (0,2)$, with $X, X' \sim P^\phi$ independent draws:
\begin{equation}
    \label{eq:energy_score}
S_E^{(\beta)}(P^\phi, y) = 2 \cdot \mathbb{E}\|X - y\|_2^\beta - \mathbb{E}\|X - X'\|_2^\beta
\end{equation}

This Energy Score is obtained from the general Kernel Score:
\begin{equation*}
    \label{eq:kernel_score} 
    S_k(P^\phi, y) = \mathbb{E}[k(X,X')] - 2\mathbb{E}[k(X,y)]
\end{equation*}

By choosing the negative power kernel $k(x,y) = -\|x - y\|_2^\beta$. 
The Kernel Score is strictly proper whenever $k$ is a positive-definite kernel and $-\|x - y\|_2^\beta$ for $\beta \in (0,2)$ is indeed positive-definite.\\

The Energy Score can be decomposed in two terms:
$$S_E^{(\beta)}(P^\phi, y) = \underbrace{2\mathbb{E}\|X - y\|_2^\beta}_{\text{accuracy term}} - \underbrace{\mathbb{E}\|X - X'\|_2^\beta}_{\text{spread reward}}$$
The accuracy term pulls generated samples toward the truth. The spread reward is subtracted, meaning the optimizer is explicitly penalized for producing a degenerate distribution where $X \approx X'$. 
This is the mechanism that prevents mode collapse: a collapsed distribution has $\|X - X'\| = 0$, contributing nothing to the spread reward, making the total score worse than a properly dispersed distribution. The strictly proper guarantee ensures the only global minimum is at $P^\phi = P^*$.\\

The paper fixes $\beta = 1$ throughout. In 1D, $S_E^{(1)}$ recovers exactly the CRPS (Continuous Ranked Probability Score), which is the standard evaluation metric in probabilistic weather forecasting:
\begin{equation}
    \label{eq:crps}
    \begin{aligned}
        \mathrm{CRPS}(F, y) =& \int_{-\infty}^{\infty} \big(F(x) - \mathbf{1}\{x \ge y\}\big)^2 \, dx\\
        &\text{ or alternatively }\\
      \mathrm{CRPS}(F, y)  =& \mathbb{E}|X - y| - \frac{1}{2}\mathbb{E}|X - X'|
    \end{aligned}
\end{equation}


$S_E$ is defined as an expectation over $P^\phi$, but is not possible to differentiate through a sampling operation.
The reparameterization trick resolves this. Because the generator defines:
$$X = G_\phi(z, \text{c}), \quad z \sim \mathcal{N}(0, I)$$
the expectation formula can be rewritten as:
$$\mathbb{E}_{X \sim P^\phi}[\|X - y\|] = \mathbb{E}_{z \sim \mathcal{N}}[\|G_\phi(z, c) - y\|]$$
The randomness now lives entirely in $z$, which is fixed before the forward pass. 
The generator $G_\phi$ is deterministic and differentiable with respect to its weights. Automatic differentiation libraries can therefore compute $\nabla_\phi S_E$ exactly.
The unbiased Monte Carlo estimator of the full gradient, as discussed in Appendix C of the paper \cite{srforgan:srforgan} is:
\begin{equation}
    \label{eq:energy_score_estimator}
    \hat{S}_E^{(\beta)}(\{x_i\}_{i=1}^m, y) = \frac{2}{m}\sum_{i=1}^m \|x_i - y\|^\beta - \frac{1}{m(m-1)}\sum_{\substack{i,k=1 \\ k \neq i}}^m \|x_i - x_k\|^\beta
\end{equation}

where each $x_i = G_\phi(z_i, c)$ with $z_i$ drawn independently. The $\frac{1}{m(m-1)}$ denominator excludes diagonal terms and makes this an unbiased estimator. 
Moreover, the paper proves that exchanging gradient and expectation is valid almost surely under mild conditions on the neural network architecture. This means that minimizing the empirical energy score via SGD gives unbiased gradient estimates of the true objective, so convergence results from standard SGD theory apply.



\begin{figure}[H]
    \centering
    \includegraphics[width=1.0\textwidth]{SRForGAN/images/srforgan_train_schema.png}
    \caption{SRForGAN training schema. The generator takes a noise vector and a conditioning variable to produce a forecast distribution, which is then scored against the true observation using the Energy Score. The critic is not used in this architecture, as the adversarial loss is replaced by the scoring rule objective.}
    \label{img:srforgan_train_schema}
\end{figure}


%
%Layer 5 — The prequential formulation: adapting to time series
%This is the paper's main theoretical contribution. For a temporal sequence $(y_1, \dots, y_T)$, the standard i.i.d. scoring rule objective (Eq. 6 in the paper) is:
%$$\min_\phi \mathbb{E}_{\theta \sim \Pi} \mathbb{E}_{y \sim P^*(\cdot|\theta)} S(P^\phi(\cdot|\theta), y)$$
%This assumes independent context-observation pairs $(\theta_i, y_i)$, which is not true for a time series. The paper replaces this with the prequential (predictive-sequential) score. Given a recorded window of length $T$, forecast each step $l$ ahead conditioned on the previous $k$ observations:
%$$\hat{\phi}_T(y_{1:T}) := \arg\min_\phi \sum_{t=k}^{T-l} S(P^\phi_{t+l}(\cdot|y_{t-k+1:t}),\; y_{t+l})$$
%This is literally how forecasting systems are evaluated in practice: for each time window in your training data, generate a conditional forecast and score it against the realization. The prequential score is then the sum of these one-at-a-time forecasting evaluations across the entire training sequence.
%Why is this still a proper scoring rule despite temporal dependence? The key is Theorem 2 (in the paper): if $S$ is strictly proper, then the prequential score $S_T$ is strictly proper for distributions over $y_{k+l:T}|y_{1:k+l-1}$ that are $k$-Markovian with lag $l$. The proof proceeds by tower property of expectations — propriety at each step propagates to the joint proper scoring of the full sequence. Intuitively, because each term $S(P^\phi_{t+l}(\cdot|y_{t-k+1:t}), y_{t+l})$ depends only on the local window of length $k$, and the generator is the same for all $t$, the sequential evaluation correctly scores the generator's ability to forecast the conditional distribution of the next step.



%Layer 6 — The consistency theorem (Theorem 3)
%This is the rigorous guarantee that makes the method scientifically sound. The empirical minimizer:
%$$\hat{\phi}_T(y_{1:T}) = \arg\min_\phi \frac{1}{T-l-k+1}\sum_{t=k}^{T-l} S(P^\phi(\cdot|y_{t-k+1:t}), y_{t+l})$$
%converges to $\phi^*_T$, the minimizer of the true expected prequential score, as $T \to \infty$, under four assumptions:
%A1 — Compactness: $\Phi$ compact. Standard regularity condition.
%A2 — Uniqueness: The true minimizer is isolated. Lemma 12 shows this is satisfied when $S$ is strictly proper and the model is well-specified (i.e., there exists $\phi^*$ such that $P^{\phi^*} = P^*$).
%A3 — Asymptotic stationarity: The marginal distribution of $(y_{t-k+1:t+l})$ averaged over time converges weakly. This is satisfied for stationary processes like GBM, which has stationary increments.
%A4 — Mixing and moment boundedness: The process is either $\alpha$-mixing or $\phi$-mixing with appropriate decay rates. GBM satisfies $\phi$-mixing since increments are independent. The moment bound for the Energy Score reduces to requiring $\mathbb{E}\|y_{t+l}\|^\beta < \infty$, which holds for log-normal GBM increments.
%The proof combines two classical results: the Uniform Law of Large Numbers for dependent data (Pötscher and Prucha 1989) applied to the sum of scoring rule evaluations, and a generic consistency result (Theorem 5.1, Skouras 1998) which promotes uniform convergence of the objective to convergence of the minimizer. The key insight is that even though consecutive terms in the sum are correlated (since the windows overlap), the mixing condition ensures this dependence decays fast enough that the ULLN still holds. This is fundamentally different from GAN training, which has no analogous consistency result for dependent data.





% ---------------------------------------------------------------
% Algorithm 1 — Unbiased Energy Score Estimate
% ---------------------------------------------------------------
\begin{algorithm}[H]
\caption{Unbiased Energy Score Estimate $\hat{S}_E$}
\label{alg:energy_score}
\begin{algorithmic}[1]
\REQUIRE Generator $G_{\phi}$, condition $c \in \mathbb{R}^l$,
         observation $y \in \mathbb{R}^d$,
         exponent $\beta \in (0,2)$
\STATE Draw $m$ independent noise vectors $\{z_j\}_{j=1}^m \sim \mathcal{N}(0, I)$
\STATE Generate $m$ samples $\{x_j\}_{j=1}^m \subset \mathbb{R}^d$ where $x_j = G_{\phi}(z_j, c)$


\STATE \textbf{// Term 1: } $\frac{2}{m}\sum_{j=1}^{m}\|x_j - y\|^{\beta}$
\STATE $T_1 \leftarrow \dfrac{2}{m}\displaystyle\sum_{j=1}^{m}\|x_j - y\|^{\beta}$

\STATE \textbf{// Term 2: }
       $\frac{1}{m(m-1)}\sum_{j \neq k}\|x_j - x_k\|^{\beta}$
\STATE $T_2 \leftarrow \dfrac{1}{m(m-1)}\displaystyle\sum_{\substack{j,k=1 \\ j \neq k}}^{m}\|x_j - x_k\|^{\beta}$

\RETURN $\hat{S}_E \leftarrow T_1 - T_2$
\end{algorithmic}
\end{algorithm}

It is important to note that the available code implementation of the scoring rule includes also the 
kernel score, which is a generalization of the energy score. The kernel score is defined as:
\begin{equation}
    \label{eq:kernel_score_estimator}
    \hat{S}_k(\{x_i\}_{i=1}^m, y) = 
    \frac{1}{m(m-1)}\sum_{\substack{i,k=1 \\ k \neq i}}^m k(x_i, x_k) 
    - \frac{2}{m}\sum_{i=1}^m k(x_i, y), 
    \quad k(x,x') = \exp\!\left(-\frac{\|x - x'\|^2}{2\gamma^2}\right)
\end{equation}
where $\gamma$ is a bandwidth parameter. The kernel score is also strictly proper and can be used as an alternative to the energy score, but the paper focuses on the energy score for its simplicity and interpretability in terms of accuracy and spread.\\
One of the major proposition arising from \cite{srforgan:srforgan} is that the sum of the energy score and the kernel score is also a strictly proper scoring rule, which they call the energy-kernel score. This allows for a flexible combination of both scores in the generator's loss function, potentially leveraging the strengths of each to further improve forecasting performance.
The energy-kernel score is defined as:
\begin{equation}
    \label{eq:energy_kernel_score_estimator}
    \hat{S}_{E+k}^{(\beta)}(\{x_i\}_{i=1}^m, y) = 
    \hat{S}_E^{(\beta)}(\{x_i\}_{i=1}^m, y) 
    + \hat{S}_k(\{x_i\}_{i=1}^m, y)
\end{equation}
where $\hat{S}_E^{(\beta)}$ is the energy score estimator from Algorithm \ref{alg:energy_score} and $\hat{S}_k$ is the kernel score estimator from Equation \ref{eq:kernel_score_estimator}.


% ---------------------------------------------------------------
% Algorithm 2 — MySRForGAN Training
% ---------------------------------------------------------------
\begin{algorithm}[H]
\caption{\textsc{SRForGAN} Training via Energy Score Minimisation}
\label{alg:srforgan_train}
\begin{algorithmic}[1]
\REQUIRE
    Training set $\mathcal{D}_{\text{train}} = \{(y_{t+l},\, y_{t-k+1:t})\}_{t=k}^{T-l}$,
    validation set $\mathcal{D}_{\text{val}}$ (optional),
    generator $G_{\phi}$ with parameters $\phi$,
    scoring rule SR $\in \{\text{energy}, \text{kernel}, \text{energy-kernel}\}$,
    number of generator draws $m$,
    learning rate $\eta$,
    maximum epochs $E_{\max}$,
    early-stopping patience $P$,
    minimum improvement threshold $\delta_{\min}$
\ENSURE Trained generator parameters $\phi^*$

\STATE Initialise Adam optimiser with parameters $\phi$ and learning rate $\eta$
\STATE $\ell^*_{\text{val}} \leftarrow +\infty$;\quad
       $p \leftarrow 0$;\quad
       $\phi^* \leftarrow \phi$

\FOR{$e = 1, \ldots, E_{\max}$}

    \STATE \textbf{// --- Training pass ---}
    \STATE Shuffle $\mathcal{D}_{\text{train}}$ into mini-batches
           $\mathcal{B}_1, \ldots, \mathcal{B}_{N_B}$ of size $B$
    \STATE $\ell_{\text{train}} \leftarrow 0$

    \FOR{each mini-batch $(Y, C) \in \mathcal{B}_1, \ldots, \mathcal{B}_{N_B}$}
        \STATE $\hat{S}_E^B = 0$
        \FOR{$b = 1$ to $B$ samples in the mini-batch}
            \STATE \text{// Compute for sample $b$ of the batch  $\hat{S}_E^b $}
            \hfill\textit{// Algorithm~\ref{alg:energy_score}}
            \STATE $\hat{S}_E^B \leftarrow \hat{S}_E^B + \hat{S}_E^b$
        \ENDFOR
        \STATE $\hat{\mathcal{L}} \leftarrow
               \textsc{SRLoss}(C,\, Y,\, G_{\phi},) = \frac{1}{N_B} \hat{S}_E^B$
               \hfill\textit{// Estimator~\ref{eq:empirical_sr}}
        \STATE $\phi \leftarrow \phi - \eta\,\nabla_{\phi}\,\hat{\mathcal{L}}$
               \hfill\textit{// Adam update, no critic needed}
        \STATE $\ell_{\text{train}} \leftarrow \ell_{\text{train}} + \hat{\mathcal{L}}$
    \ENDFOR
    \STATE $\ell_{\text{train}} \leftarrow \ell_{\text{train}} / N_B$

    \STATE \textbf{// --- Validation pass (no gradient) ---}
    \IF{$\mathcal{D}_{\text{val}}$ is provided}
        \STATE $\ell_{\text{val}} \leftarrow \textsc{SRLoss}$ evaluated on
               $\mathcal{D}_{\text{val}}$ 
        \IF{$\ell_{\text{val}} < \ell^*_{\text{val}} - \delta_{\min}$}
            \STATE $\ell^*_{\text{val}} \leftarrow \ell_{\text{val}}$;\quad
                   $p \leftarrow 0$;\quad
                   $\phi^* \leftarrow \phi$
                   \hfill\textit{// save best checkpoint}
        \ELSE
            \STATE $p \leftarrow p + 1$
            \IF{$p \geq P$}
                \STATE \textbf{break}
                       \hfill\textit{// early stopping}
            \ENDIF
        \ENDIF
    \ENDIF

\ENDFOR

\RETURN $\phi^*$
\end{algorithmic}
\end{algorithm}





% ── local colour definitions ─────────────────────────────────────────────────
\definecolor{cganblue}{RGB}{30,100,180}
\definecolor{wganred}{RGB}{180,40,40}
\definecolor{commentgray}{RGB}{110,110,110}
 
% ── local math helpers ───────────────────────────────────────────────────────
\newcommand{\Loss}{\mathcal{L}}
\newcommand{\ES}{\hat{S}_{E}}
 
% ── line-comment for the algorithmic package ─────────────────────────────────
% algorithmic has no standalone comment line; emulate it with \STATE
\newcommand{\LINECOMMENT}[1]{%
  \STATE{\color{commentgray}\small$\triangleright$\;\textit{#1}}}
 
% override inline comment style
\renewcommand{\algorithmiccomment}[1]{%
  \hfill{\color{commentgray}\small$\triangleright$\;\textit{#1}}}
 
% ─────────────────────────────────────────────────────────────────────────────

Additionally, the paper states in the conclusion that the combination of adversarial and SR minimization would be beneficial and a promising path to explore. So an additional hybrid loss term is defined by adding the Energy Score to both the Wasserstein adversarial loss and classical Jensen-Shannon adversarial loss. 
In particular, the generator loss using the hybrid approach is defined as:
\begin{equation}
    \label{eq:hybrid_wasserstein_loss}
    \mathcal{L}_G = \underbrace{-\mathbb{E}[C(G(c,z))]}_{\text{Wasserstein signal}} + \lambda \cdot \underbrace{\hat{S}_E}_{\text{Energy Score}}
\end{equation}


 
The hybrid models, which use both the adversarial losses (Binary Cross-Entropy and Wasserstein distance) and the energy score, inherit the same energy score method from \ref{alg:energy_score}.\\
The hybrid Forgan uses as adversarial term the Binary Cross-Entropy loss and the
 full hybrid generator loss for one mini-batch $B$ is:
 
\begin{equation}
  \Loss_{G}^{\mathrm{CGAN}}
  =
  \underbrace{
    -\frac{1}{B} \sum_{b=1}^B \log D(G(z^{(b)}, c^{(b)}))
  }_{\text{BCE adversarial loss}}
  \;+\;
  \underbrace{
    \frac{\lambda}{B}\;
    \sum_{b=1}^B \ES\!\left(\{\hat{x}_m\}_{m=1}^{M},\,\mathbf{Y}\right)
  }_{\text{Energy Score (\ref{alg:energy_score})}}
  \label{eq:cgan_hybrid}
\end{equation}
 
\begin{algorithm}[H]
\caption{{\color{cganblue}\texttt{ForGAN}} --- hybrid generator update
  (called when $\mathrm{step} \bmod n_{\mathrm{critic}} = 0$)}
\begin{algorithmic}[1]
\REQUIRE Generator $G_{\phi}$, Discriminator $D_{\psi}$
\REQUIRE Condition batch $\mathbf{c}$, true targets $\mathbf{y}$,
         real labels $\mathbf{1}$, samples $M$,
         ES weight $\lambda$, noise dim $N_z$
\ENSURE  Updated generator parameters $\phi$


\LINECOMMENT{Draw M independent noise vectors; generate M samples}
\STATE $\mathrm{samples} \leftarrow [\;]$
\FOR{$m = 1$ \textbf{to} $M$}
  \STATE $\mathbf{z}_m \sim \mathcal{N}(\mathbf{0},\,\mathbf{I}_{N_z})$
  \COMMENT{fresh draw each iteration}
  \STATE $\hat{y}_m \leftarrow G_{\phi}(\mathbf{c},\,\mathbf{z}_m)$
  \COMMENT{shape $[B \times D]$}
  \STATE $\mathrm{samples}.\mathrm{append}(\hat{y}_m)$
\ENDFOR
 
\LINECOMMENT{Adversarial term: pass only the first sample through D}
\STATE $\mathrm{probs}
  \leftarrow D_{\psi}\!\left(\mathrm{samples}[1],\,\mathbf{c}\right)$
\STATE $\Loss_{\mathrm{adv}}
  \leftarrow \mathrm{BCE}(\mathrm{probs},\;\mathbf{1})$
  \COMMENT{generator wants $D$ to output 1 on fakes}
 
\LINECOMMENT{Energy Score term: all M samples contribute}
\STATE $\Loss_{\mathrm{ES}}
  \leftarrow
  \textsc{EnergyScoreLoss}(\mathrm{samples},\;\mathbf{y})$
 
\LINECOMMENT{Combine, backpropagate once, update}
\STATE $\Loss_{G}
  \leftarrow \Loss_{\mathrm{adv}} + \lambda \cdot \Loss_{\mathrm{ES}}$
\STATE $\mathrm{Backward}(\Loss_{G})$
  \COMMENT{gradients from both terms accumulate in one pass}
\STATE $\mathrm{Step}(\phi)$
\RETURN $\Loss_{G}$
\end{algorithmic}
\end{algorithm}


The adversarial term is the Wasserstein generator loss
$-\mathbb{E}[D(G(z, c))]$ so in this case no labels are needed since the critic
output is an unbounded real score.
 
\begin{equation}
  \Loss_{G}^{\mathrm{CWGAN}}
  =
  \underbrace{
    -\frac{1}{B} \sum_{i=1}^B D(G(z^{(b)}, c^{(b)}))}_{\text{Wasserstein generator loss}}
  \;+\;
  \underbrace{
    \frac{\lambda}{B}\;
    \sum_{i=1}^B \ES\!\left(\{\hat{x}_m\}_{m=1}^{M},\,\mathbf{Y}\right)
  }_{\text{Energy Score (\ref{alg:energy_score})}}
  \label{eq:cwgan_hybrid}
\end{equation}
 



\begin{algorithm}[H]
\caption{{\color{wganred}\texttt{ForWGAN}} --- hybrid generator update}
\begin{algorithmic}[1]
\REQUIRE Generator $G_{\phi}$, Critic $C_{\psi}$
\REQUIRE Condition batch $\mathbf{c}$, true targets $\mathbf{y}$,
         samples $M$, ES weight $\lambda$, noise dim $N_z$
\ENSURE  Updated generator parameters $\phi$
 
 
\LINECOMMENT{Draw M independent noise vectors; generate M samples}
\STATE $\mathrm{samples} \leftarrow [\;]$
\FOR{$m = 1$ \textbf{to} $M$}
  \STATE $\mathbf{z}_m \sim \mathcal{N}(\mathbf{0},\,\mathbf{I}_{N_z})$
  \STATE $\hat{y}_m \leftarrow G_{\phi}(\mathbf{c},\,\mathbf{z}_m)$
  \COMMENT{shape $[B \times D]$}
  \STATE $\mathrm{samples}.\mathrm{append}(\hat{y}_m)$
\ENDFOR
 
\LINECOMMENT{Wasserstein term: no labels; maximise critic score on fakes}
\STATE $\Loss_{\mathrm{W}}
  \leftarrow
  -C_{\psi}\!\left(\mathrm{samples}[1],\,\mathbf{c}\right).\mathrm{mean}()$
  \COMMENT{no target labels required}
 
\LINECOMMENT{Energy Score term: all M samples contribute}
\STATE $\Loss_{\mathrm{ES}}
  \leftarrow
  \textsc{EnergyScoreLoss}(\mathrm{samples},\;\mathbf{y})$
 
\LINECOMMENT{Combine, backpropagate once, update}
\STATE $\Loss_{G}
  \leftarrow \Loss_{\mathrm{W}} + \lambda \cdot \Loss_{\mathrm{ES}}$
\STATE $\mathrm{Backward}(\Loss_{G})$
\STATE $\mathrm{Step}(\phi)$
\RETURN $\Loss_{G}$
\end{algorithmic}
\end{algorithm}
 

The critic and the discriminator are trained on real and fake samples, respectively, and do not interact with the Energy Score in any way,
so they follow the standard adversarial training procedure already described in the previous chapter \ref{alg: cgan}, \ref{alg:cwgan_gp}.



\section{Performance metrics and statistical tests}
\label{sec:performance_metrics}
In the following section the performance metrics for probability distributions used to evaluate the proposed models will be described in details.
Moreover, the Kolmogorov-Smirnov test will be presented in order to evaluate the performance of the model in generating the future probability distribution of the asset price.\\
\subsection{Distance metrics}
When dealing with discretized probability distributions represented as vectors,
where $P = [p_1, p_2, \dots, p_k]$ and $Q = [q_1, q_2, \dots, q_k]$ are probability mass functions defined over the same binning, 
different metrics capture different geometric and information-theoretic properties.\\

The Total Variation (TV) distance is the $L_1$ distance between the vectors, defined as:
\begin{equation}
    \label{eq: TV_distance}
    \delta(P, Q) = \frac{1}{2} \sum |p_i - q_i|
\end{equation}
It represents the maximum possible difference between the probabilities that the two distributions assign to the same event,
but it is "bin-blind," meaning it doesn't care if the probability mass is moved to a neighboring bin or an outlier bin.\\
The Hellinger distance, defined as:
\begin{equation}
    H(P, Q) = \frac{1}{\sqrt{2}} \sqrt{\sum (\sqrt{p_i} - \sqrt{q_i})^2}
\end{equation}
It is a bounded metric ($[0, 1]$) which is more sensitive to differences in the "tails" of the distribution than TV-distance.\\
The Jensen-Shannon (JS) Divergence is a symmetric, smoothed version of the Kullback-Leibler divergence, calculated as:
\begin{equation}
    JS(P, Q) = \frac{1}{2}KL(P || M) + \frac{1}{2}KL(Q || M)
\end{equation}
where $M = \frac{1}{2}(P + Q)$. While JS is a popular metric for GANs, it suffers from the "disjoint support" problem where it reaches a constant maximum if the distributions do not overlap.\\
In contrast, the Wasserstein distance (specifically the 1D version for discretized vectors) can be computed as the $L_1$ distance between the Cumulative Distribution Functions (CDFs): 
\begin{equation}
    W_1(P, Q) = \sum |CDF_P(i) - CDF_Q(i)|
\end{equation}
Unlike the other three, the Wasserstein distance is "geometry-aware," meaning it incorporates the underlying distance between bins, making it far more robust for Monte Carlo simulations where we care about how far the generated values are shifted from the true values.

\subsection{The Kolmogorov-Smirnov test}
The Kolmogorov-Smirnov (KS) as defined in \cite{background:ks_test}, is a goodness of fit test that evaluates whether an empirical distribution is consistent with a fully specified continuous reference distribution. 
Let $F(x)$ denote the hypothesized cumulative distribution function (CDF), 
and let $F_{n}(x)$ be the empirical CDF constructed from a sample of size $n$. 
The KS statistic is defined as:
\begin{equation}
    D_{n}=\sup_{x}\lvert F_{n}(x)-F(x)\rvert
\end{equation}
representing the maximum vertical deviation between the empirical and theoretical CDFs.
Under the null hypothesis $H_{0}: F_{n}(x)=F(x)$,
the scaled distribution $\sqrt{n}D_{n}$ converges to the Kolmogorov distribution, enabling computation of p-values by using its CDF:
\begin{equation*}
    P(\sqrt{n}D_{n}\le t)=1-2\sum_{k=1}^{\infty}(-1)^{k-1}e^{-2k^{2}t^{2}}
\end{equation*}
Large values of $D_{n}$ lead to rejection of $H_{0}$, indicating
incompatibility between the empirical and theoretical distributions.\\

When applying the KS framework to generated outputs from the proposed GAN, the empirical object is a vector of bin probabilities
$(p_{1},\dots,p_{K})$.  These represent a discretized approximation of the probability distribution. To construct the empirical CDF, 
the cumulative probabilities are computed as:
\begin{equation*}
    F_{\text{GAN}}(x_{k})=\sum_{i=1}^{k}p_{i}
\end{equation*}

where $x_k$ indexes the bin boundaries. For comparison with a normal distribution, the theoretical CDF
$F_{\mathcal N}(x_{k})$ is computed at the same grid. The KS statistic becomes:
\begin{equation*}
    D=\max_{k}\lvert F_{\text{GAN}}(x_{k})-F_{\mathcal N}(x_{k})\rvert
\end{equation*}

This yields a valid discrepancy measure, although interpreting p-values requires specifying an effective sample size.
Nonetheless, $D$ remains a rigorous, distribution-free indicator of alignment between the GAN-generated distribution and the target normal CDF.
Actually, in order to compute the p\-value of the test, the survival function of the Kolmogorov distribution is used $P(\sqrt{n}D_n \geq \sqrt{n}\hat{D_n} | H_0)$ 
where $\hat{D_n}$ is the observed test statistic. 
The effective sample size $n$ has been set equal to the number of bins $B$ used to discretize the distribution.


%\include{Background/Background}
\chapter{Experiments with Black-Scholes dynamics}
\label{sec:bs_experiments}

The Black--Scholes framework rests on: (i) a frictionless, arbitrage-free market with continuous trading and no transaction costs; (ii) a risky asset $S_t$ with \emph{constant} instantaneous drift $\mu$ and volatility $\sigma>0$; (iii) a constant risk-free rate $r$; (iv) log-normally distributed asset prices, driven by a single standard Brownian motion $W_t$ under the physical measure $\mathbb{P}$; and (v) no dividends, jumps, or stochastic volatility.

Under $\mathbb{P}$, the asset price follows a geometric Brownian motion (GBM):
\begin{equation}
dS_t = \mu S_t\, dt + \sigma S_t\, dW_t, \qquad S_0 > 0.
\end{equation}
Under the risk-neutral measure $\mathbb{Q}$, obtained via Girsanov's theorem with market price of risk $\lambda = (\mu-r)/\sigma$, the drift is replaced by $r$:
\begin{equation}
dS_t = r S_t\, dt + \sigma S_t\, dW_t^{\mathbb{Q}}.
\end{equation}


Applying It\^o's lemma to $\ln S_t$ yields the exact solution
\begin{equation}
S_t = S_0 \exp\left[\left(r - \tfrac{1}{2}\sigma^2\right)t + \sigma W_t^{\mathbb{Q}}\right],
\end{equation}
so that $\ln S_t \mid S_0 \sim \mathcal{N}\!\left(\ln S_0 + (r-\tfrac12\sigma^2)t,\; \sigma^2 t\right)$: log-returns are Gaussian and i.i.d.\ over disjoint intervals, and $S_t$ is log-normal.


Sample paths are continuous (no jumps), Markovian, and self-similar in the sense $S_{t+\Delta}/S_t \stackrel{d}{=} S_\Delta/S_0$ with instantaneous variance constant.


\section{Data and Training procedure}

For each trajectory $j = 1, ..., J$ we define 
the SDE of the log-price $X_t^j$ to be modelled using the Black-Scholes framework where the only source of uncertainty is given by the Brownian motion $W_t$:
\begin{equation}
    \label{log-price SDE}
    \begin{aligned}
       &dX_t^j = -\frac{1}{2}\sigma_j^2 dt + \sigma_j dW_t\\
         &\text{with } X_0^j = \log(S_0^j)
    \end{aligned}
    \end{equation} 

Where $S_0^j$ is the raw price at time $0$ of trajectory $j$, $\sigma_j$ is the volatility of trajectory $j$ and $W_t$ is a standard Brownian motion.
In this setting we are emulating the scenario of working under the risk-neutral measure $\mathbb{Q}$ with zero interest rate and zero dividend yield.\\
$X_0^j$ and $\sigma_j$ are randomly sampled for each trajectory such that 
$X_0^j \in [-1, 1]$ and $\sigma_j \in (0.0, 1.2]$.\\

In this context our condition $c_j$ for a general data observation $j$ is given by
a vector containing the trajectory $c_j=[X_0^j, X_{dt}^j, X_{2dt}^j, ..., X_{Ndt}^j]$.

 The SDE is discretized over a uniform time space $N dt = T$ obtaining the following form
 adopted for the simulation of the trajectories:
 \begin{equation}
    \label{eq: discretized_SDE}
    X_{t+dt}^j = X_t^j -\frac{1}{2}\sigma_j^2 dt + \sigma_j \sqrt{dt} Z_t
 \end{equation}
 Where $Z_t \sim \mathcal{N}(0,1)$.\\
The procedure to build the condition data matrix $C$ is summarized in Algorithm \ref{alg:y condition}.

\begin{algorithm}[H]
\caption{Algorithm for the condition data matrix $C$ creation}
\label{alg:y condition}
\begin{algorithmic}[1]

\STATE Define the number of simulation $J$.
\STATE Initialize an empty data matrix $C$ of shape $J\times N$.
\STATE Set iteration counter $j \leftarrow 0$.
\REPEAT
    \STATE Randomly sample $\sigma_j \in (0.1, 1.2)$ with $X_0^j = 0.0$.
    \STATE Evolve $X_t^j$ based on the discretized SDE for $0, dt, ..., Ndt$ steps.
    \STATE Store the trajectory $[X_0^j, X_{dt}^j, X_{2dt}^j, ..., X_{Ndt}^j]$ in the $j$-th row of $C$.
    \STATE Increment counter $j \leftarrow j + 1$.
\UNTIL{$j = J$}
\RETURN $C$
\end{algorithmic}
\end{algorithm}


Since the ForGAN architecture is adopted, which is trained to generate the next value of the trajectory,
the output data $X$ is a $J \times 1$ vector where each element is the $n$-th step ahead value of the trajectory, 
so the output is a single value and not a probability distribution during training time. 
In this case, the target output data matrix $X$ is computed by evolving each trajectory for the $\tau$-th step using the discretized SDE 
described in  \ref{eq: discretized_SDE} and storing the resulting value in the $j$-th position of $X$.\\

Now that the training data are defined the ForGAN model can be trained by using both the classical BCE loss, using algorithm \ref{alg: cgan}
and the Wasserstein loss with gradient penalty using algorithm \ref{alg:cwgan_gp}.
 




After the training phase, the ForGAN is able to generate the future distribution by repeatedly sampling the next value of the trajectory by using different noise realizations, as described in \ref{alg: forgan-inference}.\\
At the end of the condition trajectory, so at time $T=N dt$, we have the following information: the final log-price $X_{Ndt}^j$ of each trajectory $j$ and the corresponding volatility $\sigma_j$.
The goal of the Conditional GAN is to learn the distribution of final values of the stochastic process at next time horizon $T + \tau$. 
A theoretical closed form can be used as ground truth. It is given by a normal distribution:
\begin{equation}
    \begin{aligned}
    X_{Ndt + \tau}^j & \quad \sim \quad \mathcal{N}(X_{Ndt}^j - \frac{1}{2}\sigma_j^2 \Delta t, \sigma_j \sqrt{\Delta t})\\
    \text{where } &\Delta t = dt \times \tau \text{ with $\tau$ the number of days ahead. }
    \end{aligned}
\end{equation}


In order to provide a comparison between the proposed ForGAN architectures and 
the theoretical solution of the GBM processes, the target distribution $x$ of the next value $c_{T + \tau}$ of the process is computed by using the Monte Carlo algorithm described in \ref{alg:mc_terminal_dist}.
This is required to properly evaluate the model's performance because since the ForGAN model provides the distribution of the next value by repeated sampling, the vector of those samplings needs to be converted into a probability distribution. 
So, it is required a uniform way to compute the bins values used to build the empirical distribution.\\
In fact, after the Monte Carlo simulations are performed, the resulting terminal values are used to build the empirical distribution of the next value $c_{T + \tau}$ by using $\sqrt{n_{mc\_simulations}}$ as rule of thumb to compute the optimal number of bins for the histogram. 
The same bins are then used to compute the empirical distribution of the generated samples from the ForGAN model.\\
Alternatively, the Freedman-Diaconis rule can be used \cite{forgan:freedman_bins} to compute the optimal bin width for complex target distributions.
The Freedman-Diaconis rule provides a robust statistical method for determining optimal histogram bin width by minimizing the integrated mean squared error. 
Unlike methods sensitive to outliers, it utilizes the interquartile range (IQR) to account for data variability. 
The bin width $h$ is calculated as:$h = 2 \cdot \frac{\text{IQR}(x)}{\sqrt[3]{n}}$ where $n$ represents the number of observations.
By prioritizing the spread of the central 50\% of the distribution, 
the algorithm ensures that the resulting binning structure accurately reflects the underlying density of the empirical data without over-smoothing or producing noise-induced artifacts.\\

The algorithm to compute the theoretical distribution $p(x|c)$ given Monte Carlo simulations and the ForGAN's generated distribution are as follows:
\begin{algorithm}[H]
\caption{Monte Carlo generation of terminal distributions}
\label{alg:mc_terminal_dist}
\begin{algorithmic}[1]
\REQUIRE Number of trajectories $J$, steps ahead $N$, and MC simulations $K$, final values from which to compute the distribution $\mathbf{C}_T \in \mathbb{R}^J$, drift $\boldsymbol{\mu} \in \mathbb{R}^J$, volatility $\boldsymbol{\sigma} \in \mathbb{R}^J$, and time increment $ dt$.
\STATE \COMMENT{Generate random standard normal shocks for all trajectory $j$ and simulation $k$}
\STATE $\mathbf{Z} \leftarrow \text{Tensor}(J, K)$ where $Z_{j, k} \sim \mathcal{N}(0,1)$
\STATE $\Delta t = dt * N$
\FOR{each trajectory $j \in \{1, \dots, J\}$}
    \STATE \COMMENT{Calculate the deterministic drift component per step}
    \STATE $\mathbf{d_j} \leftarrow (\boldsymbol{\mu_j} - 0.5 \boldsymbol{\sigma_j}^2) \Delta t$ 
    \STATE \COMMENT{Compute stochastic shocks for simulation $k$}
    \STATE $\mathbf{S}_{k}^j \leftarrow \sigma_j \cdot Z_{j, k} \cdot \sqrt{\Delta t}$
    \STATE \COMMENT{Calculate increments for each simulation $k$}
    \STATE $\boldsymbol{\delta}_{k}^j \leftarrow d_j + S_{k}^j$ 
    \FOR{each simulation $k \in \{1, \dots, K\}$}
        \STATE $X_{k}^j \leftarrow C_{T,j} + \delta_{k}^j$
    \ENDFOR
\ENDFOR
\STATE $\mathbf{X} \leftarrow \text{Matrix}(J, K)$
\STATE Compute $\mathbf{D} \leftarrow \quad$ the normalized histogram of $\mathbf{X}$ to estimate the empirical distribution for each trajectory $j$ using $n_{bins} = \sqrt{K}$ or alternatively the method in \cite{forgan:freedman_bins}.
\RETURN $\mathbf{D}$, $\mathbf{bins}$ \COMMENT{The empirical probability distributions for each trajectory}
\end{algorithmic}
\end{algorithm}






\begin{algorithm}[H]
\caption{ForGAN distribution inference}
\label{alg: forgan-inference}
\begin{algorithmic}[1]
\REQUIRE Trained generator parameters $\theta_G$, conditioning history $c_j$, number of samples $V$
\STATE Initialize an empty list $\mathcal{S} \leftarrow \emptyset$ to store generated samples.
\FOR{$i = 1$ \TO $V$}
    \STATE Sample latent noise $z^{(i)} \sim p_z$ \COMMENT{Independent noise draw, same $c_t$}
    \STATE $\hat{x}_{t+h}^{(i)} \leftarrow G(z^{(i)}, c_j)$ \COMMENT{Generate candidate forecast}
    \STATE $\mathcal{S} \leftarrow \mathcal{S} \cup \{\hat{x}_{t+h}^{(i)}\}$
\ENDFOR
\RETURN $\mathcal{S}$ \COMMENT{ list of generated samples for empirical distribution of trajectory $j$}
\end{algorithmic}
\end{algorithm}

After having obtained $D$ from \ref{alg:mc_terminal_dist} and having computed the ForGAN inferred distribution by normalizing the histogram of $\mathcal{S}$ obtained from \ref{alg: forgan-inference} by using the same $bins$ from \ref{alg:mc_terminal_dist},
the two distributions can be compared by using all the metrics described in \ref{sec:performance_metrics} like Jensen-Shannon divergence, Wasserstein distance, Total Variation distance and statistical tests like the Kolmogorov-Smirnov test.


\section{ForGAN results with variable volatility trajectories}
\label{sec:bs experiment}
The following setup for the experiments has been used:\\
\begin{itemize}
    \item \textbf{Initial price range}: $X_0 = 0.0$
    \item \textbf{Drift range}: $\mu = 0.0$
    \item \textbf{Volatility range}: $\sigma \in (0.1, 1.2)$
    \item \textbf{Time horizon}: $T = 1.0 $ yearly horizon
    \item \textbf{Number of time steps}: $N = 252$
    \item \textbf{Number of simulated paths}: $J = 100{,}000$
    \item \textbf{n-steps ahead}: $\tau = 1$
\end{itemize}

The experimental setup is based on a data generation process designed to emulate realistic financial price dynamics exploiting Geometric Brownian motion described in the before in previous section. 
The asset price trajectories are simulated over a fixed time horizon of one year, discretized into $N=252$ time steps corresponding to daily observations.
For each experiment, $J=100{,}000$ independent paths are generated, ensuring sufficient sample diversity and statistical robustness.
The initial asset price $X_0$ is set for each trajectory to $0.0$, while the drift parameter $\mu$ is fixed at zero to isolate the model's ability to learn volatility-driven dynamics.
The volatility parameter $\sigma$ is sampled from a wide range $(0.1, 1.2)$ to capture varying market conditions, from near-deterministic to highly stochastic regimes.
The distribution to be learned refers to the $1$-th step ahead observation, that is, the distribution of asset prices at time $T + 1$ (or $N dt + 1$) given information up to time $T$. 
Learning distributions for longer horizons does not provide additional insights since the shape of the distribution remains Gaussian, with the only difference being a larger variance.\\

The experiment has been conducted by using the following architecture and hyperparameters:
\begin{table}[H]
\centering
\caption{ForGAN Architecture and Hyperparameters Summary}
\label{tab:forgan_summary}
\begin{tabular}{lll}
\hline
\textbf{Parameter}         & \textbf{Generator}          & \textbf{Critic (Discriminator)} \\ \hline
Input Dimension            & 64 (Latent) + 252 (Cond.)   & 1 (Data) + 252 (Cond.)          \\
Hidden Layers              & [512, 256, 128, 64]         & [512, 256, 256, 128, 64]        \\
Output Dimension           & 1                           & 1                               \\
Activation Function        & ReLU                        & Leaky ReLU                      \\
Normalization              & Batch Norm                  & None                      \\
Dropout                    & 0.0                         & 0.0                             \\
\multicolumn{3}{c}{\textbf{Global Training Settings}}                                     \\ \hline
Learning Rate              & \multicolumn{2}{l}{0.0005}                                      \\
Max Epochs                 & \multicolumn{2}{l}{100}                                       \\
\end{tabular}
\end{table}

The Kolmogorov-Smirnov test performed on 100000 real and generated probability distributions shows that
the resulting p-values are accepted 25\% of time.\\
Moreover, the average Jensen-Shannon divergence over the test samples is around 0.49,
suggesting poor performance since the Jensen-Shannon divergence computed by using the logarithm in base 2
is bounded in the range $[0, 1]$.

\begin{table}[h]
\centering
\begin{tabular}{l c}
\hline
\textbf{Metric} & \textbf{Value} \\
\hline
KS Test Acceptance (\%) & 25.00 \\
EMD Distance & 0.00855 \\
JS Distance & 0.4378 \\
Hellinger Distance & 0.4001 \\
TV Distance & 0.4129 \\
\hline
\end{tabular}
\caption{ForGAN performance metrics averaged over 100000 samples}
\end{table}

These results strongly suggest that it is worth trying to improve the performances of the proposed ForGAN by adopting the Wasserstein distance instead of the BCE loss.\\
In order to visually assess the capabilities of the model,
different samples have been drawn such that every sampling parameter is fixed except for the volatility $\sigma$.
In particular, six different samples are shown below where
the starting values is the same $X_0 = 0.0$, $\mu = 0.0$ while the volatility have different values for each sample 
$\sigma = [0.3, 0.5, 0.7, 0.9]$. From the results can be notices that also this new architecture
struggles to effectively capture the variance of the target distribution when the volatility parameter is high.

\begin{figure}[H]
    \centering
    \includegraphics[width=1.0\textwidth]{Background/images/samples_forgan_simple.png}
    \caption{Samples obtained using fixed $X_0$ and $\mu$ while varying $\sigma = [0.3, 0.5, 0.7, 0.9]$.}
    \label{forgan_simple_sample_sigma}
\end{figure}



In case of the ForWGAN architecture experiments 
just the training procedure and the architecture differs with respect to the ForGAN case, 
in fact, both training and test data remain the same.
The experiment has been conducted by using the following architecture and hyperparameters:
\begin{table}[H]
\centering
\caption{ForWGAN Architecture and Hyperparameters Summary}
\label{tab:forwgan_summary}
\begin{tabular}{lll}
\hline
\textbf{Parameter}         & \textbf{Generator}          & \textbf{Critic (Discriminator)} \\ \hline
Input Dimension            & 64 (Latent) + 252 (Cond.)   & 1 (Data) + 252 (Cond.)          \\
Hidden Layers              & [512, 256, 128, 64]         & [512, 256, 256, 128, 64]        \\
Output Dimension           & 1                           & 1                               \\
Activation Function        & ReLU                        & Leaky ReLU                      \\
Normalization              & None                        & Layer Norm                      \\
Dropout                    & 0.0                         & 0.0                             \\
\multicolumn{3}{c}{\textbf{Global Training Settings}}                                     \\ \hline
Max Epochs                 & \multicolumn{2}{l}{100}                                       \\
Learning Rate              & \multicolumn{2}{l}{0.0005}                                      \\
\end{tabular}
\end{table}


The Kolmogorov-Smirnov test performed on 100000 real and generated probability distributions shows that
the resulting p-values are accepted 34\% of times, quite higher than the ForGAN proposed models but still far from the desired level.\\
Moreover, the average Jensen-Shannon divergence over the test samples is around 0.28 which highlights an improvement with respect to the ForGAN model.\\
\begin{table}[h]
\centering
\begin{tabular}{l c}
\hline
\textbf{Metric} & \textbf{Value} \\
\hline
KS Test Acceptance (\%) & 34.00 \\
EMD distance & 0.00380 \\
JS Distance & 0.2802 \\
Hellinger Distance & 0.2510 \\
TV Distance & 0.2539 \\
\hline
\end{tabular}
\caption{ForWGAN performance metrics averaged over 100000 samples}
\end{table}


A visual inspection of different samples with various degrees of volatility is presented below
In particular, six different samples are shown below where
the starting values is the same $X_0 = 0.0$, $\mu = 0.0$ while the volatility have different values for each sample 
$\sigma = [0.3, 0.5, 0.7, 0.9]$.
From the results we can see an important improvement with respect to the simple ForGAN architecture,
suggesting that the Wasserstein loss with gradient penalty helps in stabilizing the training process 
and improving the quality of the generated distributions the same way has been demonstrated in papers like \cite{background:wgan} and \cite{background2:gp_wgan}.

\begin{figure}[H]
    \centering
    \includegraphics[width=1.0\textwidth]{Background/images/samples_forgan.png}
    \caption{Samples obtained using fixed $X_0$ and $\mu$ while varying $\sigma = [0.3, 0.5, 0.7, 0.9]$.}
    \label{forgan_sample_sigma}
\end{figure}




\section{SRForGAN experimental results}

The experimental setup is based on the same setup of the previous ForGAN section \ref{sec:bs experiment}.
The asset price trajectories are simulated over a fixed time horizon of one month, discretized into $N=22$ daily time steps.
For each experiment, $J=100{,}000$ independent paths are generated, ensuring sufficient sample diversity and statistical robustness.
The initial asset price $X_0$ is set for each trajectory to $0.0$, while the drift parameter $\mu$ is fixed at zero to isolate the model's ability to learn volatility-driven dynamics.
The volatility parameter $\sigma$ is sampled from a wide range $(0.1, 1.2)$ to capture varying market conditions, from near-deterministic to highly stochastic regimes.
The PDF to be learned refers to the $1$-th step ahead distribution, that is the distribution of asset prices at time $T + 1$ (or $N dt + 1$) given information up to time $T$.\\

Also in this case, the target distribution of the next value $c_{T + \tau}$  by using montecarlo simulations. 
The results are provided for the adversarial free SRForGAN architecture and also for the two 
hybrid approaches using the ForGAN and ForWGAN.\\



\begin{table}[H]
\centering
\caption{Generator Architecture and Hyperparameters Summary}
\label{tab:srforgan_architecture}
\begin{tabular}{lll}
\hline
\textbf{Parameter}         & \textbf{Generator}                  & \textbf{Critic (Discriminator)} \\ \hline
Input Dimension            & 32 (Latent) + 22 (Cond.)           & --                              \\
Hidden Layers              & [128, 256, 128, 64]                & --                              \\
Output Dimension           & 1                                  & --                              \\
Activation Function        & Leaky ReLU                         & --                              \\
Normalization              & BatchNormalization                 & --                              \\
Dropout                    & 0.0                                & --                              \\
\multicolumn{3}{c}{\textbf{Global Training Settings}}                                      \\ \hline
Batch Size                 & \multicolumn{2}{l}{32}                                            \\
max\_epochs                & \multicolumn{2}{l}{100}                                          \\
\hline
\end{tabular}
\end{table}


\begin{table}[h]
\centering
\begin{tabular}{l c}
\hline
\textbf{Metric} & \textbf{Value} \\
\hline
\label{tab:srforgan_results}
KS Test Acceptance (\%) & 100.00 \\
JS Distance & 0.1365 \\
Hellinger Distance & 0.1178 \\
TV Distance & 0.1099 \\
EMD Distance & 0.0018 \\
\hline
\end{tabular}
\caption{Evaluation metrics rounded to four decimal places}
\end{table}

\begin{figure}[H]
    \centering
    \includegraphics[width=1.0\textwidth]{SRForGAN/images/SRForGAN_samples.png}
    \caption{Samples obtained using fixed $X_0$, $\mu$ and $\sigma = [0.3, 0.5, 0.7, 0.9]$.}
    \label{srforgan_sample}
\end{figure}

It can be noticed that the SRForGAN architecture is able to learn the target distribution with a high degree of accuracy, as evidenced by the evaluation metrics and visual inspection of the generated samples. 
The KS test acceptance rate of 100\% indicates that the generated samples are statistically indistinguishable from the true distribution, while the low values of JS, Hellinger, TV, and EMD distances 
further confirm the model's effectiveness in capturing the underlying dynamics of the asset price trajectories. These results demonstrate
the potential of scoring rule-based generative models with respect to traditional adversarial approaches.\\
Because of these considerations further experiments have been conducted using more complicated distributions,
such as the Heston model, to assess the performance of the SRForGAN architecture in capturing complex stochastic processes.

\begin{comment}

%----------------------------------------------------------
%----------------------------------------------------------
By using a combined loss function made by both the energy score loss and the adversarial BCE or Wasserstein loss
the experiment has been repeated with the following architecture:


\begin{table}[H]
\centering
\caption{ForGAN Architecture and Hyperparameters Summary}
\label{tab:forgan_hybrid_architecture}
\begin{tabular}{lll}
\hline
\textbf{Parameter}         & \textbf{Generator}                  & \textbf{Critic (Discriminator)} \\ \hline
Input Dimension            & 32 (Latent) + 22 (Cond.)           & 1 (Data) + 22 (Cond.)           \\
Hidden Layers              & [128, 256, 128, 64]                & [128, 256, 256, 128, 64]        \\
Output Dimension           & 1                                  & 1                               \\
Activation Function        & Leaky ReLU                         & Leaky ReLU                      \\
Normalization              & BatchNormalization                 & None                            \\
Dropout                    & 0.0                                & 0.0                             \\
\multicolumn{3}{c}{\textbf{Global Training Settings}}                                      \\ \hline
Batch Size                 & \multicolumn{2}{l}{32}                                            \\
max\_epochs                & \multicolumn{2}{l}{100}                                          \\
\hline
\end{tabular}
\end{table}




\begin{table}[h]
\centering
\begin{tabular}{l c}
\hline
\textbf{Metric} & \textbf{Value} \\
\hline
\label{tab:hybrid_forgan_results}
KS Test Acceptance (\%) & 0.8530 \\
JS Distance & 0.2822 \\
Hellinger Distance & 0.2524 \\
TV Distance & 0.24029 \\
EMD Distance & 0.0102\\
\hline
\end{tabular}
\caption{Evaluation metrics rounded to four decimal places}
\end{table}
    


\begin{figure}[H]
    \centering
    \includegraphics[width=1.0\textwidth]{SRForGAN/images/Hybrid_ForGAN_samples.png}
    \caption{Samples obtained using fixed $X_0$, $\mu$ and $\sigma = [0.3, 0.5, 0.7, 0.9]$.}
    \label{hybrid_forgan_sample}
\end{figure}




\begin{table}[H]
\centering
\caption{ForWGAN Architecture and Hyperparameters Summary}
\label{tab:forwgan_hybrid_architecture}
\begin{tabular}{lll}
\hline
\textbf{Parameter}         & \textbf{Generator}                  & \textbf{Critic (Discriminator)} \\ \hline
Input Dimension            & 32 (Latent) + 22 (Cond.)           & 1 (Data) + 22 (Cond.)           \\
Hidden Layers              & [128, 256, 128, 64]                & [128, 256, 256, 128, 64]        \\
Output Dimension           & 1                                  & 1                               \\
Activation Function        & Leaky ReLU                         & Leaky ReLU                      \\
Normalization              & None                               & LayerNormalization              \\
Dropout                    & 0.0                                & 0.0                             \\
\multicolumn{3}{c}{\textbf{Global Training Settings}}                                      \\ \hline
Batch Size                 & \multicolumn{2}{l}{32}                                            \\
max\_epochs                & \multicolumn{2}{l}{100}                                          \\
\hline
\end{tabular}
\end{table}


\begin{table}[h]
\centering
\begin{tabular}{l c}
\hline
\textbf{Metric} & \textbf{Value} \\
\hline
\label{tab:hybrid_forwgan_results}
KS Test Acceptance (\%) & 0.9520 \\
JS Distance & 0.2653 \\
Hellinger Distance & 0.2348 \\
TV Distance & 0.2291 \\
EMD Distance & 0.0110\\
\hline
\end{tabular}
\caption{Evaluation metrics rounded to four decimal places}
\end{table}
   



\begin{figure}[H]
    \centering
    \includegraphics[width=1.0\textwidth]{SRForGAN/images/Hybrid_ForWGAN_samples.png}
    \caption{Samples obtained using fixed $X_0$, $\mu$ and $\sigma = [0.3, 0.5, 0.7, 0.9]$.}
    \label{hybrid_forwgan_sample}
\end{figure}

\end{comment}
\chapter{Experiments with Heston model dynamics}
\label{chap:heston}

This chapter provides experimental results  of the SRForgan applied to Heston model trajectories. 
In particular the \texttt{HestonSimulator}, a data-generation object that produces trajectories of financial log-prices under the Heston stochastic volatility model and computes the conditional probability density function of future log-prices using the Fractional Fast Fourier Transform (frFFT).
Every mathematical concept required to understand the implementation is introduced from first principles.

% ─────────────────────────────────────────────────────────────────────────────
\section{The Heston Stochastic Volatility Model}
\label{sec:heston_model}

The classical Black--Scholes model assumes that the instantaneous volatility of an asset price is a deterministic constant $\sigma$.
This assumption leads to a log-normal distribution for the terminal asset price, which is analytically tractable but empirically inadequate since financial markets exhibit a so-called
\textit{volatility smile}, whereby options at different strikes imply different values of $\sigma$ 
when inverted through the Black--Scholes formula. 
This observation alone establishes that volatility cannot be constant.

\subsection{The Heston System of SDEs}

To remedy this shortcoming, Heston \cite{heston:heston_paper} introduced in 1993 a two-factor model in which both the log-price $X_t = \ln S_t$ and the instantaneous variance $v_t$ are stochastic.
The system of stochastic differential equations (SDEs) in the \texttt{HestonSimulator} is written in \textit{log-price form}:

\begin{align}
    &dX_t = \left(\mu - \tfrac{1}{2}v_t\right)dt + \sqrt{v_t}\, dW_t^X \label{eq:logprice_sde_heston}\\
    &dv_t = \kappa(\theta - v_t)\,dt + \sigma_v\sqrt{v_t}\, dW_t^v \label{eq:var_sde_heston}\\
    &\langle dW_t^X,\, dW_t^v \rangle = \rho\, dt \label{eq:correlation_heston}
\end{align}

Each term carries a precise financial interpretation.
\begin{itemize}
    \item Log-price SDE (Eq. \ref{eq:logprice_sde_heston}). Applying It\^{o}'s lemma to $X_t = \ln S_t$ under the assumption that $S_t$ follows a geometric Brownian motion with stochastic diffusion coefficient $\sqrt{v_t}$ yields the drift $\mu - v_t/2$. The term $\mu$ is the risk-neutral drift (typically the risk-free rate), while the $-v_t/2$ correction, 
    called the \textit{It\^{o} correction}, arises from the quadratic variation of $X_t$ and ensures that $\mathbb{E}[S_t] = S_0 e^{\mu t}$.
    \item Variance SDE (Eq.\,\ref{eq:var_sde_heston}). The variance $v_t$ evolves according to a \textit{Cox--Ingersoll--Ross} (CIR) process \cite{heston:cir_process}.
     The three parameters have well-defined roles: $\kappa > 0$ is the \textit{mean-reversion speed}, 
     controlling how quickly $v_t$ is pulled back toward its long-run equilibrium; $\theta > 0$ is
      the \textit{long-run variance}, the level to which $v_t$ reverts; and $\sigma_v > 0$ is the \textit{volatility of volatility} (vol-of-vol), 
      governing the magnitude of random fluctuations in $v_t$.
      \item Correlation (Eq.\,\ref{eq:correlation_heston}). The two Brownian motions $W^X$ and $W^v$ are correlated with instantaneous correlation $\rho \in (-1, 1)$. For equities, $\rho$ is typically negative: a negative shock to the asset price tends to be accompanied by an increase in variance, a well-documented empirical phenomenon known as the \textit{leverage effect}.
\end{itemize}

A key requirement for the CIR process is that $v_t \geq 0$ at all times. it has been proved that the boundary $v_t = 0$ is never reached (i.e., $v_t > 0$ almost surely) if and only if

\begin{equation}
    2\kappa\theta > \sigma_v^2.
    \label{eq:feller}
\end{equation}

When condition the Feller condition (\ref{eq:feller}) is violated, the process touches zero with positive probability. In the simulator, this boundary is handled numerically by the \textit{full-truncation} scheme described in Section~\ref{sec:discretisation}.\\

% ─────────────────────────────────────────────────────────────────────────────

To simulate and compute the correlated pair $(dW_t^X, dW_t^v)$ over a discrete time step $\Delta t$, one exploits the fact that any bivariate Gaussian vector with correlation $\rho$ can be constructed from two independent standard normals $Z_1, Z_2 \sim \mathcal{N}(0,1)$ via the \textit{Cholesky decomposition}. Specifically, the increments

\begin{align}
    \Delta W^X &= \sqrt{\Delta t}\, Z_1 \label{eq:dWX}\\
    \Delta W^v &= \sqrt{\Delta t}\left(\rho\, Z_1 + \sqrt{1-\rho^2}\, Z_2\right) \label{eq:dWv}
\end{align}

satisfy $\mathbb{E}[\Delta W^X] = \mathbb{E}[\Delta W^v] = 0$, $\text{Var}(\Delta W^X) = \text{Var}(\Delta W^v) = \Delta t$, and $\text{Cov}(\Delta W^X, \Delta W^v) = \rho\,\Delta t$, which matches the covariance structure prescribed by Eq.~(\ref{eq:correlation_heston}). 
% ─────────────────────────────────────────────────────────────────────────────
\subsection{Numerical Discretisation of the Heston SDEs}
\label{sec:discretisation}

Since equations (\ref{eq:logprice_sde_heston})--(\ref{eq:var_sde_heston}) do not admit closed-form solutions in general, they must be integrated numerically. The proposed simulator supports two schemes:\\
The \textit{Euler--Maruyama} scheme \cite{heston:euler_maruyama} is the SDE analogue of Euler's method for ODEs. For a generic It\^{o} SDE $dY_t = a(Y_t)\,dt + b(Y_t)\,dW_t$, the scheme reads

\begin{equation}
    Y_{n+1} = Y_n + a(Y_n)\,\Delta t + b(Y_n)\,\Delta W_n.
    \label{eq:euler}
\end{equation}

Applied to the log-price SDE (\ref{eq:logprice_sde_heston}), with $\hat{v}_n = \max(v_n, 0)$ (the truncated variance):

\begin{equation}
    X_{n+1} = X_n + \left(\mu - \tfrac{1}{2}\hat{v}_n\right)\Delta t + \sqrt{\hat{v}_n}\,\Delta W_n^X.
    \label{eq:euler_logprice}
\end{equation}

The Euler--Maruyama scheme has strong convergence order $1/2$ and weak convergence order $1$ for SDEs with globally Lipschitz coefficients.
The second proposed discretization technique is the \textit{Milstein scheme} \cite{heston:milstein} improves upon Euler--Maruyama by including an additional correction term derived from It\^{o}'s formula applied to the diffusion coefficient. For an SDE with diffusion $b(Y) = \sigma_v\sqrt{Y}$, the correction is

\begin{equation}
    \frac{1}{2}b(Y_n)\frac{\partial b}{\partial Y}(Y_n)\bigl((\Delta W_n)^2 - \Delta t\bigr) = \frac{\sigma_v^2}{4}\bigl((\Delta W_n^v)^2 - \Delta t\bigr).
    \label{eq:milstein_correction}
\end{equation}

This term captures the curvature of the diffusion coefficient and raises the strong convergence order from $1/2$ to $1$.\\

Because the CIR process can reach zero when the Feller condition fails (or nearly so), a naive Milstein step may produce $v_{n+1} < 0$, which would make $\sqrt{v_{n+1}}$ undefined at the next step. A \textit{full-truncation} strategy addresses this in two stages:

\begin{enumerate}
    \item Before computing the drift and diffusion, replace $v_n$ with $\hat{v}_n = \max(v_n, 0)$ inside all coefficient functions.
    \item After the update step, apply a second truncation: $v_{n+1} \leftarrow \max(v_{n+1}, 0)$.
\end{enumerate}

The complete full-truncation Milstein step for the variance is therefore

\begin{equation}
    v_{n+1} = \max\!\Bigl(v_n + \kappa(\theta - \hat{v}_n)\Delta t + \sigma_v\sqrt{\hat{v}_n}\,\Delta W_n^v + \tfrac{1}{4}\sigma_v^2\bigl((\Delta W_n^v)^2 - \Delta t\bigr),\; 0\Bigr).
    \label{eq:milstein_full_trajectory}
\end{equation}


% ─────────────────────────────────────────────────────────────────────────────
\subsection{The Characteristic Function and the Fourier Inversion Formula}

The \textit{characteristic function} of a real-valued random variable $X$ is defined as

\begin{equation}
    \varphi(u) = \mathbb{E}\!\left[e^{iuX}\right], \quad u \in \mathbb{R},
    \label{eq:cf_def}
\end{equation}

where $i = \sqrt{-1}$. The Characteristic Function is the Fourier transform of the PDF and encodes all information about the distribution
 of $X$. Its fundamental importance stems from the fact that many stochastic models, including the Heston model, 
 admit a closed-form CF even when the PDF itself does not exist in closed form. If $\varphi$ is known analytically, the PDF can be recovered numerically via Fourier inversion.\\
For the Heston model, \cite{heston:heston_paper} it has been derived the conditional CF of the log-price $X_{t+\tau}$ given the state $(X_t, v_t)$ at time $t$.
The CF takes the affine-exponential form

\begin{equation}
    \varphi(u;\tau) = \exp\!\Bigl(iuX_0 + iu\mu\tau + A(u,\tau) + B(u,\tau)\,v_0\Bigr),
    \label{eq:cf_heston}
\end{equation}

where $X_0 = X_t$ is the current log-price, $v_0 = v_t$ is the current variance, and $\tau$ is the forecast horizon. The function $A$ captures the long-run contribution while $B$ captures the effect of the current variance. They both depend on the model parameters through an intermediate quantity $d$:

\begin{align}
    d &= \sqrt{(\kappa - \rho\sigma_v\, iu)^2 + \sigma_v^2(u^2 + iu)}, \quad \text{Re}(d) \geq 0, \label{eq:d_heston}\\
    \tilde{g} &= \frac{\kappa - \rho\sigma_v\, iu - d}{\kappa - \rho\sigma_v\, iu + d}, \label{eq:g_tilde}\\
    A(u,\tau) &= \frac{\kappa\theta}{\sigma_v^2}\left[(\kappa - \rho\sigma_v\, iu - d)\tau - 2\ln\!\frac{1 - \tilde{g}\,e^{-d\tau}}{1 - \tilde{g}}\right], \label{eq:A_heston}\\
    B(u,\tau) &= \frac{v_0}{\sigma_v^2}(\kappa - \rho\sigma_v\, iu - d)\frac{1 - e^{-d\tau}}{1 - \tilde{g}\,e^{-d\tau}}. \label{eq:B_heston}
\end{align}

The original Heston CF involves a complex-valued square root and a complex logarithm. Because the complex logarithm is multi-valued, numerical implementations must track the correct branch, which can lead to discontinuities in the CF as $u$ varies, a phenomenon known as the \textit{branch-cut problem}. This instability causes spurious oscillations in the numerically inverted PDF.

In \cite{heston:little_trap} Albrecher showed that a simple reformulation of Eq.~(\ref{eq:g_tilde}) eliminates the branch-cut issue entirely. The key insight is to define $\tilde{g}$ as the ratio $(\xi - d)/(\xi + d)$ where $\xi = \kappa - \rho\sigma_v\, iu$, and to enforce $\text{Re}(d) \geq 0$ in Eq.~(\ref{eq:d_heston}) by choosing the branch of the square root with non-negative real part. Under this convention, the logarithm in Eq.~(\ref{eq:A_heston}) remains on a consistent branch for all $u$, and the CF is smooth and numerically stable.\\




Given the CF $\varphi(u)$ of a continuous random variable $X$ with PDF $f(x)$, the \textit{Gil-Pelaez inversion theorem} \cite{heston:gil_pelaez} states that

\begin{equation}
    f(x) = \frac{1}{\pi}\int_0^\infty \text{Re}\!\left[\varphi(u)\,e^{-iux}\right]du.
    \label{eq:gilpelaez}
\end{equation}

This follows from the standard Fourier inversion theorem and the fact that $\varphi(-u) = \overline{\varphi(u)}$ for real-valued distributions, which halves the integration range. To evaluate Eq.~(\ref{eq:gilpelaez}) numerically, the integral is truncated at a finite upper limit $U = (N_\text{fft}-1)\eta$ and discretised on a uniform grid $u_j = j\eta$, $j = 0, 1, \ldots, N_\text{fft}-1$, with frequency step $\eta > 0$:

\begin{equation}
    f(x) \approx \frac{\eta}{\pi}\,\text{Re}\!\left[\sum_{j=0}^{N_\text{fft}-1} w_j\,\varphi(j\eta)\,e^{-ij\eta x}\right],
    \label{eq:gilpelaez_discrete}
\end{equation}

where $w_j$ are numerical quadrature weights.
The choice of quadrature weights $w_j$ in Eq.~(\ref{eq:gilpelaez_discrete}) determines the accuracy of the numerical integral. Simple rectangular (Euler) weights give first-order accuracy. The implementation instead uses \textit{Simpson's rule}, which achieves fourth-order accuracy by approximating the integrand as a piecewise quadratic polynomial. For $N_\text{fft}$ equally spaced points with step $\eta$, Simpson's rule assigns weights

\begin{equation}
    w_j = \frac{1}{3}\begin{cases} 1 & j = 0 \text{ or } j = N_\text{fft}-1, \\ 4 & j \text{ odd}, \\ 2 & j \text{ even}, \; j \neq 0, N_\text{fft}-1, \end{cases}
    \label{eq:simpson}
\end{equation}

The \textit{Discrete Fourier Transform} (DFT) of a sequence $\{x_j\}_{j=0}^{N-1}$ is defined as

\begin{equation}
    \hat{x}_k = \sum_{j=0}^{N-1} x_j\, e^{-i2\pi jk/N}, \quad k = 0, 1, \ldots, N-1.
    \label{eq:dft}
\end{equation}

The \textit{Fast Fourier Transform} (FFT) computes Eq.~(\ref{eq:dft}) in $\mathcal{O}(N\log N)$ operations rather than the $\mathcal{O}(N^2)$ required by direct summation. To apply the FFT to the inversion formula (\ref{eq:gilpelaez_discrete}), one evaluates Eq.~(\ref{eq:gilpelaez_discrete}) simultaneously at the grid of log-price values $x_k = b + k\lambda$, $k = 0, \ldots, N_\text{fft}-1$, for some lower bound $b$ and spacing $\lambda$. Inserting this into Eq.~(\ref{eq:gilpelaez_discrete}):

\begin{align}
    f(x_k) \approx \frac{\eta}{\pi}\,\mathrm{Re} \left[\sum_{j=0}^{N_{\mathrm{FFT}}-1}w_j\,\varphi(j\eta)\,e^{-ij\eta b}\,e^{-i2\pi\frac{\eta\lambda}{2\pi}jk}\right]
    \label{eq:fft_coupl}
\end{align}

Comparing with Eq.~(\ref{eq:dft}), the sum is a DFT if and only if the exponent is $e^{-i2\pi jk/N_\text{fft}}$, which requires

\begin{equation}
    \frac{\eta\lambda}{2\pi} = \frac{1}{N_\text{fft}} \;\Longrightarrow\; \lambda = \frac{2\pi}{N_\text{fft}\,\eta}.
    \label{eq:std_fft_constraint}
\end{equation}

This \textit{coupling constraint} links the frequency step $\eta$ and the spatial step $\lambda$: making the log-price grid finer (smaller $\lambda$) forces a larger $\eta$, which reduces the effective integration range in frequency and may introduce aliasing or truncation errors. This rigid coupling is the central limitation of the standard FFT approach to Fourier inversion.\\

The \textit{Fractional DFT} (frDFT) generalises Eq.~(\ref{eq:dft}) by replacing $1/N$ with an arbitrary real parameter $\alpha$:

\begin{equation}
    y_k = \sum_{j=0}^{N-1} x_j\, e^{-i2\pi\alpha jk}, \quad k = 0, \ldots, N-1.
    \label{eq:frdft}
\end{equation}

Comparing with Eq.~(\ref{eq:fft_coupl}), choosing $\alpha = \eta\lambda/(2\pi)$ evaluates the inversion integral at the spatial grid $\{x_k\}$ without any constraint between $\eta$ and $\lambda$. The user is therefore free to choose $\eta$ and $\lambda$ independently, completely decoupling frequency resolution from spatial resolution.
However, computing Eq.~(\ref{eq:frdft}) naively requires $\mathcal{O}(N^2)$ operations. 
Bluestein observed that the product $jk$ can be written as

\begin{equation}
    jk = \frac{j^2 + k^2 - (j-k)^2}{2},
    \label{eq:bluestein_id}
\end{equation}

which transforms Eq.~(\ref{eq:frdft}) into

\begin{equation}
    y_k = \underbrace{e^{-i\pi\alpha k^2}}_{c_k^*}\sum_j \underbrace{\bigl[x_j\, e^{-i\pi\alpha j^2}\bigr]}_{a_j = x_j c_j^*} \cdot \underbrace{e^{i\pi\alpha(j-k)^2}}_{h_{j-k}}= c_k^* \sum_j a_j\, h_{j-k}
    \label{eq:bluestein_conv}
\end{equation}

The sum over $j$ in Eq.~(\ref{eq:bluestein_conv}) is a \textit{linear convolution} of the sequence $a_j = x_j e^{-i\pi\alpha j^2}$ with the symmetric \textit{chirp kernel} $h_m = e^{i\pi\alpha m^2}$, evaluated at lags $m = j - k$. Linear convolution of length-$N$ sequences can be computed in $\mathcal{O}(N\log N)$ via the convolution theorem:
Note the FFT of $h$ is computed once so every $J$ trajectories
share one chirp and one kernel FFT, only the data-dependent part of the FFT 
operation differs per row.

\begin{equation}
    \text{Linear convolution: }\; a \star h = \mathcal{F}^{-1}\!\left[\mathcal{F}[a]\cdot\mathcal{F}[h]\right],
    \label{eq:conv_theorem}
\end{equation}

where the transforms are taken on a zero-padded array of length $M \geq 2N - 1$ (rounded up to the next power of two to allow the FFT) to prevent circular aliasing. The full frFFT algorithm is:

\begin{enumerate}
    \item Compute $a_j = x_j e^{-i\pi\alpha j^2}$ (phase-modulate input by conjugate chirp).
    \item Zero-pad $a$ to length $M = 2^{\lceil\log_2(2N-1)\rceil}$.
    \item Build the chirp kernel $h_{\text{pad}}$ of length $M$ in circular order: $h[0], h[1], \ldots, h[N-1]$ followed by zeros, then $h[N-1], \ldots, h[1]$.
    \item Compute $Y = \texttt{IFFT}(\texttt{FFT}(a_{\text{pad}}) \cdot \texttt{FFT}(h_{\text{pad}}))$.
    \item Multiply the first $N$ outputs: $y_k = e^{-i\pi\alpha k^2} Y_k$, $k = 0, \ldots, N-1$.
\end{enumerate}

\begin{algorithm}[H]
    \label{alg:frFTT}
\caption{Batched Fractional FFT via Bluestein's Convolution Identity}
\begin{algorithmic}[1]
\REQUIRE Batch $\mathbf{X}\in\mathbb{C}^{J\times N}$ (row $j$ = trajectory $j$'s frequency samples),
         fractional parameter $\alpha=\eta\lambda/(2\pi)$
\ENSURE  $\mathbf{Y}\in\mathbb{C}^{J\times N}$ with
         $Y_{j,k}=\sum_{n=0}^{N-1}X_{j,n}\,e^{-i2\pi\alpha nk}$,\; $k=0,\ldots,N-1,\; j=1,\ldots,J$

\LINECOMMENT{Identity: $nk=\tfrac12[n^2+k^2-(n-k)^2]$, so
  $e^{-i2\pi\alpha nk}=c_n^{*}c_k^{*}\,h_{n-k}$ with chirp $c_n=e^{i\pi\alpha n^2}$, $h_m=c_m$}

\LINECOMMENT{Step 0 --- precompute the chirp once, shared by every trajectory}
\STATE $c_n \leftarrow e^{i\pi\alpha n^2}$,\quad $n=0,\ldots,N-1$

\LINECOMMENT{Step 1 --- demodulate each row of the input by the conjugate chirp}
\FOR{$j=1$ \textbf{to} $J$}
    \STATE $a_{j,n}\leftarrow X_{j,n}\,c_n^{*}$,\quad $n=0,\ldots,N-1$
\ENDFOR

\LINECOMMENT{Step 2 --- padded length avoiding circular wrap-around}
\STATE $M\leftarrow 2^{\lceil\log_2(2N-1)\rceil}$

\LINECOMMENT{Step 3 --- zero-pad every row of $\mathbf{a}$ to length $M$}
\STATE $\mathbf{a}_{\mathrm{pad}}\leftarrow\mathbf{0}_{J\times M}$;\quad
       $\mathbf{a}_{\mathrm{pad}}[:,0{:}N]\leftarrow\mathbf{a}$

\LINECOMMENT{Step 4 --- build the kernel from the same chirp, laid out circularly
  ($h_{-m}=h_m$ places negative lags at the tail of length $M$)}
\STATE $\mathbf{h}_{\mathrm{pad}}\leftarrow\mathbf{0}_M$
\STATE $\mathbf{h}_{\mathrm{pad}}[0{:}N]\leftarrow[c_0,\ldots,c_{N-1}]$
       \COMMENT{lags $m=0,\ldots,N-1$}
\STATE $\mathbf{h}_{\mathrm{pad}}[M-N+1{:}M]\leftarrow[c_{N-1},\ldots,c_1]$
       \COMMENT{lags $m=-(N-1),\ldots,-1$}

\LINECOMMENT{Step 5 --- linear convolution via the convolution theorem;
  $\mathrm{FFT}(\mathbf{h}_{\mathrm{pad}})$ computed once and shared across all $J$ rows}
\STATE $\mathbf{conv}\leftarrow \mathrm{IFFT}\bigl(\mathrm{FFT}(\mathbf{a}_{\mathrm{pad}},\mathrm{axis}=1)\odot\mathrm{FFT}(\mathbf{h}_{\mathrm{pad}})\bigr)$

\LINECOMMENT{Step 6 --- demodulate by the chirp again, keep first $N$ entries}
\STATE $Y_{j,k}\leftarrow c_k^{*}\,\mathrm{conv}_{j,k}$,\quad $k=0,\ldots,N-1,\;j=1,\ldots,J$

\RETURN $\mathbf{Y}\in\mathbb{C}^{J\times N}$
\end{algorithmic}
\end{algorithm}


In the current implementation, the fractional parameter is set by default to $\alpha = 0.5/N_\text{fft}$, giving the spatial spacing

\begin{equation}
    \lambda = \frac{2\pi\alpha}{\eta} = \frac{4\pi}{N_\text{fft}\,\eta}.
    \label{eq:lambda_frfft}
\end{equation}

For $N_\text{fft} = 4096$ and $\eta = 0.25$, this yields $\lambda \approx 0.0031$, which is two orders of magnitude finer than the $\lambda_{\text{std}} = 2\pi/(N_\text{fft}\,\eta) \approx 0.0061$ that the standard FFT would enforce. The user can further tighten the grid by choosing a smaller $\alpha$ without any effect on the frequency resolution $\eta$.

% ─────────────────────────────────────────────────────────────────────────────
\subsection{The PDF Computation Pipeline}

In order to obtain the probability distribution function the code assembles the above components into a complete PDF computation. The following describes each stage of the pipeline.


\begin{itemize}
    \item \textbf{Grid Initialisation}. For each of the $J$ trajectories, a lower bound $b_j$ for the log-price evaluation grid is determined from an approximate mean and standard deviation of the forecast distribution. Since the exact moments depend on the full path of $v_t$, an approximation using the long-run variance $\theta$ is used:
    \begin{align}
        \text{approx\_mean}_j &= X_{0,j} + \left(\mu_j - \tfrac{1}{2}\theta_j\right)\tau, \label{eq:approx_mean}\\
        \text{approx\_std}_j  &= \sqrt{\max(\theta_j, 0)\,\tau}, \label{eq:approx_std}\\
        b_j &= \text{approx\_mean}_j - c\cdot\text{approx\_std}_j, \label{eq:lower_bound}
    \end{align}
    where $c = 5$ (the parameter \texttt{x\_width}) ensures that the grid spans approximately $\pm 5$ standard deviations around the forecast mean, capturing virtually all of the probability mass.
    \item \textbf{Phase Shift for the lower bound}. Equation (\ref{eq:gilpelaez_discrete}) evaluated at $x_k = b + k\lambda$ factors as
    \begin{equation}
    f(b + k\lambda) \approx \frac{\eta}{\pi}\,\text{Re}\!\left[\sum_{j=0}^{N_\text{fft}-1} \underbrace{\bigl(w_j\,\varphi(j\eta)\,e^{-ij\eta b}\bigr)}_{\text{frFFT input}} \cdot e^{-i2\pi\alpha jk}\right].
    \label{eq:phase_shifted}
    \end{equation}
    The factor $e^{-ij\eta b}$ shifts the output grid's origin to $b$ without changing the frequency-domain data. 
    \item \textbf{PDF extraction and clipping}. After the frFFT, the approximate PDF values are
    \begin{equation}
    f(x_k) \approx \frac{\eta}{\pi}\,\text{Re}\!\left[y_k\right], \quad k = 0, \ldots, N_\text{fft}-1.
    \label{eq:pdf_extract}
    \end{equation}
    Numerical inversion can produce small negative values due to discretisation error and finite truncation of the frequency integral. These artefacts are removed by clipping: \texttt{pdf\_fine = np.maximum(pdf\_fine, 0.0)}, which has negligible effect on well-resolved distributions.
    \item \textbf{Raw moments output}. When \texttt{n\_bins=None}, the method instead computes the mean and standard deviation directly from the fine frFFT grid using the rectangle rule:
    \begin{align}
    \mathbb{E}[X] &\approx \frac{\sum_k f(x_k)\,x_k\,\lambda}{\sum_k f(x_k)\,\lambda}, \label{eq:raw_mean}\\
    \text{Var}(X) &\approx \frac{\sum_k f(x_k)\,(x_k - \mathbb{E}[X])^2\,\lambda}{\sum_k f(x_k)\,\lambda}. \label{eq:raw_var}
    \end{align}
    The output is then a $(J, 2)$ array with columns containing the mean and standard deviation for each trajectory.
\end{itemize}




\begin{algorithm}[H]
    \label{alg:heston_pdf}
\caption{Heston PDF via Fractional FFT and Gil-Pelaez Inversion}
\begin{algorithmic}[1]
\REQUIRE Simulated terminal states $\{X_{T}^{(j)}, v_{T}^{(j)}\}_{j=1}^{J}$,
         Heston parameters $\{\mu^{(j)}, \kappa^{(j)}, \theta^{(j)},
         \sigma_v^{(j)}, \rho^{(j)}\}_{j=1}^{J}$
\REQUIRE Steps ahead $n_{\tau}$, time step $\Delta t$,
         FFT length $N$, frequency step $\eta$,
         fractional parameter $\alpha$, grid half-width $c$,
         number of output bins $K$ (or \texttt{None})
\ENSURE  Discretised conditional PDF or raw moments for each trajectory

\LINECOMMENT{Forecast horizon and x-domain spacing}
\STATE $\tau \leftarrow n_{\tau}\,\Delta t$
\STATE $\lambda \leftarrow 2\pi\alpha / \eta$

\LINECOMMENT{Per-trajectory lower bounds for the x-grid}
\FOR{$j = 1$ \textbf{to} $J$}
    \STATE $\tilde{\mu}^{(j)}   \leftarrow X_{T}^{(j)}
           + \bigl(\mu^{(j)} - \tfrac{1}{2}\theta^{(j)}\bigr)\tau$
           \COMMENT{approximate centre}
    \STATE $\tilde{\sigma}^{(j)} \leftarrow \sqrt{\theta^{(j)}\,\tau}$
           \COMMENT{approximate spread}
    \STATE $b^{(j)} \leftarrow \tilde{\mu}^{(j)} - c\,\tilde{\sigma}^{(j)}$
           \COMMENT{grid lower bound}
\ENDFOR

\LINECOMMENT{Frequency grid and Simpson quadrature weights}
\STATE $u_n \leftarrow n\,\eta$,\quad $n = 0, \ldots, N-1$
\STATE $w_0 \leftarrow 1/3$;\quad
       $w_n \leftarrow 4/3$ for $n$ odd;\quad
       $w_n \leftarrow 2/3$ for $n$ even, $n \notin \{0, N-1\}$;\quad
       $w_{N-1} \leftarrow 1/3$

\LINECOMMENT{Evaluate Heston characteristic function (eq. \ref{eq:cf_heston}) for all trajectories}
\STATE $\boldsymbol{\varphi} \leftarrow
       \textsc{HestonCF}\!\left(
         \mathbf{u},\,\tau,\,
         \{X_{T}^{(j)}\},\{v_{T}^{(j)}\},
         \{\mu^{(j)}\},\{\kappa^{(j)}\},\{\theta^{(j)}\},\{\sigma_v^{(j)}\},\{\rho^{(j)}\}
       \right)
       \in \mathbb{C}^{J \times N}$

\LINECOMMENT{Phase-shift to lower bound $b^{(j)}$, then weight by quadrature}
\STATE $\mathbf{Z}^{(j)}_n \leftarrow
       w_n\,\varphi^{(j)}_n\,
       e^{-i\,u_n\,b^{(j)}}$,\quad
       $j = 1,\ldots,J,\quad n = 0,\ldots,N-1$

\LINECOMMENT{Batched frFFT inversion (Algorithm~\ref{alg:frFTT})}
\STATE $\mathbf{Y} \leftarrow
       \textsc{FrFFT\_Batch}\!\left(\mathbf{Z},\,\alpha\right)
       \in \mathbb{C}^{J \times N}$

\LINECOMMENT{Recover density values and clip numerical artefacts}
\STATE $f^{(j)}_k \leftarrow
       \max\!\left(0,\;
         \dfrac{\eta}{\pi}\,\mathrm{Re}\!\left[Y^{(j)}_k\right]
       \right)$,\quad
       $k = 0,\ldots,N-1$

\LINECOMMENT{x-grid for each trajectory}
\STATE $x^{(j)}_k \leftarrow b^{(j)} + k\,\lambda$,\quad
       $k = 0,\ldots,N-1$

\IF{$K$ is \texttt{None}}
    \LINECOMMENT{Return raw moments via numerical integration}
    \FOR{$j = 1$ \textbf{to} $J$}
        \STATE $C^{(j)} \leftarrow \lambda\sum_{k} f^{(j)}_k$
               \COMMENT{normalisation constant}
        \STATE $\bar{x}^{(j)} \leftarrow
               \lambda\sum_{k} x^{(j)}_k\,f^{(j)}_k \;/\; C^{(j)}$
        \STATE $s^{(j)} \leftarrow
               \sqrt{\lambda\sum_{k}
                 \bigl(x^{(j)}_k - \bar{x}^{(j)}\bigr)^{2}
                 f^{(j)}_k \;/\; C^{(j)}}$
    \ENDFOR
    \RETURN $\bigl[\bar{x}^{(j)},\,s^{(j)}\bigr]_{j=1}^{J}$
\ENDIF

\RETURN $\mathbf{p} = f^{(j)} * x^{(j)} \in \mathbb{R}^{J \times K}$

\end{algorithmic}
\end{algorithm}

Moreover, the code provides a reference implementation based on direct Monte Carlo simulation. For each of the $J$ starting conditions $(X_{0,j}, v_{0,j})$, a large number \texttt{mc\_sims} of independent forward paths are generated using the full-truncation Euler--Maruyama scheme over \texttt{n\_steps} steps. The terminal log-prices are then histogrammed onto the shared bin structure to produce an empirical approximation of the PDF.

The Monte Carlo estimator is consistent (converges to the true PDF as \texttt{mc\_sims} $\to\infty$) but slow: its computational cost is $\mathcal{O}(J \cdot \texttt{mc\_sims} \cdot \texttt{n\_steps})$. The frFFT method, by contrast, evaluates the exact Heston CF and inverts it in $\mathcal{O}(J \cdot N_\text{fft}\log N_\text{fft})$ operations, making it orders of magnitude faster for large $J$.

\begin{algorithm}[H]
\caption{Monte-Carlo Conditional Heston PDF}
\label{alg:heston_mc_pdf}
\begin{algorithmic}[1]

\REQUIRE Current states $\{X_0^{(j)},v_0^{(j)}\}_{j=1}^{J}$,
Heston parameters
$\{\kappa^{(j)},\theta^{(j)},\sigma_v^{(j)},\rho^{(j)},\mu^{(j)}\}_{j=1}^{J}$,
forecast horizon $\tau=n_\tau\Delta t$,
Monte-Carlo sample size $M$

\ENSURE Terminal log-price samples
$\mathbf{X}\in\mathbb{R}^{J\times M}$

\FOR{$j=1$ \textbf{to} $J$}

    \STATE Initialise a CIR process from $v_0^{(j)}$ using
    $(\kappa^{(j)},\theta^{(j)},\sigma_v^{(j)})$.

    \STATE Evolve the variance process
    (\ref{eq:var_sde_heston}) over $[0,\tau]$
    using Andersen's Quadratic--Exponential (QE) scheme,
    simultaneously computing
    \[
        v_\tau,
        \qquad
        \int_0^\tau v(s)\,ds
    \]
    for all $M$ Monte-Carlo paths.

\STATE   Compute the conditional log-price increment
\[
Y=
-\frac12 I_v
+\frac{\rho}{\sigma_v}
\Bigl(v_\tau-v_0-\kappa(\theta\tau-I_v)\Bigr)
+\sqrt{(1-\rho^2)I_v}\,Z,
\]
where $I_v=\int_0^\tau v(s)\,ds$.

\STATE Set
\[
X_\tau=X_0+Y+\mu\tau.
\]

    \STATE Store the $M$ simulated terminal values as row $j$ of
    $\mathbf{X}$.

\ENDFOR

\RETURN $\mathbf{X}$

\end{algorithmic}
\end{algorithm}

So the \texttt{HestonSimulator} integrates the following chain of mathematical techniques into a single coherent pipeline. 
The Heston SDE system (\ref{eq:logprice_sde_heston})--(\ref{eq:correlation_heston}) is discretised using correlated Brownian increments constructed via Cholesky decomposition, 
with the full-truncation Milstein scheme for the variance and Euler--Maruyama for the log-price. 
The conditional PDF of future log-prices is obtained by evaluating the closed-form Heston CF in its numerically 
stable formulation and then applying the Gil-Pelaez inversion formula (\ref{eq:gilpelaez_discrete}) with Simpson quadrature weights.
The resulting discrete sum is evaluated efficiently across all $N_\text{fft}$ spatial grid points simultaneously by expressing it as a fractional DFT and computing
it via Bluestein's chirp convolution algorithm, which reduces the cost from $\mathcal{O}(N_\text{fft}^2)$ to $\mathcal{O}(N_\text{fft}\log N_\text{fft})$ 
while completely decoupling the frequency resolution $\eta$ from the spatial resolution $\lambda$.\\
The Monte Carlo estimator is validated against the frFFT method over $10000$ sample trajectories obtaining a $100\%$ of acceptance rate of the Kolmogorov-Smirnov test at a $5\%$ significance level.
Below are presented four samples for visual inspection of the two methods, showing a very good agreement between them.

\begin{figure}[H]
    \centering
    \includegraphics[width=\textwidth]{Heston/images/mc_heston_pdf.png}
    \caption{Comparison of the Heston PDF obtained via the frFFT method (red line) and the Monte Carlo estimator (blue histogram) for four sample trajectories. }
    \label{fig:heston_mc_plots}
\end{figure}



% ─────────────────────────────────────────────────────────────────────────────
% Section: Numerical Validation of the Fractional FFT via the
%          Black-Scholes Characteristic Function
% ─────────────────────────────────────────────────────────────────────────────

\section{Numerical Validation of the Fractional FFT via the
         Black-Scholes Characteristic Function}
\label{sec:frfft_validation}

The implementation of the fractional Fast Fourier Transform presented in the
preceding section constitutes a complex numerical pipeline: it involves a
quadrature approximation of an improper integral, a chirp-modulated linear
convolution evaluated through nested FFT calls, and a subsequent restriction
of the recovered density to a bounded support.  Each step introduces
approximation error, and it is therefore essential to verify that the full
pipeline produces numerically correct results before it is deployed inside the
Heston simulator.

The validation strategy adopted here exploits a key property of the
Black-Scholes model: under Geometric Brownian Motion the conditional
log-price distribution has a \emph{known closed form}.  Specifically, given
the log-price $X_t$ at time $t$, the one-step-ahead log-price
$X_{t+\tau} \mid X_t$ follows a Gaussian distribution with analytically
available parameters.  By feeding the Black-Scholes characteristic function
into the fractional FFT inversion routine and comparing the recovered density
to that closed-form Gaussian, one obtains a ground-truth benchmark that is
model-free with respect to the numerical algorithm under test.  The
experiment is therefore not a Monte Carlo approximation: it is an exact
algebraic comparison, and any discrepancy between the two densities is
attributable entirely to numerical error in the frFFT machinery.\\


Under the Black-Scholes model the log-price evolves according to the
stochastic differential equation
\begin{equation}
    dX_t \;=\; \Bigl(\mu - \tfrac{1}{2}\sigma^2\Bigr)\,dt \;+\; \sigma\,dW_t,
    \label{eq:bs_sde}
\end{equation}
where $\mu$ is the drift, $\sigma > 0$ is the volatility, and $W_t$ is a
standard Brownian motion.  Over a fixed horizon $\tau > 0$ the increment
$X_{t+\tau} - X_t$ is normally distributed with mean
$\bigl(\mu - \tfrac{1}{2}\sigma^2\bigr)\tau$ and variance $\sigma^2\tau$,
so the conditional distribution reads
\begin{equation}
    X_{t+\tau} \mid X_t \;\sim\;
    \mathcal{N}\!\bigl(\mu_{\mathrm{BS}},\,\sigma_{\mathrm{BS}}^2\bigr),
    \qquad
    \mu_{\mathrm{BS}} = X_t + \bigl(\mu - \tfrac{1}{2}\sigma^2\bigr)\tau,
    \quad
    \sigma_{\mathrm{BS}}^2 = \sigma^2\tau.
    \label{eq:bs_normal}
\end{equation}
The corresponding characteristic function (CF) is
\begin{equation}
    \varphi_{\mathrm{BS}}(u)
    \;=\;
    \mathbb{E}\bigl[e^{iuX_{t+\tau}}\bigr]
    \;=\;
    \exp\!\Bigl(iu\,\mu_{\mathrm{BS}} \;-\; \tfrac{1}{2}\sigma_{\mathrm{BS}}^2\,u^2\Bigr),
    \qquad u \in \mathbb{R}.
    \label{eq:bs_cf}
\end{equation}
This is the CF of a Gaussian and is entire in $u$, which makes the
Gil-Pelaez inversion numerically well-conditioned: the integrand decays
super-exponentially as $|u| \to \infty$, so the truncated quadrature sum
converges rapidly.

\subsection{Validation Design}
\label{subsec:validation_design}

The experiment is structured in three sequential validation steps, each
targeting a different layer of the computational pipeline.\\


A set of $J = 10000$ independent trajectories is simulated using the 
same algorithm described in section \ref{sec:bs_experiments}, in particular \ref{alg:y condition}, with $N = 252$ time steps over a horizon of
$T = 1$ year, giving a time-step size of $dt = T/N = 0.004$ years.  The
per-trajectory drift and volatility parameters are drawn uniformly:
$\mu \sim \mathcal{U}(0,\, 0.1)$ and $\sigma \sim \mathcal{U}(0.1,\, 1.2)$,
while the initial log-price satisfies $X_0 \sim \mathcal{U}(0,\, 1)$.  With
a one-step-ahead forecast horizon $\tau = dt = 0.004$ years.\\

The inversion is performed with $N_{\mathrm{FFT}} = 4096$ frequency nodes
and a frequency-domain step $\eta = 0.25$.  Following algorithm \ref{alg:frFTT}. The fractional parameter is set to
$\alpha = 2/N_{\mathrm{FFT}} \approx 4.88 \times 10^{-4}$, which yields an
$x$-domain spacing of
\begin{equation}
    \lambda \;=\; \frac{2\pi\alpha}{\eta} \;\approx\; 0.01227.
    \label{eq:lambda}
\end{equation}
The lower bound of the $x$-grid for each trajectory $k$ is placed at
$b_k = \mu_{\mathrm{BS},k} - 5\,\sigma_{\mathrm{BS},k}$, centring the
grid on the Gaussian peak and ensuring that $\pm 5\sigma$ tails are always
covered.\\

\paragraph{Step A: Analytical moments sanity check.}
Before engaging the frFFT at all, the first validation step confirms that the
method used to obtain the analytical probability distribution function correctly implements
equation~\eqref{eq:bs_normal}.  The
method returns a $(J \times 2)$ array whose columns are the analytical mean
$\mu_{\mathrm{BS},k}$ and standard deviation $\sigma_{\mathrm{BS},k}$ for
each trajectory $k$.  These are compared against the values computed directly
from the simulation parameters and the closed-form expressions in
equation~\eqref{eq:bs_normal}.

The results, summarised in Table~\ref{tab:step_a}, show that the analytical
moments returned by the code implementation are \emph{bit-for-bit identical} to the
directly computed values: the absolute difference in both mean and standard
deviation is exactly zero across all four trajectories.  This is expected,
since both quantities are computed from the same floating-point operations on
the same arrays; nevertheless it confirms that there is no implementation
error in the formula for the analytical parameters.

\begin{table}[H]
\centering
\caption{Step A: comparison between directly computed analytical parameters
and those returned by code simulation for the first 4 trajectories.
All differences are zero to double-precision.}
\label{tab:step_a}
\begin{tabular}{ccccccc}
\toprule
$k$ & $\mu_{\mathrm{BS},k}$ & \texttt{computed} mean & $|\Delta\mu|$
    & $\sigma_{\mathrm{BS},k}$ & \texttt{computed} std & $|\Delta\sigma|$ \\
\midrule
0 & 0.8533110 & 0.8533110 & $0$ & 0.0232011 & 0.0232011 & $0$ \\
1 & 0.1042671 & 0.1042671 & $0$ & 0.0459892 & 0.0459892 & $0$ \\
2 & 0.6879010 & 0.6879010 & $0$ & 0.0403615 & 0.0403615 & $0$ \\
3 & 0.7827374 & 0.7827374 & $0$ & 0.0796743 & 0.0796743 & $0$ \\
\bottomrule
\end{tabular}
\end{table}

\paragraph{Step B: frFFT Inversion.}
The core validation step applies the Gil-Pelaez inversion formula to the
Black-Scholes CF of equation~\eqref{eq:bs_cf}.  The inversion approximates
the probability density function as
\begin{equation}
    f(x_k) \;\approx\;
    \frac{\eta}{\pi}\,
    \mathrm{Re}\!\left[
        \sum_{j=0}^{N_{\mathrm{FFT}}-1}
        w_j\,\varphi_{\mathrm{BS}}(j\eta)\,e^{-ij\eta x_k}
    \right],
    \qquad x_k = b + k\lambda,
    \label{eq:gp_discrete}
\end{equation}
where $w_j$ are Simpson quadrature weights.  The sum over $j$ is evaluated
simultaneously at all $k = 0, \ldots, N_{\mathrm{FFT}} - 1$ grid points
using the batched fractional FFT described in algorithm \ref{alg:frFTT},
which implements Bluestein's identity.  Concretely, defining the frequency-domain
input
\begin{equation}
    \tilde{x}_j \;=\; w_j\,\varphi_{\mathrm{BS}}(j\eta)\,e^{-ij\eta b},
    \label{eq:fft_input}
\end{equation}
equation~\eqref{eq:gp_discrete} reduces to the fractional DFT
\begin{equation}
    Y[k] \;=\;
    \sum_{j=0}^{N-1} \tilde{x}_j\,e^{-i2\pi\alpha jk},
    \qquad
    f(x_k) = \frac{\eta}{\pi}\,\mathrm{Re}[Y[k]],
    \label{eq:frfft_sum}
\end{equation}

An important numerical subtlety arises from the choice $\alpha = 2/N$.
In this regime, the frFFT evaluates the standard DFT at even-indexed
frequencies: $Y[k] = \hat{X}[2k]$, where $\hat{X}$ denotes the length-$N$
DFT of the input sequence.  A consequence of this sub-sampling in frequency
space is that the second half of the output grid aliases onto the first half,
effectively doubling the contribution of the Gaussian peak.  As a result, the
raw integral
\begin{equation}
    \int_{-\infty}^{\infty} f(x)\,dx
    \;\approx\;
    \lambda \sum_{k=0}^{N-1} f(x_k)
    \;\approx\; 2,
    \label{eq:raw_integral}
\end{equation}
instead of the theoretically correct value of one.  The numerical output
confirms this precisely: the raw integral equals $2.0000$ for every
trajectory, as shown below.

\begin{center}
\texttt{t0:2.0000 \quad t1:2.0000 \quad t2:2.0000 \quad t3:2.0000}
\end{center}

This factor of two is \emph{systematic and entirely predictable}: it does
not represent an error in the inversion formula, but rather a known property
of the fractional parameter choice $\alpha = 2/N$.  The \texttt{HestonSimulator}
handles it transparently by normalising the recovered density after
restricting it to the relevant support, exactly as one would normalise a
histogram.  Specifically, the density is zeroed outside the interval
$[\mu_{\mathrm{BS}} - 5\sigma_{\mathrm{BS}},\,
  \mu_{\mathrm{BS}} + 5\sigma_{\mathrm{BS}}]$
and then rescaled so that $\lambda \sum_k f(x_k) = 1$.  This restriction also
has a direct physical interpretation: a $\pm 5\sigma$ interval captures
$99.9999\%$ of the Gaussian probability mass, so no meaningful probability is
discarded.\\

Following normalisation, the mean and standard deviation of the recovered
density are computed numerically from the frFFT grid as
\begin{equation}
    \hat{\mu}_k = \lambda\sum_{j \in \mathcal{S}_k} f(x_j)\,x_j,
    \qquad
    \hat{\sigma}_k = \sqrt{\lambda\sum_{j \in \mathcal{S}_k}
                           f(x_j)\,(x_j - \hat{\mu}_k)^2},
    \label{eq:frfft_moments}
\end{equation}
where $\mathcal{S}_k$ is the index set of grid points within the support.
Table~\ref{tab:step_b} reports these moments alongside the analytical
parameters and the corresponding absolute errors.

\begin{table}[H]
\centering
\caption{Step B: moment accuracy of the frFFT-recovered density for the first 4 trajectories.  The
frFFT mean matches the analytical mean to within $\approx 10^{-7}$ and the
standard deviation to within $\approx 10^{-7}$.}
\label{tab:step_b}
\begin{tabular}{ccccccc}
\toprule
$k$ & $\hat{\mu}_k$ & $\mu_{\mathrm{BS},k}$ & $|\Delta\mu|$
    & $\hat{\sigma}_k$ & $\sigma_{\mathrm{BS},k}$ & $|\Delta\sigma|$ \\
\midrule
0 & 0.8533109 & 0.8533110 & $7.02 \times 10^{-8}$ & 0.0232009 & 0.0232011 & $2.05 \times 10^{-7}$ \\
1 & 0.1042671 & 0.1042671 & $2.94 \times 10^{-8}$ & 0.0459890 & 0.0459892 & $2.38 \times 10^{-7}$ \\
2 & 0.6879009 & 0.6879010 & $7.36 \times 10^{-8}$ & 0.0403612 & 0.0403615 & $3.08 \times 10^{-7}$ \\
3 & 0.7827373 & 0.7827374 & $8.18 \times 10^{-8}$ & 0.0796737 & 0.0796743 & $5.95 \times 10^{-7}$ \\
\bottomrule
\end{tabular}
\end{table}

The absolute error in the mean is at most $8.2 \times 10^{-8}$, while the
error in the standard deviation does not exceed $6.0 \times 10^{-7}$.  Both
figures are several orders of magnitude smaller than $\lambda \approx 0.012$,
the grid spacing of the frFFT output, which represents the fundamental
resolution limit of the discretisation.  The moment errors are therefore
entirely consistent with the quadrature discretisation error of the
Gil-Pelaez inversion and contain no evidence of a systematic implementation
defect.

\paragraph{Step C: Distribution Metrics.}
\label{subsec:step_c}

The final validation step compares the full shape of the frFFT-recovered
density with the analytical Gaussian at every active grid point within the
support.  Because the frFFT grid spacing $\lambda \approx 0.012$ is
considerably coarser than the Gaussian standard deviations
$\sigma_{\mathrm{BS}} \in [0.023, 0.080]$, there are on average only
$2 \times 5\sigma_{\mathrm{BS}} / \lambda \approx 38$ active grid points per
trajectory within the $\pm 5\sigma$ support.  Re-binning onto a finer histogram
would leave most bins empty, artificially inflating divergence metrics such as
Jensen-Shannon distance.  The correct approach is therefore to treat the two
densities as discrete probability vectors evaluated at the \emph{same}
$38$ quadrature points, converting density values to probability masses via
$p_j = f(x_j) \cdot \lambda$.

Four distribution metrics are computed: the Kolmogorov-Smirnov (KS)
statistic $D = \max_j |F_{\mathrm{frFFT}}(x_j) - F_{\mathrm{Gauss}}(x_j)|$,
where $F$ denotes the discrete CDF; the Jensen-Shannon distance
$\mathrm{JS}(p \,\|\, q) = \sqrt{D_{\mathrm{KL}}(p \,\|\, m) / 2
+ D_{\mathrm{KL}}(q \,\|\, m)/2}$ with $m = (p+q)/2$; the total variation
distance $\mathrm{TV} = \tfrac{1}{2}\sum_j |p_j - q_j|$; and the
Hellinger distance $H = \sqrt{\tfrac{1}{2}\sum_j (\sqrt{p_j} - \sqrt{q_j})^2}$.\\
The average values of the distribution metrics between the frFFT density and the
analytical Gaussian for the $10,000$ trajectories are in the order of $10^{-8}$ for the Jensen-Shannon, TV distance and Hellinger distance
while the average value of the KS test's pvalues is around $1.0$ suggesting no relevant difference between the 
analytical distribution and the one obtained from the fractional fft transform.\\

This result has a precise
interpretation: at each of the 38 active quadrature points, the probability
mass $p_j = f_{\mathrm{frFFT}}(x_j) \cdot \lambda$ and
$q_j = f_{\mathrm{Gauss}}(x_j) \cdot \lambda$ agree to within
double-precision floating-point resolution, so no measurable distance exists
between the two discrete distributions.\\

To probe the numerical error more finely the plots of the probability densities obtained from the frFFT and the analytical Gaussian are shown in Figure~\ref{fig:frfft_plots}. 
Moreover, Figure~\ref{fig:frfft_residuals}
plots the pointwise residual
$r(x_k) = f_{\mathrm{frFFT}}(x_k) - f_{\mathrm{Gauss}}(x_k)$
at every active grid point in the support for four randomly sampled trajectories.

\begin{figure}[H]
    \centering
    \includegraphics[width=\textwidth]{Heston/images/fft_plots.png}
    \caption{Comparison of the frFFT-recovered density (red dashed) and the analytical Gaussian (blue) for four randomly sampled trajectories.}
    \label{fig:frfft_plots}
\end{figure}


\begin{figure}[H]
    \centering
    \includegraphics[width=\textwidth]{Heston/images/fft_experiment.png}
    \caption{Pointwise residuals between the normalised frFFT density and the
    analytical Gaussian density, evaluated at the active fine-grid points
    within $\pm 5\sigma$ of each trajectory's mean.  The vertical axes carry
    exponents of $1^{-15}$, confirming that the residuals are at the level of double-precision
    machine epsilon.}
    \label{fig:frfft_residuals}
\end{figure}

The figure reveals several important features.  First, the magnitude of the
residuals is around $1^{-15}$, which lies within one to
two orders of magnitude of the double-precision machine epsilon
$\varepsilon_{\mathrm{mach}} = 2^{-52} \approx 2.22 \times 10^{-16}$. In other words,
the frFFT recovers the Gaussian density \emph{to the last representable bit}.

Second, the residuals are visibly oscillatory and change sign multiple times
across the support, with no discernible systematic trend or bias.  This
sign-changing, zero-mean pattern is the hallmark of floating-point rounding
error, which is essentially random in sign and magnitude, as opposed to
truncation error or quadrature bias, which would manifest as a residual that
is consistently positive or negative, or that grows monotonically toward the
boundaries.  The complete absence of such structured deviations confirms that
the Simpson quadrature weights and the frequency truncation at $j\eta_{\max}
= N_{\mathrm{FFT}} \cdot \eta = 1024$ are adequate for the narrow Gaussian
distributions encountered here: the CF has already decayed to
$|\varphi_{\mathrm{BS}}(1024)| = \exp(-\tfrac{1}{2}\sigma_{\mathrm{BS}}^2
\cdot 1024^2) \approx 10^{-2400}$, making the truncation error entirely
negligible.


\subsection{Summary and Interpretation}
\label{subsec:validation_summary}

The three-step experiment establishes the correctness of the frFFT
implementation with high confidence.  The main findings can be stated as
follows.

\begin{enumerate}
    \item \textbf{Moment accuracy}: the frFFT recovers the mean and standard
    deviation of the Gaussian to within $10^{-7}$, which is well below the
    grid-spacing resolution $\lambda \approx 0.012$, demonstrating that the
    numerical quadrature of the Gil-Pelaez integral is essentially exact.

    \item \textbf{Distributional identity}: all distribution metrics (KS, TV,
    Hellinger) are identically zero at the level of double-precision
    arithmetic, confirming that the recovered discrete probability masses
    agree point-by-point with those of the analytical Gaussian.

    \item \textbf{Machine-epsilon residuals}: the pointwise density residual
    lies between $10^{-15}$ and $10^{-14}$, consistent with the rounding
    error floor of machine precision, with no systematic bias.

    \item \textbf{Aliasing is benign}: the factor-of-two inflation of the raw
    integral, which is a predictable consequence of the choice
    $\alpha = 2/N_{\mathrm{FFT}}$, is fully corrected by the
    normalisation step and has no effect on the final PDF.
\end{enumerate}

These results validate the frFFT inversion routine for use in the Heston
model, where the characteristic function is considerably more complex but the
numerical pipeline is identical.  Since the Black-Scholes CF is an extremely
well-conditioned test case -- it is entire, super-exponentially decaying, and
possesses an exact closed-form inverse -- the experiment places an upper bound
on the algorithmic error: any additional inversion error encountered in the
Heston case will arise from the slower decay of the Heston CF and from the
presence of the stochastic variance term, not from a defect in the frFFT
machinery itself.





















\section{Scoring Rule ForGAN results}

The experiment has been conducted using the Heston trajectories and the frFFT-based PDF estimator described in \ref{sec:frfft_validation}.\\
The training procedure implemented is designed to run on a high-performance computing (HPC) of the CINECA nodes. In particular, the training script is executed on a node equipped with \textbf{8 CPU cores and 1 GPU}, as available on the Leonardo Booster partition as explained in the setup discussion \ref{setup}.\\
The script follows a fully automated pipeline that (i) generates a large synthetic Heston dataset in chunks, (ii) stores it in a binary streaming format, (iii) builds a memory-mapped dataset to avoid loading the full sample set into RAM, (iv) trains the SRForGAN generator on the GPU, and (v) saves the learned weights together with the full training configuration for reproducibility.\\

The code implements the following workflow:
\begin{enumerate}
    \item Generate synthetic Heston trajectories in chunks by using equation described in \ref{eq:milstein_full_trajectory}.
    \item Store each chunk in a flat binary file with a JSON header describing the record size.
    \item Build a memory-mapped dataset that reads trajectories lazily from disk.
    \item Split the dataset into training and validation subsets.
    \item Instantiate the SRForGAN model and select an LSTM generator.
    \item Train the model on the GPU using the Energy Score objective as defined in algorithm \ref{alg:srforgan_train}.
    \item Save the trained weights and configuration in the model directory.
    \item Test the model on unseen data and evaluate performance metrics.
\end{enumerate}

\subsection{Steps for experiment reproducibility}
The first stage constructs a large synthetic dataset based on the Heston stochastic-volatility model. 
The code simulates \text{10,000,000} trajectories, each of length \text{252 + 1}. To avoid memory overflow, the trajectories are generated in chunks of size $100,000$.

Each chunk is produced by the Heston simulator described in previous section \ref{eq:milstein_full_trajectory} and appended to a binary file.
A companion JSON file stores:
\begin{itemize}
    \item the record length, equal to \texttt{N\_STEPS = 252 + 1} where the 252 values are used as the conditioning window and the last value is used as the prediction target,
    \item the number of stored trajectories \text{J = 100,000}.
\end{itemize}

This structure allows the file to be accessed sequentially or through memory mapping without loading the entire dataset into memory.
The following parameters have been used to generate Heston trajectories:
\begin{table}[H]
    \label{tab:heston_training_params}
\centering
\begin{tabular}{lc}
\hline
Parameter & Value \\
\hline
$x_0$      & 0.0 \\
$\mu$      & 0.0 \\
$v_0$      & 0.5 \\
$\kappa$   & 2.0 \\
$\theta$   & 0.5 \\
$\sigma_v$ & 0.9 \\
$\rho$     & -0.7 \\
\texttt{Condition\_trajectory\_length} & 252 \\
\texttt{Forecast\_horizon}     & 10 \\
\hline
\end{tabular}
\caption{Heston parameters}
\end{table}
Where $x_0$ is the initial value of the log-price process, $\mu$ is the drift, $v_0$ is the initial variance, $\kappa$ is the mean-reversion speed of the variance, $\theta$ is the long-term mean of the variance, $\sigma_v$ is the volatility of the variance process, and $\rho$ is the correlation between the Brownian motions driving the log-price and variance.
From the parameters values, the Feller condition is satisfied (in fact, $\kappa \theta = 2.0 \times 0.5 = 1.0 > 0.9 = \sigma_v^2$), ensuring that the variance process remains positive. 


\begin{algorithm}[H]
\caption{Chunked generation of the Heston dataset}
\label{alg:heston_experiment}
\begin{algorithmic}[1]
\REQUIRE parameter ranges, time horizon $T$, number of steps $N$, total number of trajectories $J_{total}$, chunk size $J$
\STATE create \texttt{data/} and \texttt{models/} directories if needed
\STATE initialize binary file and JSON header
\FOR{$i = 1$ to $J_{total}/J$}
    \STATE simulate $J$ Heston trajectories with fixed seed $42+i$
    \STATE convert trajectories to 32-bit floating point format
    \STATE append the raw bytes to the binary file
    \STATE update the number of stored records in the JSON header
\ENDFOR
\end{algorithmic}
\end{algorithm}

A memory-mapped dataset is constructed from the binary file. Each stored record corresponds to a full trajectory
\[
[X_0, X_1, \dots, X_N],
\]
where the first $N$ values are used as the conditioning window and the last value is used as the prediction target. In other words, the dataset returns:
\[
c = [X_0, X_1, \dots, X_{N-10}], \qquad y = X_N.
\]


\begin{algorithm}[H]
\caption{Lazy access to the memory-mapped dataset}
\label{alg:memmap_dataset}
\begin{algorithmic}[1]
\REQUIRE binary trajectory file, JSON header with record length and number of records
\STATE read metadata from the JSON header
\STATE store the file path and defer memmap creation
\STATE \textbf{function} getitem(index)
\IF{memmap has not been opened yet}
    \STATE open the binary file as a NumPy memmap
\ENDIF
\STATE read trajectory $[X_0, \dots, X_N]$ at the requested index
\STATE set condition $c \leftarrow [X_0, \dots, X_{N-10}]$
\STATE set target $y \leftarrow X_N$
\STATE \textbf{return} $(y,c)$
\end{algorithmic}
\end{algorithm}


Once the dataset has been created, it is split into training and validation subsets using $80\%$ of the dataset for training and the remaining $20\%$ for validation.
The split is randomized with a fixed seed to preserve reproducibility.\\

The implemented SRForGAN model \ref{alg:srforgan_train} is initialized using the following configuration:
\begin{itemize}
    \item latent noise dimension \texttt{Z\_NOISE\_DIM = 64},
    \item scoring rule \texttt{energy} \ref{alg:energy_score},
    \item early stopping patience $ 10$ epochs,
    \item batch size  $128$,
    \item learning rate \texttt{$LR = 10^{-3}$}.
\end{itemize}

Moreover, given the complexity of the Heston process with respect to the Black-Scholes case, 
the generator architecture is designed to capture temporal dependencies in the conditioning window
by using the LSTM variant of the generator. In this variant the generator first extract hidden states from the conditioning window through an LSTM layer and then concatenates the final hidden state with 
the latent noise vector which will become the input of a feedforward network to produce the output prediction.
The hidden dimension of the LSTM layer is set to 64.\\

\begin{table}[H]
\centering
\caption{Generator Architecture and Hyperparameters Summary}
\label{tab:srforgan_architecture}
\begin{tabular}{lll}
\hline
\textbf{Parameter}         & \textbf{Generator}                  & \textbf{Critic (Discriminator)} \\ \hline
Input Dimension            & 64 (Latent) + 252 (Cond.)           & --                              \\
LSTM Layers                & 128                                 & --                              \\
Feedforward Layers         & 128+64                                 & --                              \\
Output Dimension           & 1                                  & --                              \\
Activation Function        & Leaky ReLU                         & --                              \\
Dropout                    & 0.0                                & --                              \\
\multicolumn{3}{c}{\textbf{Global Training Settings}}                                      \\ \hline
Batch Size                 & \multicolumn{2}{l}{32}                                            \\
max\_epochs                & \multicolumn{2}{l}{100}                                          \\
early\_stopping\_patience & \multicolumn{2}{l}{10}                                           \\
learning\_rate            & \multicolumn{2}{l}{$10^{-3}$}                           \\         \\
\hline
\end{tabular}
\end{table}



\subsection{Test results}
The trained model is evaluated on a test set given by \text{100,000} trajectories that were not seen during training by using the same Heston parameters adopted for the training process and described in table \ref{tab:heston_training_params}. 
The evaluation metrics are computed by comparing the distribution of the generated predictions with the distribution of the true targets in the test set, using the Heston Monte Carlo simulations (and the frFFT-based PDF estimator for the visualization) described in \ref{sec:frfft_validation}.\\

\begin{table}[h]
\centering
\begin{tabular}{lc}
\hline
Metric & Value \\
\hline
KS test acceptance rate & 87.500\% \\
Jensen--Shannon distance & 0.121850 \\
Hellinger distance & 0.104270 \\
Total variation distance & 0.102780 \\
Earth mover's distance & 0.000707 \\
\hline
\end{tabular}
\caption{Distribution comparison metrics}
\end{table}




\begin{figure}[H]
    \centering
    \includegraphics[width=1.0\textwidth]{Heston/images/srforgan_samples_2.png}
    \caption{Samples obtained using fixed $x_0$, $\mu$, $v_0$, $\theta$, $\kappa$, $\sigma_v$, $\rho$.}
    \label{srforgan_heston_sample2}
\end{figure}



To validate the distributional predictions of the generative model against the
theoretical Heston dynamics, an analytical consistency check is performed on the
conditional mean and standard deviation of the ten-step-ahead log-price
distribution.  The analysis compares three quantities: the theoretical moments
derived from the Heston characteristic function, the empirical moments of the
Monte Carlo ground-truth distribution, and the empirical moments of the
generated distribution.\\

Under the Heston model the log-price follows the Stochastic Differential Equation described in equation \ref{eq:logprice_sde_heston},
so that the exact conditional mean over a forecast horizon $\tau$ is
\begin{equation}
    \mathbb{E}\bigl[X_{T+\tau}\mid X_T,\,v_T\bigr]
    \;=\;
    X_T \;+\; \mu\tau \;-\; \tfrac{1}{2}\,\bar{v}\,\tau,
    \qquad
    \bar{v} = \frac{\mathbb{E}\!\left[\int_0^\tau v_{T+s}\,ds\,\Big|\,v_T\right]}{\tau},
    \label{eq:heston_mean}
\end{equation}
where the integrated variance is given in closed form by the CIR expectation
\begin{equation}
    \mathbb{E}\!\left[\int_0^\tau v_{T+s}\,ds\,\Big|\,v_T\right]
    \;=\;
    \theta\tau + (v_T - \theta)\,\frac{1 - e^{-\kappa\tau}}{\kappa}.
    \label{eq:integral_v}
\end{equation}
Since $\mu = 0$ in our parametrisation, the only deviation from $X_T$ is the
It\^{o} correction $-\tfrac{1}{2}\mathbb{E}[\int v_s\,ds]$, which is always
negative.  For the ten-step horizon $\tau = 10/252 \approx 0.040$\,yr with
$\kappa\tau = 0.079 \ll 1$, equation~\eqref{eq:integral_v} reduces to
$\approx v_T\tau$ with negligible error, so the mean shift from $X_T$ is
$-\tfrac{1}{2}v_T\tau$.

The theoretical standard deviation is dominated by the same integrated variance:
\begin{equation}
    \sigma_{\mathrm{th}} \;=\; \sqrt{\,\mathbb{E}\!\left[\int_0^\tau v_{T+s}\,ds\right]},
    \label{eq:std_th}
\end{equation}
which equals $\sqrt{v_T\tau}$ to leading order.  The vol-of-vol correction
term, proportional to $\sigma_v^2\tau^2$, contributes an additional
$\approx 12.5\%$ to the standard deviation across all four trajectories;
this small and uniform inflation reflects the brevity of the horizon relative
to the mean-reversion timescale.\\


Table~\ref{tab:mean_consistency} reports the theoretical mean from
equation~\eqref{eq:heston_mean}, together with the empirical means of the
true and generated Monte Carlo distributions.

\begin{table}[H]
\centering
\caption{Conditional mean consistency.  $\hat{\mu}_{\mathrm{true}}$ and
$\hat{\mu}_{\mathrm{gen}}$ are the empirical means of the Monte Carlo ground
truth and of the generated distribution, respectively.
$|\Delta_{\mathrm{true}}|=|\hat{\mu}_{\mathrm{true}}-\mu_{\mathrm{th}}|$,
$|\Delta_{\mathrm{gen}}|=|\hat{\mu}_{\mathrm{gen}}-\mu_{\mathrm{th}}|$.}
\label{tab:mean_consistency}
\begin{tabular}{cccccccc}
\toprule
Traj. & $X_T$ & $v_T$ & $\mu_{\mathrm{th}}$
      & $\hat{\mu}_{\mathrm{true}}$ & $|\Delta_{\mathrm{true}}|$
      & $\hat{\mu}_{\mathrm{gen}}$\phantom{$|$}
      & $|\Delta_{\mathrm{gen}}|$ \\
\midrule
0 & $-1.219$ & $0.879$ & $-1.236$ & $-1.237$ & $0.001$ & $-1.243$ & $0.006$\\
1 & $-0.975$ & $1.094$ & $-0.996$ & $-0.983$ & $0.014$ & $-1.086$ & $0.09$\\
2 & $\phantom{-}0.721$ & $0.367$  & $\phantom{-}0.714$ & $\phantom{-}0.712$ & $0.001$ & $\phantom{-}0.675$ & $0.039$ \\
3 & $-0.159$ & $0.498$ & $-0.169$ & $-0.174$ & $0.006$ & $-0.147$ & $0.022$\\
\bottomrule
\end{tabular}
\end{table}

The true distribution means match the theoretical predictions to within
$0.014$ in all cases, confirming that the Monte Carlo simulation correctly
implements the Heston dynamics.\\
A noteworthy feature of equation~\eqref{eq:heston_mean} is that, with $\mu=0$,
the distribution centre lies slightly \emph{left} of the current log-price
$X_T$ by the It\^{o} correction $-\tfrac{1}{2}v_T\tau \in [-0.021,\,-0.007]$.\\


Table \ref{tab:std_consistency} compares the theoretical standard deviation
from equation~\eqref{eq:std_th} against the empirical counterparts.

\begin{table}[H]
\centering
\caption{Conditional standard deviation consistency.
$\sigma_{\mathrm{th}} = \sqrt{\mathbb{E}[\int_0^\tau v_s\,ds]}$;
$\hat{\sigma}_{\mathrm{true}}$ and $\hat{\sigma}_{\mathrm{gen}}$ are the
empirical standard deviations of the true and generated distributions.}
\label{tab:std_consistency}
\begin{tabular}{cccccccc}
\toprule
Traj. & $v_T$ & $\sqrt{v_T\tau}$ & $\sigma_{\mathrm{th}}$
      & $\hat{\sigma}_{\mathrm{true}}$ & $|\Delta_{\mathrm{true}}|$
      & $\hat{\sigma}_{\mathrm{gen}}$ & $|\Delta_{\mathrm{gen}}|$ \\
\midrule
0 & $0.879$ & $0.1868$ & $0.1852$ & $0.1830$ & $0.0022$ & $0.2068$ & $0.0216$ \\
1 & $1.094$ & $0.2084$ & $0.2062$ & $0.2129$ & $0.0067$ & $0.1989$ & $0.0073$ \\
2 & $0.367$ & $0.1207$ & $0.1215$ & $0.1227$ & $0.0012$ & $0.1160$ & $0.0055$ \\
3 & $0.498$ & $0.1406$ & $0.1406$ & $0.1443$ & $0.0037$ & $0.1377$ & $0.0029$ \\
\bottomrule
\end{tabular}
\end{table}

The true empirical standard deviations deviate from the theoretical values
by at most $0.007$, well within Monte Carlo sampling noise at $10^3$
simulations.  Crucially, the \emph{ordering} imposed by $v_T$ is preserved:
trajectory~2 ($v_T = 0.367$, lowest variance) yields the narrowest
distribution and trajectory~1 ($v_T = 1.094$, highest) the widest, in both
the theoretical and empirical cases.  The generated standard deviations are also
in reasonable agreement with theory, with trajectory~0 exhibiting the largest
absolute deviation ($0.022$).\\


The analytical assessment confirms full consistency between the Heston
parameter set, the terminal states $(X_T,\,v_T)$, and the Monte Carlo ground-truth
distributions: theoretical means match empirical means to within $0.014$ and
theoretical standard deviations to within $0.007$.  The generated distributions
replicate the correct distributional shape and ordering as suggested by the 
distance metrics and Kolmogorov-Smirnov test as well.\\


The same kind of experiment is run on Heston dataset obtained by allowing 
the different parameters to change.\\
The first set of experiments allows just the volatility of the volatility to change. In this case the parameters of the trajectories are sampled as follows:
\begin{table}[H]
\centering
\begin{tabular}{lcc}
\hline
Parameter & Value training & Value testing \\
\hline
$x_0$      & 0.0 & 0.0\\
$\mu$      & 0.0 & 0.0\\
$v_0$      & 0.5 & 0.5\\
$\kappa$   & 2.0 & 2.0\\
$\theta$   & 0.5 & 0.5\\
$\sigma_v$ & (0.2, 1.3) & (0.3, 1.2) \\
$\rho$     & -0.7 & -0.7\\
\texttt{Condition\_trajectory\_length} & 252 & 252\\
\texttt{Forecast\_horizon}     & 10 & 10\\
\hline
\end{tabular}
\caption{Heston parameters vol-of-vol}
\end{table}

The performance metrics obtained over \text{10000} samples are slightly decreased
with respect to the fixed parameters experiment:

\begin{table}[h]
\centering
\begin{tabular}{lc}
\hline
Metric & Value \\
\hline
KS test acceptance rate & 83.00\% \\
Jensen--Shannon distance & 0.149905 \\
Hellinger distance & 0.128144 \\
Total variation distance & 0.128569 \\
Earth mover's distance & 0.001136 \\
\hline
\end{tabular}
\caption{Distribution comparison metrics}
\end{table}

\begin{figure}[H]
    \centering
    \includegraphics[width=1.0\textwidth]{Heston/images/srforgan_samples_volvol.png}
    \caption{Samples obtained using fixed $x_0$, $\mu$, $v_0$, $\theta$, $\kappa$, $\rho$ and variable $\sigma_v$}
    \label{srforgan_heston_sample_volvol}
\end{figure}

The last experiment allows all the different parameters to change for each trajectory,
while keeping the Feller condition always satisfied, 
in particular:

\begin{table}[H]
\centering
\begin{tabular}{lcc}
\hline
Parameter & Value training & Value testing \\
\hline
$x_0$      & 0.0  & 0.0 \\
$\mu$      & 0.0 & 0.0 \\
$v_0$      & 0.5 & 0.5 \\
$\kappa$   & (1.5, 5.0) & (2.0, 4.0) \\
$\theta$   & (0.57, 1.5) & (0.6, 1.4) \\
$\sigma_v$ & (0.2, 1.3) & (0.3, 1.2) \\
$\rho$     & -0.7 & -0.7 \\
\texttt{Condition\_trajectory\_length} & 252 & 252 \\
\texttt{Forecast\_horizon}     & 10 & 10\\
\hline
\end{tabular}
\caption{Heston with all variable parameters}
\end{table}


The performance metrics does not differ significantly from the previous experiment where just the vol-of-vol 
is allowed to change for each trajectory. This suggests that the model is able to 
model also trajectories guided by different underlying parameters 
\begin{table}[h]
\centering
\begin{tabular}{lc}
\hline
Metric & Value \\
\hline
KS test acceptance rate & 82.10\% \\
Jensen--Shannon distance & 0.150024 \\
Hellinger distance & 0.129049 \\
Total variation distance & 0.129361 \\
Earth mover's distance & 0.001198 \\
\hline
\end{tabular}
\caption{Distribution comparison metrics}
\end{table}



\begin{figure}[H]
    \centering
    \includegraphics[width=1.0\textwidth]{Heston/images/srforgan_samples_all_params.png}
    \caption{Samples obtained using fixed $x_0$, $\mu$, $v_0$, but different $\theta$, $\kappa$, $\rho$ and $\sigma_v$}
    \label{srforgan_heston_sample_all_var}
\end{figure}



\section{Pricing-Based Validation of the Forecast Distribution}
\label{sec:pricing_validation}


The distributional metrics and experiments ran before quantify how close 
the generated distribution $P^\phi(\cdot\mid\mathbf c)$
is to the true distribution$P^\star(\cdot\mid\mathbf c)$ in a generic, application-agnostic environment.\\
The following experiments in this section evaluate the SRForGAN generated distribution in a pricing setting by using European options.
If $G_\phi$ has genuinely learned the conditional law of
$S_{t+l}$, discounted expectations of option payoffs computed under
$P^\phi$ should agree with the corresponding Heston-model prices, across
a grid of strikes spanning the smile.\\

A single price comparison, may conflates two distinct sources of
discrepancy, that is, the sampling noise inherent to any Monte Carlo estimator of
an expectation, and genuine mismatch between $P^\phi$ and $P^\star$. To
disentangle them, three pricing engines are used side by side.
\begin{itemize}
    \item An \emph{oracle} price is obtained semi-analytically from the Heston
characteristic function using the Fourier inversion method already described in \ref{alg:frFTT} 
which is the closest available proxy to the true,
noise-free price, but one that requires knowledge of the latent
instantaneous variance $v_t$, which $G_\phi$ never observes.
    \item A \emph{Heston Monte Carlo} price is obtained with the same
discounted-payoff sample-average estimator as the GAN price, but drawing
from the true Heston dynamics started from the
same $(S_t,v_t)$ like in \ref{alg:heston_mc_pdf}.
    \item A \emph{GAN} price is obtained with the identical
sample-average estimator, drawing instead from $G_\phi(\cdot\mid\mathbf
c)$.
\end{itemize}
   
Because the Heston-MC and GAN routes share the same estimator
family, their standard errors are directly comparable and
$\text{price}^{\text{GAN}}-\text{price}^{\text{MC}}$ isolates the
generator's distributional error net of Monte Carlo noise; because the
oracle and Heston-MC routes both target the true Heston distribution,
$\text{price}^{\text{MC}}-\text{price}^{\text{oracle}}$ measures pure
numerical noise and acts as a control, in fact, a large gap here would flag a
problem with the benchmark itself rather than with $G_\phi$.

\subsection{Experimental design}

\paragraph{Test set.} $J=300$ fresh Heston trajectories (Milstein scheme)
are simulated with parameters drawn from the training ranges, that allow the Feller condition to be always satisfied, of the previous experiment,
$\kappa\in[1.5,5.0]$, $\theta\in[0.57,1.5]$, $\sigma_v\in[0.2,1.3]$, and
$\rho=-0.7$, $v_0=0.5$ fixed. Each
trajectory carries a conditioning window of $k=252$ observed log-price
steps ($\approx1$ trading year) and a forecast horizon of $l=10$ steps
($\approx2$ weeks), matching the training configuration. The reference
time $t$ is the last index of the conditioning window; $S_t=e^{X_t}$ and
$v_t$ (both known from simulation, but $v_t$ withheld from $G_\phi$)
anchor all three pricing engines to the same starting state.

\paragraph{Strike grid.} For every trajectory, strikes are built on a
common moneyness grid $\mathcal M=\{0.80,0.90,0.95,1.00,1.05,1.10,1.20\}$,
$K=m\cdot S_t$, so each trajectory contributes the same relative coverage
of the smile regardless of its own realised spot level. Following
standard smile-construction practice, options are priced as puts for
$K<S_t$ and calls for $K\ge S_t$, keeping every quoted contract
out-of-the-money and away from the numerically unstable deep-ITM region
of Black--Scholes implied-volatility inversion.

\begin{table}[H]
\centering
\begin{tabular}{cc}
\hline
$K/S_t$ & \textbf{Contract} \\
\hline
0.80 & Put \\
0.90 & Put \\
0.95 & Put \\
1.00 & Call \\
1.05 & Call \\
1.10 & Call \\
1.20 & Call \\
\hline
\end{tabular}
\end{table}



\paragraph{Discounting.} The risk-free rate is set to $r=0$, matching the
risk-neutral drift $\mu=0$ used throughout the simulated data, so $S_t$
is a martingale under both the true and the forecast measure and the
three engines remain directly comparable.

\paragraph{Oracle price.} Writing $\tau=l\Delta t$, the Heston
characteristic function $\varphi(u)=\mathbb E[e^{iuX_{t+\tau}}\mid
X_t,v_t]$ is inverted with a fractional FFT into a discretised density
$f(x)$ of the terminal log-price.
\[
C^{\text{oracle}}(K)=e^{-r\tau}\int (e^{x}-K)^+ f(x)\,dx .
\]

Similarly for the put options with $K<S_t$:

\[
P^{\text{oracle}}(K)=e^{-r\tau}\int (K-e^{x})^+ f(x)\,dx
\]

During the rest of the discussion the notation $\hat C^{\text{°}}(K)$ will be used to refer to both call and put options, with the understanding that the payoff function is chosen according to the moneyness of the strike $K$.


\paragraph{Heston Monte Carlo price.} $10{,}000$ terminal log-price draws
per trajectory are simulated with the exact/QE CIR scheme already used
in this thesis, started from $(X_t,v_t)$:
\[
\hat C^{\text{MC}}(K)=\frac{e^{-r\tau}}{N}\sum_{i=1}^N
\left(e^{X_{t+\tau}^{(i)}}-K\right)^{+},\qquad
\widehat{\text{se}}^{\text{MC}}(K)=\frac{e^{-r\tau}}{\sqrt N}\,
\widehat{\mathrm{sd}}\!\left[\left(e^{X_{t+\tau}}-K\right)^{+}\right].
\]


\paragraph{GAN price.} The identical estimator is applied to $N=10{,}000$
draws $\hat X_{t+\tau}^{(i)}\sim G_\phi(\cdot\mid\mathbf c)$ from the
trained generator, conditioned only on the observed price window.
Similarly to the Heston-MC case, the standard error of the GAN price is estimated as:
\[
\hat C^{\text{GAN}}(K)=\frac{e^{-r\tau}}{N}\sum_{i=1}^N
\left(e^{X_{t+\tau}^{(i)}}-K\right)^{+},\qquad
\widehat{\text{se}}^{\text{GAN}}(K)=\frac{e^{-r\tau}}{\sqrt N}\,
\widehat{\mathrm{sd}}\!\left[\left(e^{X_{t+\tau}}-K\right)^{+}\right].
\]

\paragraph{Implied volatility.} Every price (oracle, MC, GAN) is
converted into a Black--Scholes implied volatility $\sigma_{\text{impl}}$
by Brent root-finding. By using this monotone re-parametrisation the prices errors are placed
 on a scale comparable across strikes and that directly exposes smile-shape mismatches, which is the
feature a stochastic-volatility model is specifically built to capture.

\paragraph{Error and diagnostic metrics.} For each (trajectory, strike):
price and IV errors of GAN vs.\ oracle, GAN vs.\ MC, and MC vs.\ oracle;
and a $z$-score
testing whether the GAN/MC price gap exceeds sampling noise. Separately,
for each trajectory the first four empirical moments of the raw MC and
GAN log-price samples (mean, standard deviation, skewness, excess
kurtosis) are compared directly, to diagnose which feature of the
distribution drives any pricing mismatch.\\

Let $i=1,\dots,n$ ($n=2{,}100$) index the
(trajectory, strike) pairs underlying Table~\ref{tab:overall}. For pair
$i$, let $\hat C_i^{\text{oracle}}$, $\hat C_i^{\text{MC}}$,
$\hat C_i^{\text{GAN}}$ denote the oracle, Heston-MC, and GAN prices, and
$\hat\sigma_i^{\text{oracle}}$, $\hat\sigma_i^{\text{GAN}}$ the
corresponding Black--Scholes implied volatilities. Define the residuals
\[
e_i = \hat C_i^{\text{GAN}}-\hat C_i^{\text{oracle}}, \qquad
e_i^{\text{MC}} = \hat C_i^{\text{GAN}}-\hat C_i^{\text{MC}}, \qquad
\varepsilon_i = \hat\sigma_i^{\text{GAN}}-\hat\sigma_i^{\text{oracle}},
\]
and the standardised price gap
\[
z_i = \frac{\hat C_i^{\text{GAN}}-\hat C_i^{\text{MC}}}
{\sqrt{\big(\widehat{\text{se}}_i^{\text{GAN}}\big)^2+\big(\widehat{\text{se}}_i^{\text{MC}}\big)^2}},
\]
where $\widehat{\text{se}}_i^{\text{GAN}}$ and $\widehat{\text{se}}_i^{\text{MC}}$
are the Monte Carlo standard errors of the GAN and Heston-MC price
estimators defined in the "Heston Monte Carlo price"/"GAN price"
paragraphs above, and $\mathbf 1\{\cdot\}$ is the indicator function.

\subsection{Results}

Table~\ref{tab:overall} reports the aggregate error statistics over all
$J\times|\mathcal M|=2{,}100$ (trajectory, strike) pairs.
Table~\ref{tab:moneyness} breaks these down by moneyness, and
Table~\ref{tab:moments} compares the empirical moments of the GAN and
Heston-MC terminal distributions.\\



\begin{table}[H]
\centering
\small
\caption{Aggregate pricing-error statistics ($J=300$ trajectories,
$7$ strikes, $n=2{,}100$).}
\label{tab:overall}
\begin{tabular}{lcc}
\toprule
Metric & Formula & Value \\
\midrule
Bias, price vs.\ oracle          & $\dfrac1n\sum_i e_i$                                   & 0.0005 \\[1.2ex]
\hline
RMSE, price vs.\ oracle          & $\sqrt{\dfrac1n\sum_i e_i^2}$                         & 0.0312 \\[1.2ex]
\hline
WAPE, price vs.\ oracle           & $\frac{\sum \vert{}e_i\vert{}}{\sum C_i^{\text{oracle}}}$  & 0.3718 \\[1.2ex]
\hline
RMSE, price vs.\ Heston MC       & $\sqrt{\dfrac1n\sum_i \big(e_i^{\text{MC}}\big)^2}$    & 0.0311 \\[1.2ex]
\hline
WAPE, price vs.\ Heston MC        & $\frac{\sum \vert{}e_i^{MC}\vert{}}{\sum C_i^{\text{oracle}}}$             & 0.3731 \\[1.2ex]
\hline
Bias, implied vol.\ vs.\ oracle  & $\dfrac1n\sum_i \varepsilon_i$                         & 0.0045 \\[1.2ex]
\hline
RMSE, implied vol.\ vs.\ oracle  & $\sqrt{\dfrac1n\sum_i \varepsilon_i^2}$                & 0.3320 \\[1.2ex]
\hline
WAPE, implied vol.\ vs.\ oracle   & $\frac{\sum \vert{}\varepsilon_i\vert{}}{\sum \hat\sigma_i^{\text{oracle}}}$                      & 0.2601 \\[1.2ex]
\hline
Share of $|z_{\text{GAN vs.\ MC}}|>1.96$ & $\dfrac1n\sum_i \mathbf 1\{|z_i|>1.96\}$       & 91.3\% \\
\bottomrule
\end{tabular}
\end{table}


\begin{table}[H]
\centering
\caption{Pricing-error statistics by moneyness $K/S_t$. \\(*** = $p<0.01$, ** = $p<0.05$, * = $p<0.1$).}
\label{tab:moneyness}
\begin{tabular}{lrrrr}
\toprule
$K/S_t$ & Bias, price & RMSE, IV & WAPE, IV & Bias, IV \\
\midrule
0.80 & $-0.0023$\text{***} & 0.2256 & 0.4696 & $-0.0574$\text{***}\\
0.90 & $-0.0041$\text{***} & 0.2861 & 0.3494 & $-0.0534$\text{**} \\
0.95 & $-0.0048$\text{**} & 0.3300 & 0.2961 & $-0.0517$\text{**} \\
1.00 & $\phantom{-}0.0049$\text{*} & 0.2688 & 0.3557 & $\phantom{-}0.0662$\text{*} \\
1.05 & $\phantom{-}0.0042$\text{**} & 0.3133 & 0.2951 & $\phantom{-}0.0541$\text{*} \\
1.10 & $\phantom{-}0.0035$\text{**} & 0.3679 & 0.2522 & $\phantom{-}0.0432$\text{**} \\
1.20 & $\phantom{-}0.0021$\text{***} & 0.4851 & 0.1985 & $\phantom{-}0.0303$\text{***} \\
\bottomrule
\end{tabular}
\end{table}

\begin{table}[H]
\centering
\caption{GAN vs.\ Heston-MC terminal log-price moments (bias
$=\text{GAN}-\text{MC}$, averaged over $J=300$ trajectories; correlation
across trajectories).}
\label{tab:moments}
\begin{tabular}{lrrr}
\toprule
Moment & Bias & WAPE & Correlation \\
\midrule
Mean               & 0.0095  & 0.0516 & 0.998 \\
Std.\ deviation     & $-0.0103$ & 0.101 & 0.814 \\
Skewness            & 0.0725  & 0.4964 & $-0.099$ \\
Excess kurtosis     & 0.1078  & 1.5267 & 0.374 \\
\bottomrule
\end{tabular}
\end{table}

The oracle and Heston-MC benchmarks agree closely with each other, both
in aggregate (RMSE $0.0312$ vs.\ $0.0311$) and at every moneyness level;
this confirms the two ground-truth routes are mutually consistent, so
essentially all the reported GAN discrepancy reflects genuine
distributional error rather than benchmark disagreement. The aggregate
price bias is small ($0.0005$), but Table~\ref{tab:moneyness} shows this
is the average of a sign-flipping pattern: prices are systematically
\emph{below} the oracle for $K<S_t$ (puts) and \emph{above} it for
$K\ge S_t$ (calls), with the effect largest near the money. For $91.3\%$
of (trajectory, strike) pairs the GAN/MC price gap exceeds what sampling
noise alone would explain at the $95\%$ level ($|z|>1.96$), so the
mismatch is not an artefact of Monte Carlo variance. 
Moreover, for each moneyness level the bias in the prices and in the implied volatilities is statistically significant at least at the $10\%$ confidence level.\\

Table~\ref{tab:moments} identifies the mechanism. The GAN reproduces the
trajectory-conditional \emph{mean} of the terminal log-price almost
perfectly (correlation $0.998$) and the \emph{standard deviation}
reasonably well (correlation $0.814$), but shows essentially no
association with the true \emph{skewness} ($-0.099$) or \emph{kurtosis}
($0.374$) across trajectories. A small positive mean bias ($+0.0095$,
under $10\%$ of one standard deviation at this horizon) shifts the whole
distribution to the right, mechanically depressing put prices and
inflating call prices, exactly the sign pattern observed. The
Heston leverage effect ($\rho=-0.7$) induces genuine, trajectory-specific
negative skew in $P^\star$; since the GAN does not track skewness
across regimes, its implied smile comes out flatter and less negatively
skewed than the truth, which is precisely why implied-volatility errors
are largest at the money (where price is most sensitive to dispersion
and to skew-driven asymmetry) and shrink towards both wings, rather than
following the pattern of noisier deep-OTM/ITM inversion that one might
otherwise expect. A plausible contributing cause is that Energy-Score
training here used only $n_{\text{samples}}=10$ generator draws per
gradient step: mean and variance are estimated far more precisely than
skewness or kurtosis at that sample size, so the training signal is
inherently weaker for exactly the moments found deficient here. 
In addition, no hyperparameter tuning was performed to improve the GAN's ability to capture higher-order moments.



\chapter{Conclusion}

This thesis proposes to test and evaluate whether generative adversarial networks,
and their adversarial-free variant, can be used to learn the conditional distribution of a financial time series.
The proposed scoring rule based final model, named SRForGAN, has been tested to check its capabilities to learn the
forecast distribution of a financial time series accurately enough to be
useful, moving progressively from Geometric Brownian Motion to the
stochastic-volatility Heston model, and ultimately evaluating the learned
distributions in a concrete option pricing task.\\

The results consistently favour the proposed approach over its
GAN-based alternatives. Under Black--Scholes dynamics, the vanilla ForGAN
trained with the BCE loss struggled visibly (25\% KS acceptance,
Jensen--Shannon distance $\approx0.49$), the Wasserstein variant improved
matters but remained far from satisfactory (34\% KS acceptance,
JS $\approx0.28$), while SRForGAN, trained by minimising the Energy
Score, with no critic and no adversarial game, matched the analytical
Gaussian benchmark almost exactly (99\% KS acceptance, JS $\approx0.14$).
This gap is the central practical advantage demonstrated throughout the
thesis: replacing the adversarial min-max objective with a strictly proper
scoring rule removes the well-known instability and mode-collapse
tendencies of GAN training, at no cost in distributional accuracy and, in
this setting, at a considerable gain. The same architecture and training
recipe then transferred, with only a modest and expected drop in
performance, to the substantially harder Heston setting: 87.5\% KS
acceptance with fixed parameters, and 82--83\% once $\kappa$, $\theta$,
$\sigma_v$ and $\rho$ were all allowed to vary across trajectories, showing
that a single trained generator can reasonably generalise across a family
of volatility regimes rather than requiring re-estimation for each one,
which is precisely the property one would want from a model intended for
practical deployment. The machine-epsilon-level validation of the frFFT
inversion pipeline further ensures that all of these comparisons rest on a
numerically trustworthy Heston benchmark rather than on an artefact of the
evaluation code itself.

The pricing-based validation introduced in the final chapter is arguably
the thesis's most informative contribution, precisely because it complicates
this otherwise positive picture. Distributional distances such as KS,
Jensen--Shannon, Hellinger and Total Variation are dominated by the bulk of
the distribution and, as such, rewarded SRForGAN generously. Pricing a
volatility smile, by contrast, is disproportionately sensitive to exactly
the features that live in the tails and in the asymmetry of the
distribution. The moment diagnostics make the mechanism explicit: the
generator reproduces the trajectory-conditional mean of the terminal
log-price almost perfectly (correlation $0.998$) and the standard deviation
well (correlation $0.814$), but shows essentially no ability to track
skewness ($-0.189$) or excess kurtosis ($0.264$) across the different Heston
regimes. Since the leverage effect ($\rho=-0.7$) is precisely what endows
the true distribution with a regime-dependent negative skew, its absence in
the generated distribution produces a smile that is systematically flatter
than the truth, with a statistically significant bias (at worst $p<0.001$)
that underprices out-of-the-money puts and overprices out-of-the-money
calls. In other words, a model that looks close to indistinguishable from
the ground truth under generic distributional metrics can still be
economically unsuitable for a skew-sensitive financial application, a
distinction that only becomes visible once the forecast is stress-tested
against the task it is meant to serve, and one that future generative
forecasting work in finance should not overlook.

This shortcoming also delineates the clearest avenues for future work.
Training used only ten Monte Carlo draws per gradient step to estimate the
Energy Score, a sample size sufficient to estimate a mean or a variance
precisely but a considerably weaker basis for estimating skewness or
kurtosis; whether increasing this sample size, or adopting a scoring rule
more explicitly sensitive to higher-order moments (for instance a variogram
score or a weighted combination of scoring rules), closes this gap is a
natural and immediately testable next step, since no hyperparameter tuning
was directed at this issue in the present work. Likewise, the generator was
deliberately restricted to conditioning on the observed price history
alone, mirroring realistic deployment where the instantaneous variance is
unobservable; enriching the conditioning set with realised-volatility
proxies, extending the analysis to longer horizons and more extreme
parameter regimes, and testing the same pricing-based methodology on other
volatility models are all promising directions building directly on the
framework developed here.
%conclusion
\begin{appendices}


\chapter{Joint distributions as transport plans}


The key idea is that a \textbf{joint distribution} does more than describe two random variables simultaneously---it specifies \textbf{how probability mass from one distribution is matched to probability mass in another}. That matching is exactly what is meant by a \textbf{transport plan} in optimal transport.\\

Suppose you have

\begin{itemize}
    \item a real data distribution $\mu$ on $\mathcal{X}$,
    \item a generated distribution $\nu$ on $\mathcal{Y}$.
\end{itemize}

By themselves, these distributions tell you only \textbf{where the mass is located}. They do \textbf{not} tell you which part of the mass in $\mu$ should be transported to which part of the mass in $\nu$.

A \textbf{joint distribution}

\[
\pi(x,y)
\]

defined on the product space

\[
\mathcal{X}\times\mathcal{Y}
\]

provides exactly this missing information.\\


A transport plan specifies the amount of mass moved from every point $x$ to every point $y$.

More precisely,

\[
\pi(A\times B)
\]

represents the amount of probability mass originally in the subset $A\subseteq\mathcal{X}$ that is transported to the subset $B\subseteq\mathcal{Y}$.

Equivalently, if $\pi$ has a density, then

\[
\pi(x,y)\,dx\,dy
\]

is the infinitesimal amount of mass transported from a neighborhood around $x$ to a neighborhood around $y$.

Thus every pair $(x,y)$ corresponds to one possible transportation route, and the joint distribution tells us how much mass uses each route.\\

A transport plan cannot create or destroy mass.

Therefore,

\[
\int_{\mathcal Y}\pi(x,y)\,dy=\mu(x)
\]

states that all mass leaving $x$ equals the amount of mass originally located at $x$.

Similarly,

\[
\int_{\mathcal X}\pi(x,y)\,dx=\nu(y)
\]

states that all mass arriving at $y$ equals the amount of mass required at $y$.

These are exactly the marginal constraints

\[
\pi\in\Pi(\mu,\nu).
\]

They enforce \textbf{mass conservation}.

\section*{A simple discrete example}

Suppose

\[
\mu=(0.6,\,0.4)
\]

and

\[
\nu=(0.5,\,0.5).
\]

One admissible joint distribution is

\[
\pi=
\begin{pmatrix}
0.5 & 0.1\\
0 & 0.4
\end{pmatrix}.
\]

Its row sums are

\[
(0.6,\;0.4)=\mu,
\]

and its column sums are

\[
(0.5,\;0.5)=\nu.
\]

This matrix has a natural transportation interpretation:

\begin{itemize}
    \item transport $0.5$ units of mass from $x_1$ to $y_1$,
    \item transport $0.1$ units from $x_1$ to $y_2$,
    \item transport $0.4$ units from $x_2$ to $y_2$.
\end{itemize}

Nothing else is transported.

Notice that the matrix is simultaneously

\begin{itemize}
    \item a joint probability distribution,
    \item a feasible transport plan.
\end{itemize}

The two viewpoints are mathematically identical.

\section*{Connection to the Kantorovich formulation}

The Kantorovich problem minimizes the total transportation cost

\[
\inf_{\pi\in\Pi(\mu,\nu)}
\int_{\mathcal X\times\mathcal Y}
c(x,y)\,d\pi(x,y).
\]

Here,

\begin{itemize}
    \item $c(x,y)$ is the cost of transporting one unit of mass from $x$ to $y$,
    \item $d\pi(x,y)$ is the amount of mass transported along that route.
\end{itemize}

Therefore,

\[
c(x,y)\,d\pi(x,y)
\]

is the infinitesimal transportation cost, and integrating over all pairs $(x,y)$ gives the total work required to transform $\mu$ into $\nu$.

\section*{Connection to WGANs}

In WGANs, the real distribution $P_r$ and the generated distribution $P_g$ are generally known only through samples. The transport plan $\pi$ is \textbf{not computed explicitly}. Instead, the Kantorovich duality shows that the same optimal transport cost can be obtained by optimizing over 1-Lipschitz functions (the critic), avoiding the need to solve directly for $\pi$.

Nevertheless, conceptually, the transport plan remains the object that specifies \textbf{how probability mass from the real distribution is matched to probability mass in the generated distribution}. The critic can be viewed as evaluating the cost of this optimal matching without ever constructing the joint distribution itself.





\chapter{Verification of the assumptions underlying the Energy Score as loss function}
\label{ch:condition_proofs}

The convergence analysis of the scoring-rule-based learning framework proposed by Pacchiardi et al.~\cite{srforgan:srforgan} 
relies on a number of regularity assumptions 
concerning the data generating process and the 
neural network defining the predictive distribution. 
This appendix verifies that these assumptions listed in ~Chapter\ref{ES_conditions} are satisfied by the simulation framework adopted in this work.



\section{Finite-moment assumption}
The Energy Score is defined for distributions possessing finite moments of order $\beta$, where

\[
0<\beta<2.
\]
So it require $\mathbb{E}_{X\sim P^\phi}\|X\|^\beta<\infty$ for the strict propriety of $S_E^{(\beta)}$ and for its unbiased estimator to have finite variance.
Both the generated and real data distributions are required to satisfy this condition.\\

Starting from the simulated processes, we verify that this requirement is met.
For the Black--Scholes model,

\[
X_t
=
X_0
-
\frac12\sigma^2t
+
\sigma W_t,
\]

where

\[
W_t
\sim
\mathcal N(0,t).
\]

Hence,

\[
X_t
\sim
\mathcal N
\left(
X_0-\frac12\sigma^2t,
\sigma^2t
\right),
\]

which possesses finite moments of every order.

Similarly, under the Heston model, provided that the parameters satisfy the usual admissibility conditions

\[
\kappa>0,\qquad
\theta>0,\qquad
\xi>0,
\]

the variance process remains non-negative and the log-price admits finite moments over finite forecasting horizons under the parameter ranges considered in this work.

Since every parameter is sampled from a bounded domain,

\[
\vartheta\in\Theta,
\]

where $\Theta$ is compact, there exists a finite constant

\[
M_\beta<\infty
\]

such that

\[
\sup_{\vartheta\in\Theta}
\mathbb E
\left[
\|X_t\|^\beta
\mid
\vartheta
\right]
\le
M_\beta.
\]

Applying the law of total expectation,

\[
\mathbb E
\|X_t\|^\beta
=
\int
\mathbb E
[
\|X_t\|^\beta
\mid
\vartheta
]
\,d\Pi(\vartheta)
\le
M_\beta
<
\infty.
\]


\vspace{3cm}
Coming to the generated distribution,
if the generator's output layer is bounded (tanh, sigmoid), this is immediate, in fact, $\|X\|$ is almost surely bounded, 
so every moment is trivially finite. 
Our generator instead uses a \emph{linear} output layer, 
so $X=G_\phi(z,c)$ is, in principle, 
unbounded as $z$ ranges over $\mathbb{R}^{d_z}$.\\
$G_\phi(z,c)$ is Lipschitz in $z$. By fixing $\phi$ and $c$ it can be seen that
the map $z\mapsto G_\phi(z,c)$ is a finite composition of two kinds of layers:
\begin{itemize}
\item affine maps $h\mapsto Wh+b$, like the linear output layer,
which are Lipschitz with constant $\|W\|_{\mathrm{op}}$ (the operator norm) finite since $W$ is a finite matrix. Here in our case $h$ is the concatenation of the recurrent embedding $h(c)$ and the latent variable $z$.
\item pointwise nonlinear functions (ReLU, LeakyReLU, tanh, GELU, \ldots), each Lipschitz with a known finite constant $\ell$ (e.g.\ $\ell=1$ for ReLU).
\end{itemize}
A composition of $L_1$ and 
$L_2$-Lipschitz maps is $(L_1L_2)$-Lipschitz, 
so $G_\phi(z,c)$ is $L_\phi(c)$-Lipschitz, 
with $L_\phi(c)$ the finite product of all per-layer 
constants. 
This uses only the composition structure of 
the network, so it holds regardless of whether the 
\emph{output} layer is bounded, and applies equally 
after the recurrent embedding of $c$ has collapsed to a 
fixed finite vector $h(c)$: with $c$ fixed, 
$z\mapsto G_\phi(z,c)$ is a finite composition 
of dense layers acting on $(h(c),z)$.\\

Now by definition of Lipschitz continuity and the 
triangle inequality,
\[
\|G_\phi(z,c)\| \;\le\; \|G_\phi(0,c)\| + L_\phi(c)\,\|z\|,
\]
where $\|G_\phi(z,c)\|$ is the norm of the generator output, $\|G_\phi(0,c)\|$ is a fixed finite 
number (the network evaluated at $z=0$, no randomness 
involved) and $L_\phi(c)$ is a fixed finite 
Lipschitz constant.\\

For $Z\sim\mathcal N(0,I_{d_z})$, the squared norm is a chi-squared random variable with $d_z$ degrees of freedom $\|Z\|^2\sim\chi^2_{d_z}$, and
\[
\mathbb{E}\|Z\|^\beta = 2^{\beta/2}\,\frac{\Gamma\!\left(\tfrac{d_z+\beta}{2}\right)}{\Gamma\!\left(\tfrac{d_z}{2}\right)} <\infty \qquad\text{for every }\beta>0,
\]
where $\Gamma(t)=\int_0^\infty x^{t-1}e^{-x}\,dx$ is Euler's Gamma function, finite for every $t>0$. 
Both arguments above are positive since $d_z,\beta>0$.
Finiteness itself needs no closed form but can be derived from the fact that $\|Z\|^2$ has a density with Gaussian tails,
and Gaussian tails decay faster than any polynomial.\\

For $a,b\ge0$, $(a+b)^\beta\le C_\beta(a^\beta+b^\beta)$, 
with $C_\beta=1$ if $\beta\le1$ (subadditivity of the concave 
map $t^\beta$) and $C_\beta=2^{\beta-1}$ if $\beta>1$ (convexity, 
via Jensen). Apply this with $a=\|G_\phi(0,c)\|$, $b=L_\phi(c)\|z\|$, 
the two terms bounding $\|G_\phi(z,c)\|$, 
then take $\mathbb{E}_{Z\sim\mathcal N(0,I_{d_z})}$ of both sides:

\[
\mathbb{E}_{Z\sim\mathcal N(0,I_{d_z})}\|G_\phi(Z,c)\|^\beta
\;\le\;
C_\beta\Big(\|G_\phi(0,c)\|^\beta + L_\phi(c)^\beta\,\mathbb{E}\|Z\|^\beta\Big)
\;<\;\infty,
\]

finite because $\|G_\phi(0,c)\|$ and $L_\phi(c)$ are 
fixed finite numbers for fixed $\phi,c$ 
and $\mathbb{E}\|Z\|^\beta<\infty$ as seen before. 
Since $X=G_\phi(Z,c)$ 
under the reparameterisation trick, 
this proves $\mathbb{E}_{X\sim P^\phi(\cdot|c)}\|X\|^\beta<\infty$ 
for every fixed $\phi\in\Phi$ and $c$: 
a linear output layer poses no obstruction, 
because linearity is itself a Lipschitz property, 
and the resulting at-most-linear growth in $\|z\|$ 
is dominated by the finiteness of all Gaussian moments.\\

The bound above holds for one fixed $\phi$; 
Lemmas~13--14 of \cite{srforgan:srforgan} 
additionally require it \emph{uniformly} over 
$\phi\in\Phi$, ultimately because their Theorem~3 
assumes $\Phi$ compact (Assumption~A1). 
This does not follow from the architecture 
alone: $L_\phi(c)$ and $\|G_\phi(0,c)\|$ both 
depend on the current weights, and nothing 
prevents $L_\phi(c)\to\infty$ along an 
unconstrained training trajectory. In practice this 
is enforced, not assumed by the use of weight decay or spectral-norm 
regularisations which bounds $\|W\|_{\mathrm{op}}$ 
layer-by-layer, and gradient clipping bounds 
the trajectory travelled by $\phi$, together 
giving a finite $\bar L=\sup_{\phi\in\Phi}L_\phi(c)$ 
over any compact $\Phi$ the optimiser is confined to. 
The moment-boundedness condition of Sec.~\ref{sec:energy_score_theory} 
is therefore satisfiable by construction 
for the proposed generator, 
but it is a property of the \emph{training regime} as a whole.\\


Therefore, the finite-moment assumption required for the Energy Score is satisfied.

%%%%%%%%%%%%%%%%%%%%%%%%%%%%%%%%%%%%%%%%%%%%%%%%%%%%%%%%%%
\section{Unbiasedness of the estimator}

For one condition $c$, let

\[
X_1,\ldots,X_m
\stackrel{\mathrm{i.i.d.}}{\sim}
P_\phi(\cdot\mid c)
\]

denote independent samples generated by the neural network.

The empirical Energy Score used during optimization is

\[
\widehat S_E^{(\beta)}
=
\frac{2}{m}
\sum_{i=1}^m
\|X_i-y\|^\beta
-
\frac{1}{m(m-1)}
\sum_{i\neq j}
\|X_i-X_j\|^\beta.
\]

Using linearity of expectation,

\[
\mathbb E
\left[
\frac{2}{m}
\sum_{i=1}^m
\|X_i-y\|^\beta
\right]
=
2
\mathbb E
\|X-y\|^\beta.
\]

Likewise,

\[
\mathbb E
\left[
\frac1{m(m-1)}
\sum_{i\neq j}
\|X_i-X_j\|^\beta
\right]
=
\mathbb E
\|X-X'\|^\beta,
\]

because every ordered pair
$(X_i,X_j)$,
$i\neq j$,
has the same joint distribution as two independent draws
$X,X'\sim P_\phi(\cdot\mid c)$.

Therefore,
\[
\mathbb E
\left[
\widehat S_E^{(\beta)}
\right]
=
S_E^{(\beta)},
\]


showing that the estimator implemented during training is unbiased, exactly as established in Appendix C of Pacchiardi et al.~\cite{srforgan:srforgan}.

%%%%%%%%%%%%%%%%%%%%%%%%%%%%%%%%%%%%%%%%%%%%%%%%%%%%%%%%%%

\section{Independence of the generated observations}


Although different trajectories may be generated using different model parameters (for example different values of the volatility parameter in the Black--Scholes model or different Heston parameters), this does not violate the theoretical assumptions underlying scoring-rule minimization. Indeed, if the parameters $\vartheta$ of the underlying data generation process are independently 
sampled from a probability distribution $\Pi(\vartheta)$ before simulating each trajectory, every observation is generated according to a hierarchical model. In fact, the first step of the generative process is to sample the parameters $\vartheta_j$ for each trajectory $j$ from $\Pi$, and then to sample the condition--target pair $(c_j,y_j)$ from the conditional law $P(C,Y \mid\vartheta)$ defined by those parameters. Formally:
\[
\vartheta_j\sim\Pi,
\qquad
(c_j,y_j)\sim P(C,Y \mid\vartheta),
\]
which induces the marginal joint distribution. In fact, starting from the joint law and the fact that the $J$ trajectories 
are independently generated:

\[
d \mathbb{P}\left(\vartheta_1, c_1, y_1, \ldots, \vartheta_J, c_J, y_J\right)=
\prod_{j=1}^J\left[P\left(d c_j, d y_j \mid \vartheta_j\right) 
d \Pi\left(\vartheta_j\right)\right]
\]

then by marginalizing out the latent parameters $\vartheta_j$ we obtain the marginal law of the observed data:

\[
d \mathbb{P}\left(c_1, y_1, \ldots, c_J, y_J\right)=
\int \cdots \int \prod_{j=1}^J\left[P\left(d c_j, d y_j 
\mid \vartheta_j\right) d \Pi\left(\vartheta_j\right)\right]
\]

by applying Tonelli's theorem and integrating out the latent parameters $\vartheta_j$:
\[
=\prod_{j=1}^J 
\underbrace{\left(\int P\left(d c_j, d y_j \mid 
\vartheta_j\right) d \Pi\left(\vartheta_j\right)\right)}_{=P\left(d c_j, d y_j\right)}=\prod_{j=1}^J 
P\left(d c_j, d y_j\right)
\]


the label $\vartheta$
is merely the (dummy) variable of integration, 
and every trajectory shares the same 
support $\Pi$ for the latent generating process's parameters 
and the same conditional family 
$P(C,Y\mid\vartheta)$,
the marginal law of the observed data is simply a product of $J$ identical measures $P(C,Y)$, each of which is the marginal law of a single condition, target pair $(c_j,y_j)$ obtained by integrating out the latent parameters $\vartheta$ from the conditional law.
So every factor in the product is the identical measure 
$P(C,Y)=
\int
P(C,Y\mid\vartheta)
\,d\Pi(\vartheta)$, only evaluated at a different observed point.

Therefore,
\[
(c_1,y_1),\ldots,(c_J,y_J)
\stackrel{\mathrm{i.i.d.}}{\sim}
P(C,Y),
\]

So, even though each trajectory may have been simulated with different underlying stochastic parameters. The variation in the parameters is absorbed into the marginal data distribution and does not affect the strict propriety of the scoring rule, which is defined conditionally on the observed conditioning variable.

%%%%%%%%%%%%%%%%%%%%%%%%%%%%%%%%%%%%%%%%%%%%%%%%%%%%%%%%%%



\section{Differentiability of the learning objective}

The generator defines the predictive distribution through the reparameterization

\[
X
=
G_\phi(z,c),
\qquad
z\sim\mathcal N(0,I).
\]

The latent variable is independent of the network parameters, while the mapping

\[
G_\phi(\cdot,\cdot)
\]

is differentiable almost everywhere because it is composed of affine mappings and standard neural-network activation functions.

Consequently,

\[
\mathbb E
\left[
\|G_\phi(z,c)-y\|^\beta
\right]
\]

is differentiable with respect to $\phi$.

Under the bounded-moment conditions verified above, the Dominated Convergence Theorem justifies exchanging differentiation and expectation,

\[
\nabla_\phi
\mathbb E
[S_E]
=
\mathbb E
[
\nabla_\phi S_E
].
\]

Therefore, the stochastic gradients computed through automatic differentiation constitute unbiased estimators of the gradient of the population objective. Standard stochastic gradient descent convergence results therefore apply, exactly as shown in Theorem 1 and Appendix C of Pacchiardi et al.~\cite{srforgan:srforgan}.

%%%%%%%%%%%%%%%%%%%%%%%%%%%%%%%%%%%%%%%%%%%%%%%%%%%%%%%%%%

The previous results demonstrate that the simulation framework adopted in this work satisfies the assumptions required by the theory of strictly proper scoring-rule minimization. Although individual trajectories are generated using different realizations of the underlying stochastic model parameters, these parameters are sampled independently from a common distribution before each simulation. Consequently, the resulting condition--target pairs constitute an i.i.d. sample from a joint distribution, the Energy Score remains strictly proper, its Monte Carlo approximation is unbiased, and the optimization procedure implemented in this work inherits the theoretical guarantees established by Pacchiardi et al.

\end{appendices}



%\nocite{*}
\printbibliography

\end{document}